\documentclass[12pt,a4paper,oneside]{book}


% --- Language and Fonts ---
\usepackage{lmodern}
\usepackage[T1]{fontenc}
\usepackage[utf8]{inputenc}
\usepackage[polish]{babel}


% --- Layout ---
\usepackage{geometry}
\geometry{margin=2.5cm}
\setlength{\marginparwidth}{3.5cm}  % wider for todonotes compatibility

% --- Headers and Footers ---
\usepackage{fancyhdr}
\pagestyle{fancy}
\fancyhf{}
\fancyhead[L]{\nouppercase{\leftmark}}
\fancyhead[R]{\thepage}
\renewcommand{\headrulewidth}{0.4pt}

% --- Math and Symbols ---
\usepackage{amsmath,amsfonts,amssymb}
\usepackage{amsthm}


% --- Graphics and Tables ---
\usepackage{graphicx}
\usepackage{booktabs}
\usepackage{array}
\usepackage{longtable}
\usepackage{multirow}
\usepackage{caption}
\usepackage{float}


% --- Quotes ---
\usepackage{csquotes}

% --- Algorithms ---
\usepackage{algorithm}
\usepackage{algpseudocode}
\usepackage{tikz}
\usetikzlibrary{arrows.meta, shapes.geometric, patterns}
\usepackage{xcolor}

% --- Utilities ---
\usepackage{verbatim}
\usepackage{enumitem}
\usepackage{titlesec}
\usepackage{authblk}
\usepackage{todonotes}


% \usepackage{csvsimple-13}
\usepackage{csvsimple}
% Upewnij się, że separator jest zdefiniowany globalnie
\csvset{separator=semicolon}
% SetWatermark
\usepackage{draftwatermark}
\SetWatermarkText{DRAFT}        % tekst znaku wodnego
\SetWatermarkScale{1}           % skala (1 = domyślna)
\SetWatermarkColor[gray]{0.65}  % kolor (jasny szary)
\SetWatermarkAngle{45}          % kąt obrotu

% --- Bibliography ---
\usepackage[
  backend=biber,
  style=verbose,
  autocite=footnote,
  sorting=none
]{biblatex}
\DeclareBibliographyAlias{legislation}{misc}
\addbibresource{bib-law.bib}
\addbibresource{bib-other.bib}


% hyperref must be loaded last (before custom settings)
\usepackage[hidelinks]{hyperref}

% --- Theorem Environments ---
\newtheorem{definition}{Definicja}[chapter]
\newtheorem{algotheorem}{Algorytm}[chapter]


% --- Chapter Title Format ---
\titleformat{\chapter}[hang]
  {\normalfont\huge\bfseries}
  {\thechapter.}
  {0.5em}
  {}

\titlespacing*{\chapter}{0pt}{0pt}{40pt}



% --- Metadata ---
\title{ORION. Ocena Ryzyka Informatycznego Organizacji Nadzorowanych\\
Architektura, bezpieczeństwo i suwerenność danych\\
w regulowanych ekosystemach cyfrowych}
\title{Metodologia ORION.\\
Formalny model oceny ryzyka informatycznego\\
w organizacjach nadzorowanych}

\author[1]{Bartosz Lewandowski \thanks{ORCID: 0009-0004-1158-0693}}
\affil[1]{Department of Clinical Engineering, Academy of Silesia in Katowice}
\date{} 

%\includeonly{rozdzial-przeprowadzenie-audytu}

\begin{document}

%\frontmatter
\maketitle
\tableofcontents

%\mainmatter
\chapter{Preambuła}

Luka kompetencyjna na styku ocen rozwiązań informatycznych przez zespoły prawne oraz techniczne jest przyczyną nieprawidłowej klasyfikacji systemów informatycznych. 
Co więcej, organizacje nieprawidłowo zarządzają łańcuchem dostawców i poddostawców, nawet gdy są do tego zobowiązane przez rozporządzenie DORA oraz dyrektywę NIS2. Szczególnie jest to widoczne w relacjach z hiperskalerami objętymi jurysdykcją Cloud Act, gdzie bardzo często pomija się ryzyka transferu danych do USA, co stoi w sprzeczności z RODO.

Ocena Ryzyka Informatycznego Podmiotów Nadzorowanych, w skrócie zwana ORION, jest nową metodą pozwalającą na wykrywanie i adresowanie wskazanych luk przez organizacje.
Metoda adresowana jest przede wszystkim do podmiotów objętych nadzorem regulacyjnym, instytucji finansowych, operatorów usług kluczowych oraz dostawców usług cyfrowych, dla których zgodność z wymogami DORA, NIS2 i RODO stanowi obowiązek prawny, a nie dobrowolny standard. Warto jednak wykorzystać medodę ORION również w pozostałych przedsiębiorstwach, traktując podnoszone przez metodologię wyzwania wynikające z implementacji DORA oraz NIS2 jako zestaw dobrych praktych, równocześnie upewniając się, że działania podmiotu są zgodne z obowiązującymi przepisami np. w zakresie RODO.

Niniejsza publikacja przedstawia założenia teoretyczne metody ORION, opisuje jej strukturę klasyfikacyjną oraz ilustruje jej zastosowanie praktyczne na przykładzie wybranej organizacji.

\chapter{Wstęp}

Transformacja cyfrowa przedsiębiorstw i instytucji publicznych, postępująca intensywnie od końca XX wiek,  zmieniła charakter ryzyk i wyzwań stawianych przed organizacjami. Systemy informatyczne przestały być wyłącznie narzędziami wspierającymi codzienną pracę, stały się jej elementem strategii przedsiębiorstw, bez którego realizacja procesów jest praktycznie niemożliwa. Dotyczy to w równym stopniu banku przetwarzającego miliony transakcji konsumentów, szpitala prowadzącego elektroniczną dokumentację medyczną pacjentów, operatora infrastruktury energetycznej zarządzającego siecią dystrybucyjną, uczelni i szkoły prowadzącej cyfrowe ewidencje postępów podopiecznych, firm produkcyjnych opierających ciągłość działania na cyfrowych algorytmach. We wszystkich tych przypadkach niedostępność systemu informatycznego awarią o ograniczonych skutkach, jest zdarzeniem operacyjnym, prawnym często niosącym negatywne skutki finansowe.

Ustawodawca zdefiniował standardy bezpieczeństwa opisane w przepisach prawa. Do kluczowych aktów prawnych, które wymagają prowadzenia szczegółowej analizy ryzyka, należą przede wszystkim Ustawa o Krajowym Systemie Cyberbezpieczeństwa \autocite{uksc2018} oraz Rozporządzenie w sprawie operacyjnej odporności cyfrowej sektora finansowego, potocznie zwane DORA \autocite{dora2022}. Choć powyższe dokumenty jasno określają, kogo dotyczą ich zapisy, warto potraktować te publikacje jako uniwersalne źródło wiedzy. Wskazane w nich procedury oraz opisane wektory zagrożeń mogą pomóc zabezpieczyć każdą organizację przed skutkami cyberataków.

Równolegle z rosnącym uzależnieniem organizacji od technologii informatycznych rozwinął się złożony rynek dostawców, poddostawców i partnerów technologicznych, tworzący łańcuchy zależności ICT (\textit{Information and Communications Technology}). Pojedyncza aplikacja biznesowa może dziś opierać się na wielu elementach dostarczanych przez szereg dostawców: platformach chmurowych, bibliotekach oprogramowania, usługach uwierzytelniania, systemach płatności, zewnętrznych centrach przetwarzania danych czy dostawcach łączności. Każde z tych elementów jest potencjalnym źródłem ryzyka, technicznego (awarii), prawnego (brak zgodności z obowiązującym prawem) i operacyjnego. To organizacja ponosi odpowiedzialność regulacyjną za skutki zdarzeń w całym tym łańcuchu, niezależnie od tego, czy posiada nad nim bezpośrednią kontrolę. Złożoność łańcuchów dostaw ICT sprawia, że ich pełna inwentaryzacja, ocena i nadzór są zadaniem, z którym większość organizacji radzi sobie niewystarczająco,  nie z braku dobrej woli, lecz z braku odpowiednich narzędzi metodycznych.

Problem ten nabiera szczególnego wymiaru w kontekście intensywnej aktywności legislacyjnej europejskiego prawodawcy ostatniej dekady. Ochrona danych stanowi jeden z fundamentalnych elementów współczesnego porządku prawnego Unii Europejskiej. Dynamiczny rozwój technologii, w szczególności sztucznej inteligencji oraz gromadzenia  informacji, a co za tym idzie szeroka cyfryzacja usług publicznych i prywatnych oraz rosnąca skala międzynarodowego przetwarzania danych osobowych sprawiły, że zagadnienie ochrony danych stało się przedmiotem prac legislacyjnych zarówno na poziomie unijnym, jak i krajowym. Unia Europejska oraz jej państwa członkowskie zwracają szczególną uwagę na jakość ochrony danych, uznając ją za istotny element ochrony praw jednostki, zaufania do gospodarki oraz bezpiecznego funkcjonowania społeczeństwa informacyjnego. Regulacje te, omówione szczegółowo w rozdziale \textit{Krajobraz regulacyjny}, tworzą środowisko, w którym zgodność z prawem stała się warunkiem koniecznym prowadzenia działalności w sektorach takich jak finanse, ochrona zdrowia, energetyka czy transport. Niespełnienie wymogów wiąże się z sankcjami finansowymi wobec przedsiębiorstw, a odpowiedzialność za naruszenia spoczywa bezpośrednio na organach zarządczych podmiotów nadzorowanych.

W tym kontekście szczególnego znaczenia nabiera pytanie badawcze niniejszej monografii: w jaki sposób organizacja może w sposób wiarygodny, powtarzalny i oparty na obiektywnych dowodach ocenić rzeczywisty stan zgodności swoich systemów informatycznych i ich łańcucha dostawców z obowiązującymi wymogami prawnymi? Jest to teoretycznie proste pytanie. Praktyka audytowa i badawcza pokazuje bowiem, że istniejące metodyki, choć wartościowe w swoich dziedzinach, nie dają na nie zadowalającej odpowiedzi. Powstała więc luka, której wymiary nie były do tej pory szacowane, jednak luka jest realna, dotychczas nieadresowana, jednocześnie wskazuje na braki prawidłowej oceny zgodności zarządzania systemami informatycznymi.
Deklarowana zgodność w wymagami regulatorów rynku często rozmija się z weryfikowalnym stanem faktycznym. Organizacja posiadają certyfikaty i polityki bezpieczeństwa, lecz nie potrafią wskazać, które konkretne systemy przetwarzają dane szczególnej kategorii, kto jest faktycznym podddostawcą kluczowej platformy chmurowej ani czy dostawca zewnętrzny kiedykolwiek był przedmiotem oceny zgodności. Zjawisko to, opisywane w literaturze anglojęzycznej jako \textit{compliance theater}, wynika ze słabości systemów oceny opartych wyłącznie na deklaracjach i samoocenie. Jest ono równie szkodliwe, jak jawna niezgodność: daje organizacji fałszywe poczucie bezpieczeństwa.

\section{Postulat: Ocena Ryzyka Nadzorowanych Systemów Informatycznych }
Prezentowana praca opisuje metodologię Oceny Ryzyka Informatycznego Organizacji Nadzorowanych (ORION), oryginalny framework metodyczny zaprojektowany w celu wypełnienia opisanej luki. ORION nie jest kolejną implementacją istniejących standardów, nie zastępuje audytów ISO czy NIST, nie jest też systemem informatycznym w technicznym sensie tego słowa. Jest formalnie zdefiniowaną metodyką oceny zgodności systemów ICT i ich łańcucha dostaw, opartą na trzech innowacjach konceptualnych.

Fundamentem metodywki jest wieloperspektywiczność, którą można wytłumaczyć jako segmentacje kompetencyjną. Celem jest rozdzielenie procesu oceny na płaszczyzny: prawną, biznesową (użytkową) oraz techniczną. Zapewnia to, że odpowiedzi udzielają wyłącznie osoby posiadające realne kompetencje w danej dziedzinie.

Kolejny istotny element to rygor dowodowy. Kluczem jest zastąpienie deklaratywności ocen weryfikacją artefaktów (logi, certyfikaty). System wprowadza mechanizm penalizacji scoringowej za brak materialnych dowodów, co eliminuje zjawisko fasadowej zgodności (compliance theater).

Finalnie wynik mapowany do jako sieć powiązań w modelu grafowym. Dzięki algorytmom propagacji ryzyka, ORION pozwala wyliczyć, jak błędy poddostawcy wpływają na bezpieczeństwo i ciągłość działania podmiotu nadrzędnego.

\chapter{Krajobraz regulacyjny}

\section{Krajowe ramy prawne cyberbezpieczeństwa i~ochrony danych}

Polska, jako państwo członkowskie Unii Europejskiej, od lat aktywnie uczestniczy w budowaniu europejskiego systemu ochrony danych i~cyberbezpieczeństwa. Punktem wyjścia dla krajowych regulacji w~tym zakresie stało się wdrożenie unijnej dyrektywy NIS z~2016 roku \autocite{dyrektywa_nis}, która zobowiązała państwa członkowskie do ustanowienia krajowych norm cyberbezpieczeństwa. W~odpowiedzi na te wymagania, Polska uchwaliła ustawę o~krajowym systemie cyberbezpieczeństwa, przyjętą przez Sejm RP dnia 5~lipca 2018~roku i~podpisaną przez Prezydenta RP dnia 1~sierpnia 2018~r., która weszła w~życie 28~sierpnia 2018~r.
\autocite{uksc2018}.

\subsection{Zakres obowiązków wynikający z ustawy o~krajowym systemie cyberbezpieczeństwa z~2018 r.}

Ustawa o~krajowym systemie cyberbezpieczeństwa~\autocite{uksc2018} nakłada kompleksowe obowiązki przede wszystkim na operatorów usług kluczowych (\textit{OUK}) oraz dostawców usług cyfrowych (\textit{DUC}), działających w~sektorach takich jak energetyka, transport, bankowość, ochrona zdrowia czy infrastruktura cyfrowa.

\begin{table}[H]
\centering
\caption{Główne obowiązki nałożone ustawą o~KSC z~2018~r.}
\label{tab:ksc_obowiazki}
\begin{tabular}{>{\bfseries}p{4cm} p{9.5cm}}
\hline
Obszar obowiązku & Opis \\
\hline
Zarządzanie ryzykiem
  & Regularna identyfikacja i~ocena ryzyk cyberbezpieczeństwa; opracowanie
    i~wdrożenie strategii ochrony adekwatnej do specyfiki działalności. \\[4pt]
Środki techniczne i~org.
  & Zapewnienie bezpiecznej eksploatacji systemów, w~tym bezpieczeństwa
    fizycznego, kontroli dostępu oraz ciągłości świadczenia usług. \\[4pt]
Monitorowanie i~zgłaszanie incydentów
  & Bieżące monitorowanie systemów informatycznych; raportowanie
    incydentów do właściwych zesp.\ CSIRT (Zespół reagowania na incydenty bezpieczeństwa komputerowego)
    (CSIRT GOV, CSIRT NASK, CSIRT MON). \\[4pt]
Dokumentacja
  & Opracowanie i~przechowywanie przez co~najmniej 2~lata dokumentacji
    dotyczącej cyberbezpieczeństwa systemu informacyjnego. \\[4pt]
Audyty
  & Przeprowadzanie regularnych audytów bezpieczeństwa oraz szkoleń
    dla pracowników. \\[4pt]
Współpraca z~organami
  & Koordynacja działań z~organami właściwymi ds.\ cyberbezpieczeństwa;
    udostępnianie wyników audytów w~razie kontroli. \\
\hline
\end{tabular}
\end{table}

\subsection{Nowelizacja UKSC z~23 stycznia 2026~roku\ (implementacja NIS2)}

W~dniu 23~stycznia 2026~r.\ Sejm uchwalił nowelizację ustawy o~krajowym systemie cyberbezpieczeństwa~\autocite{uksc2026}, stanowiącą implementację
unijnej dyrektywy NIS2~\autocite{dyrektywa_nis2} do polskiego porządku prawnego. Ustawa została następnie podpisana przez Prezydenta RP, który jednocześnie
skierował ją do Trybunału Konstytucyjnego, podnosząc wątpliwości dotyczące m.in.\ obowiązku wymiany sprzętu bez odszkodowania oraz sposobu podejmowania
decyzji przez organy nadzorcze.

\begin{table}[H]
\centering
\caption{Kluczowe zmiany wprowadzone nowelizacją UKSC z~2026~r.}
\label{tab:ksc_nowelizacja}
\begin{tabular}{>{\bfseries}p{5cm} p{8.5cm}}
\hline
Zmiana & Opis \\
\hline
Rozszerzenie zakresu podmiotowego
  & Poszerzenie katalogu branż i~obniżenie progu wejścia do
    systemu \\[4pt]
Obowiązek samoidentyfikacji
  & Organizacja samodzielnie  ma obowiązek ustalić swój
    status, dokonać kwalifikacji i~złożyć wniosek o~wpis do
    właściwego wykazu. \\[4pt]
Odpowiedzialność zarządcza
  & Odpowiedzialność przestaje być wyłącznie operacyjna i~staje się
    elementem ładu korporacyjnego oraz osobistej odpowiedzialności
    członków zarządu. \\[4pt]
Kary finansowe
  & Podmioty kluczowe: do 10~mln~EUR lub 2\% rocznego światowego obrotu;
    podmioty ważne: do 7~mln~EUR lub 1,4\% obrotu. \\[4pt]
Odroczenie kar
  & Kary pieniężne za naruszenia mogą być nakładane dopiero po upływie
    2~lat od wejścia w~życie ustawy. \\[4pt]
Modyfikacje definicji
  & Zawężenie definicji organizacji badawczej i~doprecyzowanie zakresu
    podmiotowego wobec środowisk naukowych. \\
\hline
\end{tabular}
\end{table}

\subsection{Ustawa o~ochronie danych osobowych z~10 maja 2018~r.}

Z punktu widzenia bezpieczeństwa, istotną ustawą w~krajowym prawodawstwie jest Ustawa z~dnia 10~maja 2018~r.\
o~ochronie danych osobowych~\autocite{uodo2018}, która stanowi akt wykonawczy wobec rozporządzenia RODO~\autocite{rodo} i~reguluje m.in.\ zasady działania Prezesa Urzędu Ochrony Danych Osobowych (PUODO) jako krajowego organu nadzorczego, tryb postępowania przed tym organem, szczegółowe zasady przetwarzania danych oraz krajowe przepisy uzupełniające RODO tam, gdzie rozporządzenie
pozostawia swobodę regulacyjną państwom członkowskim.

\subsubsection{Odpowiedzialność administracyjna na gruncie RODO}

Rozporządzenie RODO~\autocite{rodo} przewiduje dwustopniowy system administracyjnych kar
pieniężnych (art.~83 RODO).

\begin{table}[H]
\centering
\caption{Administracyjne kary pieniężne na mocy art.~83 RODO}
\label{tab:rodo_kary}
\begin{tabular}{>{\bfseries}p{4.5cm} p{9cm}}
\hline
Wysokość kary & Przykładowe naruszenia \\
\hline
Do 10~mln~EUR lub 2\% rocznego obrotu
  & Brak rejestru czynności przetwarzania; brak wyznaczenia inspektora ochrony
    danych w~wymaganych przypadkach; brak współpracy z~organem nadzorczym. \\[4pt]
Do 20~mln~EUR lub 4\% rocznego obrotu
  & Złamanie podstawowych zasad przetwarzania danych; naruszenie praw osób,
    których dane dotyczą; nielegalny transfer danych do państw trzecich. \\
\hline
\end{tabular}
\end{table}

Organ nadzorczy dysponuje również wachlarzem środków nacisku (art.~58 RODO): ostrzeżenia i~upomnienia, nakazanie spełnienia żądań osoby,
której dane dotyczą, nakazanie dostosowania operacji przetwarzania do przepisów RODO, a~w~najpoważniejszych przypadkach, czasowe lub całkowite ograniczenie
przetwarzania, w~tym zakaz przetwarzania danych oraz nakaz ich usunięcia lub powiadomienia o~naruszeniu.

\subsection{Ustawa o~zwalczaniu nieuczciwej konkurencji (1993)}

Ustawa z~dnia 16~kwietnia 1993~r.\ o~zwalczaniu nieuczciwej konkurencji~\autocite{uznk1993} stanowi istotne uzupełnienie systemu ochrony danych, zwłaszcza w~zakresie ochrony tajemnicy przedsiębiorstwa. Zgodnie z~art.~11 u.z.n.k., czynem nieuczciwej konkurencji jest ujawnienie, wykorzystanie lub pozyskanie cudzych informacji stanowiących tajemnicę przedsiębiorstwa. Za tajemnicę przedsiębiorstwa uznaje się informacje techniczne,
technologiczne, organizacyjne lub inne posiadające wartość gospodarczą, które nie są powszechnie znane i~co do których przedsiębiorca podjął należyte środki
w~celu zachowania ich poufności.

W~kontekście naruszeń danych i~cyberbezpieczeństwa ustawa ta nabiera szczególnego znaczenia: wyciek danych stanowiących tajemnicę przedsiębiorstwa (np.\ w~wynik ataku hakerskiego, działania nielojalnego pracownika lub wadliwego zabezpieczenia systemów) może jednocześnie stanowić naruszenie RODO oraz czyn nieuczciwej konkurencji z~u.z.n.k. Ustawa przewiduje cywilnoprawne środki ochrony, w~tym roszczenia o~zaniechanie, usunięcie skutków, naprawienie szkody i~wydanie bezpodstawnie uzyskanych korzyści~\autocite{uznk1993}.

\subsection{Odpowiedzialność cywilna wg kodeksu cywilnego}

Odpowiedzialności z~tytułu naruszenia przepisów o~ochronie danych nie można rozpatrywać wyłącznie przez pryzmat RODO, kluczową rolę odgrywa tu również
Kodeks cywilny~\autocite{kc1964}.

\begin{table}[H]
\centering
\caption{Reżimy odpowiedzialności cywilnej w~obszarze ochrony danych}
\label{tab:odpowiedzialnosc_cywilna}
\begin{tabular}{>{\bfseries}p{5cm} p{8.5cm}}
\hline
Podstawa prawna & Opis \\
\hline
Art.~415 k.c.: odpowiedzialność deliktowa
  & Administrator lub podmiot przetwarzający, który ze swojej winy wyrządził
    szkodę osobie fizycznej poprzez naruszenie zasad ochrony danych,
    zobowiązany jest do jej naprawienia. Przesłankami są: bezprawne działanie
    lub zaniechanie, wina, szkoda oraz związek przyczynowo-skutkowy. \\[4pt]
Art.~23, 24 i~448 k.c.: dobra osobiste
  & Prywatność jest traktowana jako dobro osobiste. W~razie jego naruszenia
    pokrzywdzony może żądać zaniechania naruszeń i~zasądzenia
    zadośćuczynienia (sądy orzekają w~zbiegu z~art.~82 RODO). \\[4pt]
Art.~82 ust.~4 RODO w~zw.\ z~k.c.: solidarność
  & Jeżeli w~tym samym przetwarzaniu uczestniczy więcej niż jeden podmiot,
    odpowiadają oni solidarnie za całą szkodę, co zapewnia poszkodowanemu
    rzeczywiste uzyskanie odszkodowania. \\[4pt]
Zbieg podstaw odpowiedzialności
  & Nałożenie kary przez PUODO nie zwalnia administratora z~równoległej
    odpowiedzialności cywilnej wobec osób poszkodowanych. \\
\hline
\end{tabular}
\end{table}

\subsection{Odpowiedzialność karna: Kodeks karny}

Odrębnym, lecz komplementarnym reżimem jest odpowiedzialność karna przewidziana w~Kodeksie karnym~\autocite{kk1997}, w~szczególności w~rozdziale XXXIII
,,Przestępstwa przeciwko ochronie informacji'' (art.~265--269b k.k.).

\begin{table}[H]
\centering
\caption{Przestępstwa z~zakresu ochrony informacji (art.~265--269b k.k.)}
\label{tab:kk_przestepstwa}
\begin{tabular}{>{\bfseries}p{3.5cm} p{10cm}}
\hline
Przepis & Opis \\
\hline
Art.~267 k.k.\: hacking
  & Kara grzywny, ograniczenia wolności albo pozbawienia wolności do 2~lat
    za bezprawne uzyskanie dostępu do informacji poprzez przełamanie
    zabezpieczeń elektronicznych, podpięcie do sieci lub stosowanie
    oprogramowania szpiegującego; penalizowane również ujawnianie
    tak pozyskanych informacji. \\[4pt]
Art.~268 i~268a k.k. & Niszczenie lub uszkadzanie danych informatycznych i~udaremnianie osobie
    uprawnionej zapoznania się z~informacją; kara do 3~lat pozbawienia
    wolności, a~przy znacznej szkodzie majątkowej od 3~miesięcy do 5~lat. \\[4pt]
Art.~269 k.k.: sabotaż komputerowy
  & Uszkodzenie danych informatycznych lub oprogramowania o~szczególnym
    znaczeniu dla obronności lub bezpieczeństwa państwa;
    kara pozbawienia wolności od 6~miesięcy do 8~lat. \\[4pt]
Art.~269a k.k.: zakłócenie systemu
  & Istotne zakłócenie pracy systemu komputerowego lub sieci
    teleinformatycznej przez transmisję, usunięcie lub zmianę danych;
    kara pozbawienia wolności od 3~miesięcy do 5~lat. \\[4pt]
Art.~269b k.k.: narzędzia ataków
  & Wytwarzanie, pozyskiwanie lub udostępnianie narzędzi informatycznych
    przeznaczonych do popełnienia powyższych przestępstw. \\[4pt]
Art.~23 u.z.n.k.\ (uzupełniająco)
  & Ujawnienie tajemnicy przedsiębiorstwa przez osobę zobowiązaną do jej
    zachowania, jeżeli wyrządza to poważną szkodę: kara grzywny, ograniczenia
    wolności albo pozbawienia wolności do 2~lat. \\
\hline
\end{tabular}
\end{table}

Warto podkreślić, że reżim karny funkcjonuje równolegle do odpowiedzialności administracyjnej (kary RODO/UKSC) i~cywilnej, skazanie sprawcy nie eliminuje
roszczeń odszkodowawczych poszkodowanych ani nie uchyla administracyjnych kar nałożonych na administratora danych~\autocite{kk1997}\autocite{rodo}.


\subsection{Dyrektywa NIS2 jako fundament europejskiego cyberbezpieczeństwa}

Dyrektywa Parlamentu Europejskiego i~Rady (UE) 2022/2555 z dnia 14~grudnia 2022~r., znana jako NIS2~\autocite{dyrektywa_nis2}, reformuje unijny system cyberbezpieczeństwa. NIS2 uchyliła i~zastąpiła dyrektywę NIS z~2016~r.~\autocite{dyrektywa_nis}, która  okazała się niewystarczająca wobec  wzrostu zagrożeń cybernetycznych oraz rosnącej zależności gospodarki i~administracji od systemów informatycznych.
Dyrektywa NIS2 weszła w życie 16~stycznia 2023~r., zaś państwa członkowskie zobowiązane były do dokonania jej transpozycji do krajowych porządków prawnych
do 17~października 2024~r. W Polsce implementacji dyrektywy dokonała nowelizacja ustawy o Krajowym Systemie Cyberbezpieczeństwa z~dnia 23~stycznia 2026~r.~\autocite{uksc2026}.

\subsection*{Zakres podmiotowy: podmioty kluczowe i~ważne}

Dyrektywa NIS2 odchodzi od dotychczasowego podziału na operatorów usług kluczowych
i dostawców usług cyfrowych, dyrektywa definiuje nowe kategorie: {podmioty kluczowe} (ang.\ \textit{essential entities}) oraz {podmioty ważne}
(ang.\ \textit{important entities}). 
Dyrektywą objęte są podmioty będące co najmniej średnim przedsiębiorstwem wg definicji MŚP\autocite{deficja_msp_2018} (tj.\ zatrudniające min. 50 osób i~osiągające obrót roczny powyżej 10~mln~EUR), działające w sektorach wskazanych w załącznikach i~i II do dyrektywy. Jednocześnie przewidziano szereg wyjątków,
na mocy których dyrektywa stosuje się niezależnie od wielkości podmiotu.

\begin{table}[H]
\centering
\caption{Sektory objęte dyrektywą NIS2 według kategorii podmiotów}
\label{tab:nis2_sektory}
\begin{tabular}{>{\bfseries}p{4.5cm} p{8.5cm}}
\hline
Kategoria & Sektory (wybrane) \\
\hline
Podmioty kluczowe (załącznik I)
 & Energetyka (energia elektryczna, gaz, ropa, wodór); transport (lotniczy,
    kolejowy, wodny, drogowy); bankowość i~infrastruktura rynków finansowych;
    ochrona zdrowia; woda pitna i~ścieki; infrastruktura cyfrowa (IXP, DNS,
    usługi chmurowe, centra danych, sieci CDN); zarządzanie usługami ICT (B2B);
    administracja publiczna; przestrzeń kosmiczna. \\[4pt]
Podmioty ważne (załącznik II)
 & Usługi pocztowe i~kurierskie; gospodarka odpadami; produkcja
    i~dystrybucja chemikaliów; produkcja i~dystrybucja żywności; produkcja
    (wyroby medyczne, sprzęt elektroniczny, pojazdy); dostawcy usług cyfrowych
    (platformy handlowe, wyszukiwarki, serwisy społecznościowe); organizacje
    badawcze. \\
\hline
\end{tabular}
\end{table}

Zakresem dyrektywy nie są objęte podmioty administracji publicznej prowadzące działalność w dziedzinie bezpieczeństwa narodowego, obronności lub ścigania
przestępstw, podmioty wyłączone ze stosowania rozporządzenia DORA~\autocite{dora2022} (stosownie do art.~2 ust.~4 DORA), a~także dostawcy świadczący usługi
wyłącznie na rzecz zwolnionych organów administracji publicznej. Z obowiązku zwolnienia nie można jednak skorzystać wobec dostawców usług zaufania.

\subsection*{Obowiązki podmiotów kluczowych i~ważnych}

Dyrektywa NIS2 wprowadza identyczne obowiązki dla obu kategorii podmiotów (kluczowe oraz ważne), rozróżniając jedynie intensywność nadzoru i~poziom kar. Obowiązki te mają być stosowane w sposób proporcjonalny do specyfiki, wielkości i~poziomu narażenia danego podmiotu na ryzyko.

Podmioty muszą zaimplementować \textit{System Zarządzania Bezpieczeństwem Informacji}, w tym zarządzania ryzykiem. Wymaga do opracowania, wdrożenia i~testowania planów ciągłości działania (\textit{BCP}) oraz politykę bezpieczeństwa łańcucha dostaw, polegająceł na szacowaniu ryzk współpracy w dostawcami oraz poddostawcami.

Każdy podmiot kluczowy i~ważny zobowiązany jest do wdrożenia odpowiednich i~proporcjonalnych środków technicznych, operacyjnych i~organizacyjnych,
obejmujących co najmniej:

\begin{table}[H]
\centering
\caption{Minimalne środki zarządzania ryzykiem wymagane przez art.~21 NIS2}
\label{tab:nis2_ryzyko}
\begin{tabular}{>{\bfseries}p{4.5cm} p{8.5cm}}
\hline
Obszar & Wymagania minimalne \\
\hline
Polityki bezpieczeństwa i~analizy ryzyka
  & Opracowanie i~wdrożenie polityk bezpieczeństwa systemów informatycznych;
    regularna ocena ryzyka uwzględniająca zagrożenia, prawdopodobieństwo
    i~dotkliwość incydentów oraz skutki społeczne i~gospodarcze. \\[4pt]
Obsługa incydentów
  & Ustanowienie procedur wykrywania, analizy, reakcji i~odtwarzania
    po incydentach bezpieczeństwa. \\[4pt]
Ciągłość działania
  & Wdrożenie planów ciągłości działania (BCP) i~odtwarzania po awarii (DRP),
    w tym zarządzanie kopiami zapasowymi i~zarządzanie kryzysowe. \\[4pt]
Bezpieczeństwo łańcucha dostaw
  & Ocena i~zarządzanie ryzykiem ze strony dostawców i~usługodawców ICT;
    stosowanie klauzul bezpieczeństwa w umowach z dostawcami. \\[4pt]
Bezpieczeństwo sieci i~systemów
  & Środki kontroli dostępu, uwierzytelnianie wieloskładnikowe (MFA),
    szyfrowanie danych w spoczynku i~w tranzycie, segmentacja sieci,
    zarządzanie podatnościami i~aktualizacjami. \\[4pt]
Szkolenia i~świadomość
  & Regularne szkolenia dla zarządu i~pracowników; organy zarządzające
    obowiązane do zatwierdzania środków zarządzania ryzykiem i~odbycia
    wymaganych szkoleń. \\[4pt]
Kryptografia i~szyfrowanie
  & Wdrożenie polityk stosowania kryptografii i~szyfrowania w odniesieniu
    do przetwarzanych danych i~komunikacji. \\[4pt]
Bezpieczeństwo zasobów ludzkich
  & Procedury zarządzania dostępem i~polityki kontroli dostępu oparte
    na zasadzie najmniejszych uprawnień; bezpieczeństwo personelu. \\
\hline
\end{tabular}
\end{table}

Środki zarządzania ryzykiem muszą być zatwierdzone przez organ zarządzający podmiotu. W przypadku stwierdzenia niezgodności z wymaganiami dyrektywy,
podmiot zobowiązany jest do bezzwłocznego wdrożenia środków naprawczych.

\subsubsection{Zgłaszanie incydentów (art.~23 NIS2)}

Podmioty kluczowe i~ważne zobowiązane są do zgłaszania {poważnych incydentów} do właściwego zespołu CSIRT lub organu nadzorczego. Za poważny
incydent uznaje się zdarzenie, które powoduje lub może spowodować znaczne zakłócenia operacyjne lub straty finansowe, albo wpływa na inne osoby lub
organizacje powodując znaczne szkody.

\begin{table}[H]
\centering
\caption{Harmonogram obowiązkowego raportowania incydentów na mocy art.~23 NIS2}
\label{tab:nis2_raportowanie}
\begin{tabular}{>{\bfseries}p{4cm} p{9.5cm}}
\hline
Termin & Zakres informacji \\
\hline
Max. 24 h:  wczesne ostrzeżenie
  & Niezwłoczne poinformowanie o poważnym incydencie z~wskazaniem: czy
    incydent mógł być wywołany działaniem bezprawnym lub złym zamiarem;
    czy mógł wywrzeć wpływ transgraniczny. \\[4pt]
Max. 72 h:  zgłoszenie incydentu
  & Pełna aktualizacja wczesnego ostrzeżenia z~wstępną oceną incydentu,
    jego dotkliwości i~skutków oraz wskaźnikami integralności systemu.
    Dostawcy usług zaufania dokonują zgłoszenia w~ciągu 24~h. \\[4pt]
Na wniosek organu: sprawozdanie okresowe
  & Aktualizacje statusu w trakcie obsługi incydentu, stosownie do żądań
    CSIRT lub właściwego organu. \\[4pt]
Max. 1 miesiąc: sprawozdanie końcowe
  & Szczegółowy opis incydentu, rodzaj zagrożenia i~pierwotna przyczyna,
    zastosowane środki ograniczające ryzyko, transgraniczne skutki incydentu.
    Jeżeli incydent trwa - w~terminie miesiąca od jego zakończenia. \\
\hline
\end{tabular}
\end{table}

Podmioty mają również obowiązek poinformowania odbiorców usług o poważnych incydentach mogących negatywnie wpłynąć na świadczone usługi, a~także
o~poważnych cyberzagrożeniach i~środkach zaradczych, jakie mogą podjąć. Samo zgłoszenie incydentu nie nakłada na podmiot zgłaszający zwiększonej
odpowiedzialności prawnej.

\subsection{Nadzór i~egzekwowanie przepisów}

Dyrektywa NIS2 wprowadza wyraźne rozróżnienie reżimów nadzoru w zależności
od kategorii podmiotu.

\begin{table}[H]
\centering
\caption{Różnice w nadzorze nad podmiotami kluczowymi i~ważnymi}
\label{tab:nis2_nadzor}
\begin{tabular}{>{\bfseries}p{4.5cm} p{4.5cm} p{4.5cm}}
\hline
Aspekt & Podmioty kluczowe & Podmioty ważne \\
\hline
Charakter nadzoru
  & Nadzór ex ante (prewencyjny, ciągły) i~ex post
  & Nadzór ex post (reaktywny, po incydencie lub sygnale) \\[4pt]
Inspekcje i~audyty
  & Regularne inspekcje na miejscu, zdalne audyty, testy
  & Inspekcje na żądanie, po stwierdzeniu naruszeń \\[4pt]
Maks. kara administracyjna
  & Do 10~mln~EUR lub 2\% rocznego światowego obrotu
  & Do 7~mln~EUR lub 1,4\% rocznego światowego obrotu \\[4pt]
Odpowiedzialność osobista
  & Możliwy zakaz pełnienia funkcji kierowniczych przez osoby fizyczne
  & Możliwy zakaz pełnienia funkcji kierowniczych przez osoby fizyczne \\
\hline
\end{tabular}
\end{table}

\subsection{Odpowiedzialność zarządcza}

Jedną z~najistotniejszych zmian wprowadzanych dyrektywą NIS2 jest \textit{bezpośrednia odpowiedzialność organów zarządzających} za przestrzeganie wymogów cyberbezpieczeństwa. Oznacza to, że członkowie zarządu i~inne osoby pełniące funkcje kierownicze mogą być osobiście pociągnięci do odpowiedzialności
za poważne naruszenia przepisów dyrektywy. Organy zarządzające podmiotów zobowiązane są do: 
\begin{enumerate}[label=(\roman*)]
  \item zatwierdzania środków zarządzania ryzykiem,
  \item nadzoru nad ich wdrożeniem,
  \item odbywania regularnych szkoleń z~zakresu cyberbezpieczeństwa. 
\end{enumerate}

W~skrajnych przypadkach właściwy organ może zwrócić się do właściwego sądu o~tymczasowe zawieszenie osoby fizycznej
w~pełnieniu funkcji kierowniczej w~danym podmiocie.


\subsection{Rejestracja i~samoidentyfikacja}

Dyrektywa NIS2 wprowadza mechanizm {samoidentyfikacji}: to sam podmiot, a nie organ nadzorczy, ma obowiązek ustalić, czy mieści się w zakresie
stosowania dyrektywy, a następnie dokonać rejestracji u właściwego organu, przekazując dane adresowe, kontaktowe, sektor i~podsektor oraz wykaz państw
członkowskich, w których świadczy usługi. O~każdej zmianie tych danych podmiot zobowiązany jest do powiadomienia organu w terminie dwóch tygodni.
W polskiej implementacji dokonanej Krajową Ustawą o Cyberbezpieczeństwie\autocite{uksc2026} termin na zgłoszenie przedsiębiorstwa do wykazu
podmiotów kluczowych i~ważnych wynosi 6~miesięcy od wejścia w~życie ustawy, zaś na wdrożenie środków zarządzania ryzykiem 12~miesięcy.

\section{Regulacje związane z rynkiem finansowym}

Rynek finansowy w Polsce podlega nadzorowi Komisji Nadzoru Finansowego (\textit{KNF}). Do stycznia 2025~roku w kontekście bezpieczeństwa i~odporności cyfrowej,
podstawowym dokumentem obowiązującym banki była {Rekomendacja D KNF}\autocite{knf_rekomendacja_d_2013}.
KNF wydaje rekomendacje na podstawie art.~137 pkt~5 ustawy Prawo bankowe z dnia 29~sierpnia 1997~roku, Rekomendacje to dokumenty o charakterze nadzorczym,
działające na zasadzie \blockquote{zastosuj albo wyjaśnij}, które efiniuje i~opisują oczekiwania KNF wobec instytucji finansowych w zakresie
zarządzania ryzykiem, bezpieczeństwa IT oraz ochrony danych klientów. Formalnie nie są aktami prawnymi, w praktyce traktowane są przez nadzorowane podmioty jako  wymogi, ponieważ ich niestosowanie skutkuje negatywną oceną w ramach Badania i~Oceny Nadzorczej (BION)\autocite{knf_bion_2024}.

Rekomendacja D została przyjęta uchwałą KNF nr~7/2013 z dnia 8~stycznia 2013~roku, następnie  uchylona z dniem 17~stycznia 2025~r. i~zastąpiona
bezpośrednio stosowanym \textit{Rozporządzeniem Parlamentu Europejskiego i~Rady (UE) 2022/2554 z dnia 14~grudnia 2022~r. w sprawie operacyjnej odporności cyfrowej sektora finansowego i~zmieniającym rozporządzenia (WE) nr~1060/2009, (UE) nr~648/2012, (UE) nr~600/2014, (UE) nr~909/2014
oraz (UE) 2016/1011}, powszechnie zwanym DORA (\textit{Digital Operational Resilience Act}) \autocite{dora2022}.

Wytyczne i~dobre praktyki zawarte w Rekomendacji~D zostały w większości zachowane i~przeniesione do DORA oraz towarzyszących jej regulacyjnych
standardów technicznych (RTS) i~wykonawczych standardów technicznych (ITS) opracowanych przez Europejski Urząd Nadzoru Bankowego (EBA).

\subsection{Prawo bankowe}
Ustawa z dnia 29 sierpnia 1997 r. - {Prawo bankowe} to akt prawny uchwalony przez Sejm RP 29 sierpnia 1997~r., opublikowana w Dzienniku Ustaw 1997 nr 140 poz. 939, weszła w życie 1 stycznia 1998~roku. Od tamtej pory była wielokrotnie nowelizowana, aktualny tekst jednolity obejmuje zmiany wprowadzone m.in. przez przepisy implementujące dyrektywy unijne (CRD IV, PSD2, AML).

Art. 104 ustawy Prawo bankowe nakłada na bank, jego pracowników oraz wszystkich pośredników wykonujących czynności bankowe bezwzględny obowiązek zachowania tajemnicy bankowej, obejmuje ona wszelkie informacje o czynnościach bankowych uzyskane podczas negocjacji, zawierania i~realizacji umowy z klientem, a więc m.in. salda rachunków, historię transakcji, dane o kredytach, zabezpieczeniach czy zdolności kredytowej. Ujawnienie tych informacji osobom trzecim jest dopuszczalne wyłącznie za pisemną zgodą klienta lub w ściśle określonych przypadkach przewidzianych przez prawo (np. na żądanie prokuratury lub organów podatkowych). W praktyce oznacza to, że każdy podmiot zewnętrzny, w tym dostawca chmury czy firma technologiczna któremu bank powierza przetwarzanie danych, przejmuje obowiązek zachowania tajemnicy bankowej i~musi gwarantować jej nienaruszalność kontraktowo oraz technicznie.

\section{Rozporządzenie DORA: cyfrowa odporność operacyjna sektora finansowego}

Rozporządzenie Parlamentu Europejskiego i~Rady (UE) 2022/2554, znane jako {DORA} (\textit{Digital Operational Resilience Act}), weszło w~życie 16~stycznia 2023~roku, a~jego stosowanie stało się obowiązkowe od 17~stycznia 2025~roku.
Rozporządzenie stanowi odpowiedź na rosnącą zależność sektora finansowego od technologii informacyjno-komunikacyjnych (ICT) oraz systematyczny wzrost liczby cyberincydentów dotykających instytucje finansowe w~całej Unii Europejskiej~\autocite{dora2022}.

Rozporządzenie strukturyzuje obowiązki wokół pięciu wzajemnie powiązanych filarów, opisanych szczegółowo poniżej.

\subsection*{Filar I, Zarządzanie ryzykiem ICT}

Pierwszy filar zobowiązuje instytucje finansowe do ustanowienia kompleksowych i~udokumentowanych ram zarządzania ryzykiem ICT, za które odpowiedzialność ponosi
bezpośrednio kierownictwo wyższego szczebla~\autocite{dora2022}. Ramy te muszą obejmować szczegółowy rejestr wszystkich aktywów ICT wspierających funkcje krytyczne, polityki bezpieczeństwa informacji, mechanizmy ciągłego monitorowania anomalii oraz plany ciągłości działania i~odtwarzania po awarii. Podmiot finansowy zobowiązany jest do regularnego przeglądu i~aktualizacji tych ram, a~organ nadzoru może w~każdej chwili
zażądać ich udostępnienia~\autocite{resilia_dora}.

\subsection*{Filar II, Zarządzanie incydentami ICT i~ich raportowanie}

Drugi filar wprowadza ujednolicone, rygorystyczne reguły klasyfikowania, rejestrowania i~zgłaszania incydentów związanych z~ICT. W~przypadku wykrycia
poważnego incydentu podmiot finansowy zobowiązany jest do przekazania wstępnego powiadomienia właściwemu organowi nadzorczemu w~ciągu 4~godzin,
raportu pośredniego w~ciągu 72~godzin oraz raportu końcowego w~ciągu miesiąca.
Naruszenie bezpieczeństwa ICT prowadzące do wycieku danych osobowych podlega równolegle obowiązkowi zgłoszenia do UODO na podstawie art.~33 RODO, co~czyni z~DORA i~RODO komplementarne reżimy sprawozdawcze.

\subsection*{Filar III, Testowanie cyfrowej odporności operacyjnej}

Trzeci filar nakłada obowiązek regularnego, ustrukturyzowanego testowania odporności systemów ICT na zagrożenia. Co najmniej raz w~roku podmioty finansowe
muszą przeprowadzać testy obejmujące m.in.\ oceny podatności, testy scenariuszowe i~testy penetracyjne. Duże, systemowo istotne instytucje zobowiązane są do wykonywania zaawansowanych testów penetracyjnych sterowanych zagrożeniami (\textit{Threat-Led Penetration Testing, TLPT}) z~udziałem zewnętrznych certyfikowanych podmiotów, przy czym zewnętrzny dostawca ICT musi w~pełni współpracować przy ich przeprowadzaniu.

\subsection*{Filar IV, Zarządzanie ryzykiem zewnętrznych dostawców ICT}

Czwarty filar stanowi najbardziej przełomowy element rozporządzenia. DORA nakłada bezpośrednie obowiązki na zewnętrznych dostawców ICT, w~tym hiperscalerów chmurowych takich jak Google Cloud, Microsoft Azure czy Amazon Web Services, obejmując ich nadzorem jako dostawców kluczowych (ang. \textit{Critical ICT Third-Party Providers
(CTPP)}). Instytucje finansowe muszą przeprowadzać ocenę ryzyka przed zawarciem każdej umowy z~dostawcą ICT, a~sama umowa musi zawierać postanowienia
o~prawie do audytu, lokalizacji przetwarzania danych, planach wyjścia (ang. \textit{exit plan}) oraz minimalnych okresach przejściowych przy zmianie dostawcy~\autocite{eiopa_dora}.
Korzystanie z~niezatwierdzonych narzędzi chmurowych (ang. \textit{shadow IT}) stanowi naruszenie tego filaru, a~jednocześnie naruszenie zasady rozliczalności z~art.~5 ust.~2 RODO.

\subsection*{Filar V, Wymiana informacji o~cyberzagrożeniach}

Piąty filar zachęca, a~w~określonych przypadkach zobowiązuje, podmioty finansowe do uczestnictwa w~schematach wymiany informacji o~zagrożeniach cybernetycznych, podatnościach i~technikach ataku w~ramach zaufanych społeczności sektorowych~\autocite{eiopa_dora}.
Celem jest budowanie zbiorowej odporności sektora: wczesne ostrzeżenie jednej instytucji o~nowym wektorze ataku pozwala innym podmiotom na podjęcie środków zapobiegawczych zanim incydent zdąży się rozprzestrzenić. Uczestnictwo w~takich schematach musi jednak odbywać się zgodnie z~przepisami o~ochronie danych osobowych oraz z~zachowaniem tajemnicy zawodowej~\autocite{resilia_dora}.

\subsection*{Sankcje}

Za naruszenie rozporządzenia Wiodący Organ Nadzorczy (ang \textit{Lead Overseer}) może nakładać na kluczowych dostawców ICT kary w~wysokości do 1\% średniego dziennego światowego obrotu za każdy dzień naruszenia, przez okres do sześciu miesięcy. Dla podmiotów finansowych nadzorowanych przez KNF sankcje określa prawo krajowe wdrażające uprawnienia nadzorcze wynikające z~rozporządzenia~\autocite{knf_dora}.

Wiodący Organ Nadzorczy to instytucja wyznaczona na poziomie UE do bezpośredniego nadzoru nad kluczowymi dostawcami ICT  w ramach rozporządzenia DORA.
Rolę Wiodącego Organu Nadzorczego pełni jedna z trzech unijnych instytucji:
\begin{enumerate}[label=(\roman*)]
  \item EBA (European Banking Authority) - dla dostawców ICT obsługujących głównie sektor bankowy,
\item ESMA (European Securities and Markets Authority) - dla dostawców obsługujących rynki kapitałowe,
\item EIOPA (European Insurance and Occupational Pensions Authority) - dla dostawców obsługujących ubezpieczenia i fundusze emerytalne.
\end{enumerate}
Przypisanie konkretnego dostawcy do jednego z tych organów zależy od tego, jaki typ podmiotów finansowych jest jego głównym klientem.
Wiodący Organ Nadzorczy ma szerokie uprawnienia wobec krytycznych dostawców ICT: może żądać informacji i dokumentacji, przeprowadzać inspekcje, wydawać zalecenia naprawcze oraz nakładać kary finansowe.

\section{Ochrona danych medycznych, regulacje prawne i~nadzór}

Dane dotyczące zdrowia należą do szczególnej kategorii danych osobowych w~rozumieniu art.~9 rozporządzenia (UE) 2016/679 (\textit{RODO})\autocite{rodo}, co oznacza, że ich przetwarzanie jest co~do zasady zakazane, a~wyjątki od tego zakazu są ściśle i~enumeratywnie określone w~przepisach~\autocite{przewodnik_rodo_zdrowie}. Do kategorii tej zalicza się wszelkie informacje o~stanie zdrowia fizycznego lub psychicznego osoby, dane genetyczne i~biometryczne, a~także informacje pośrednie (takie jak numer recepty, historia wizyt lekarskich czy przynależność do grupy ubezpieczenia zdrowotnego)  jeżeli pozwalają one na wyciągnięcie wniosków o~stanie zdrowia danej osoby.
Podwyższony reżim ochrony danych zdrowotnych uzasadniony jest ich szczególną wrażliwością: ujawnienie informacji o~chorobie, leczeniu psychiatrycznym czy uzależnieniu może prowadzić do dyskryminacji w~zatrudnieniu, wykluczenia z~rynku ubezpieczeń lub stygmatyzacji społecznej.

\subsection*{Krajowe ramy prawne i~wielość nadzorców}

Ochrona danych medycznych w~Polsce regulowana jest przez trzy nakładające się
warstwy prawa, za których egzekwowanie odpowiadają odrębne organy nadzorcze.

\begin{enumerate}[label=(\roman*)]
  \item Pierwszą warstwę stanowi RODO wraz z~krajową ustawą o~ochronie danych osobowych z~2018~roku. Nadzór nad jej przestrzeganiem sprawuje Prezes Urzędu Ochrony Danych Osobowych (UODO), uprawniony do przeprowadzania kontroli w~placówkach medycznych, wydawania decyzji nakazowych i~nakładania kar administracyjnych sięgających 20~milionów euro lub 4\% rocznego obrotu. Każda placówka medyczna przetwarzająca dane zdrowotne na dużą skalę, zobowiązana jest do wyznaczenia Inspektora Ochrony Danych (IOD), który monitoruje zgodność przetwarzania z~RODO, szkoli personel i~pełni rolę punktu kontaktowego dla pacjentów oraz organu nadzoru~\autocite{przewodnik_rodo_dane_zdrowotne}.

  \item Drugą warstwę tworzy ustawa z~dnia 6~listopada 2008~roku o~prawach pacjenta i~Rzeczniku Praw Pacjenta. Nadzór nad prawem pacjenta do dostępu do dokumentacji medycznej (art.~23--29 ustawy) sprawuje Rzecznik Praw Pacjenta, który rozpatruje skargi pacjentów i~może kierować zalecenia do podmiotów leczniczych. Ustawa nakłada na każdy podmiot udzielający świadczeń zdrowotnych obowiązek prowadzenia, przechowywania i~udostępniania dokumentacji medycznej w~ściśle określony sposób, a~dokumentacja ta podlega ochronie zarówno na podstawie ustawy, jak i~RODO~\autocite{rpp_dokumentacja}.

  \item Trzecią warstwę stanowi system akredytacji i~nadzoru nad podmiotami leczniczymi, który sprawuje {Ministerstwo Zdrowia} we współpracy z~\textit{Narodowym Funduszem Zdrowia} oraz, w~zakresie standardów informatycznych, {Centrum e-Zdrowia}, odpowiadające za systemy P1. Systemy P1 to centralny system e-zdrowia w Polsce, formalnie określany jako \textit{Elektroniczna Platforma Gromadzenia, Analizy i~Udostępniania Zasobów Cyfrowych o Zdarzeniach Medycznych}. Jest to państwowa platforma teleinformatyczna rozwijana przez Centrum e-Zdrowia, służąca do gromadzenia, przetwarzania i~udostępniania danych o zdarzeniach medycznych oraz do obsługi usług takich jak e-recepta, e-skierowanie i~wymiana EDM ~\autocite{sc_prawa_pacjenta}.
\end{enumerate}  

\subsection*{Europejska Przestrzeń Danych Zdrowotnych (EHDS)}

Nową regulacją jest \textit{Europejska Przestrzeń Danych Zdrowotnych} (\textit{European Health Data Space, EHDS}), której rozporządzenie weszło w~życie 26~marca 2025~roku~\autocite{ehds_gov_zdrowie}. EHDS tworzy
ogólnounijne ramy pierwotnego i~wtórnego wykorzystania elektronicznych danych zdrowotnych: pacjent uzyskuje prawo dostępu do swoich danych w~standardowym europejskim formacie (\textit{EEHRxF}), możliwość kontrolowania kto ma do nich wgląd oraz prawo do ograniczenia dostępu do wybranych części dokumentacji~\autocite{ehds_komisja}.
Pełne wdrożenie systemu  przewidziane jest do 26~marca 2027~roku~\autocite{ehds_gov_zdrowie}.

EHDS działa równolegle do RODO, nie zastępując go, lecz nadając sektorowi zdrowotnemu dodatkowy, sektorowy wymiar regulacyjny. Oznacza to, że placówki medyczne będą podlegać jednocześnie: RODO (nadzór UODO), ustawie o~prawach pacjenta (nadzór Rzecznika Praw Pacjenta) oraz EHDS (nadzór krajowego organu ds.\ elektronicznych danych zdrowotnych, który Polska wyznaczy do 2027~roku).
Naruszenie bezpieczeństwa systemu ICT przechowującego dane medyczne będzie tym samym incydentem, który należy raportować do wielu organów jednocześnie, analogicznie do zbieżności DORA i~RODO w~sektorze finansowym.


\chapter{Konflikt jurysdykcji i~suwerenność danych}

Informacja jest współczesną walutą, jest to bardzo popularne, jednocześnie prawdziwe stwierdzenie. Dane, które posiadają przedsiębiorstwa, uczelnie oraz rządy stanowią o sile gospodarczej i~przewadze technologicznej. Nic więc dziwnego w działaniach, które mają na celu ochronę informacji przed nieautoryzowanym dostępem.
Do szpiegostwa gospodarczego ze strony takich państw jak Rosja czy Chiny dochodziło od wielu lat \autocite{putter2023espionage}, ale nowością jest wykorzystywanie informacji i~technologii przeciwko sojusznikom, gdy decyzje podejmowane nie podobają się jednej ze stron.

Flagowym przykładem była sytuacja związana z działaniami administracji amerykańskiej w stosunku do Międzynarodowego Trybunału Karnego w Hadze, a źródłem ich była decyzja sądu, której treść nie była zgodna z oczekiwaniami administracji Donalda Trumpa.

Sankcje USA wobec MTK przełożyły się na odcięcie dostępu do niektórych usług amerykańskich, w szczególności do usług Microsoft w odniesieniu do prokuratora Karima Khana, a także do ograniczeń finansowych i~wizowych. Podstawę tych działań stanowił nakaz (\textit{executive order}) USA z 5 lutego 2025 roku, przewidujący m.in. blokowanie majątku i~aktywów oraz restrykcje wizowe wobec urzędników MTK i~ich rodzin. \autocite{whitehouse_icc_sanctions_2025}

Jednocześnie sankcje zostały rozszerzone w czerwcu 2025 na kolejne osoby, Reine Adelaide Sophie Alapini Gansou, sędzia MTK i~druga wiceprezes MTK, oraz sędziowe Solomy Balungi Bossa, Luz del Carmen Ibáñez Carranza oraz Beti Hohler, sędzia MTK.

Decyzja zaostrza napięcie między USA a państwami popierającymi porządek oparty na prawie międzynarodowym. Najważniejszy efekt polityczny to osłabienie niezależności Trybunału przez pokazanie. Ta decyzja buduje precedens, że państwo może używać sankcji nie tylko wobec przeciwników geopolitycznych, lecz także wobec międzynarodowego sądu i~osób współpracujących z nim. Taki precedens może zachęcać inne rządy do podobnych działań wobec instytucji międzynarodowych, gdy ich decyzje staną się politycznie niewygodne.

Incydent związany z odcięciem amerykańskich usług cyfrowych dla osób objętych sankcjami USA nie wydaje się stanowić jedynej przyczyny działań UE w obszarze suwerenności danych, lecz stał się katalizatorem politycznym wzmacniającym unijną debatę o ryzykach wynikających z zależności od dostawców podlegających jurysdykcji Stanów Zjednoczonych.

\section{Problem CLOUD Act i~FISA 702}

Podstawą konfliktu prawnego w zakresie przetwarzania i~wykorzystywania danych w chmurze oraz dostępu do nich przez władze Stanów Zjednoczonych są dwie amerykańskie ustawy: pierwsza to \textit{Clarifying Lawful Overseas Use of Data Act} (CLOUD Act, 2018)  \autocite{cloudact2018}, druga to sekcja 702 ustawy \textit{Foreign Intelligence Surveillance Act} (FISA 702) \autocite{fisa702}. Obie regulacje pozwalają władzom amerykańskim na żądanie dostępu do danych przechowywanych przez przedsiębiorstwa podlegające pod amerykańską jurysdykcję, bez względu na to, gdzie dane te się znajdują.

Amerykański ustawodawca zdefiniował jurysdykcję w CLOUD Act, określając, że dotyczy ona podmiotu, a nie danych. Oznacza to, że Amazon, Microsoft i~Google, które łącznie kontrolują około 70 procent europejskiego rynku usług w chmurze \autocite{synergy2025eu}, mogą zostać zobowiązane nakazem sądowym do ujawnienia danych swoich klientów, nawet jeśli serwery znajdują się na terytorium Unii Europejskiej \autocite{edpb2019cloudact}. Co istotne, ani osoba, której dane
dotyczą, ani też żaden europejski organ nadzoru nie muszą być powiadamiani. Ustawa FISA~702 działa równolegle, upoważniając agencje wywiadowcze, bez nakazu, do gromadzenia komunikacji osób spoza Stanów Zjednoczonych.


\section{Suwerenność danych: Gaia-X, suwerenne chmury obliczeniowe i~poufne przetwarzanie danych}

Europejską odpowiedzią na strukturalną zależność od dostawców spoza Unii Europejskiej jest opracowanie koncepcji {suwerenności danych cyfrowych}. Składa się ona z trzech wymiarów: rezydencja danych (ang. \textit{data residency}), suwerenność prawna (ang. \textit{data sovereignty}) oraz  suwerenność operacyjna (ang. \textit{operational sovereignty}). Kluczowym problemem jest fakt, że
fizyczna lokalizacja serwerów na terytorium Unii Europejskiej nie przesądza o jurysdykcji  prawnej. Stąd UE podjęła szereg inicjatyw: inicjatywę GAIA-X, suwerenne chmury, poufne przetwarzanie.

\subsection*{Gaia-X}
Inicjatywa Gaia-X, ogłoszona przez Niemcy i~Francję w 2019~r. i~formalnie zainicjowana przez Komisję Europejską w 2020 roku, ma na celu stworzenie interoperacyjnej, zdecentralizowanej infrastruktury danych zgodnej z europejskimi normami prawnymi i~wartościami \autocite{braud2021gaiax}. Jego główne cele to: zapewnienie użytkownikom pełnej kontroli nad swoimi danymi i~ich wykorzystaniem, standaryzacja usług chmurowych w oparciu o otwarte standardy oraz
stworzenie alternatywy dla dominujących trzech amerykańskich hiperskalerów: Amazon, Google oraz Microsoft \autocite{braud2021gaiax}.

Realizacja projektu jest kontrowersyjna. Krytycy zauważają, że do konsorcjum Gaia-X dołączyli sami hiperskalerzy, Amazon, Google i~Microsoft, co rodzi pytania i~wątpliwości o faktyczną suwerenność usług w ramach inicjatywy.

\subsection*{Operator Chmury Krajowej (OChK)}
Polska inicjatywa, czyli Operator Chmury Krajowej (OChK), jest poddawana podobnej krytyce jak Gaia-X. OChK jest spółką z ograniczoną odpowiedzialnością, w której po 50\% udziałów posiadają PKO Bank Polski oraz Polski Fundusz Rozwoju (PFR), oba podmioty z dominującym udziałem Skarbu Państwa. Spółka została powołana w celu utworzenia suwerennego, krajowego środowiska chmurowego
dla polskiego sektora publicznego i~finansowego.

W roku 2019 OChK zawarł strategiczne partnerstwo z Google Cloud. W jego ramach OChK stał się certyfikowanym sprzedawcą usług Google w Polsce, a Google otworzył region chmurowy w Warszawie. Partnerstwo jest analogiczne do krytykowanego wejścia amerykańskich hiperskalerów do Gaia-X: OChK, który miał stanowić polską alternatywę dla dostawców spoza UE, stał się pośrednikiem sprzedaży ich usług.

OChK w październiku 2025 roku ogłosiło, że platforma \blockquote{osiągnęła pełną suwerenność cyfrową zgodnie z definicją Cloud Sovereignty Framework Komisji Europejskiej} i~że usługi podlegają \blockquote{całkowitej kontroli podmiotów z UE oraz unijnemu prawu, bez krytycznych zależności od dostawców spoza Europy}\autocite{ochk2025sovereignty}. Twierdzenie to wymaga jednak zastrzeżenia: dotyczy ono własnej warstwy platformowej OChK, natomiast nie eliminuje faktu, że klienci korzystający za pośrednictwem OChK z usług Google Cloud w dalszym ciągu korzystają z infrastruktury podmiotu podlegającego prawu USA, w tym CLOUD Act i~FISA~702. W praktyce OChK umożliwia dwa różne modele, usługi własne (suwerenne) oraz usługi Google dystrybuowane przez OChK, które są niesuwerenne w sensie prawnym, i~to od odbiorcy usługi i~wybranej przez niego konfiguracji zamówienia oraz świadomości użytkownika zależy, z którego modelu faktycznie korzysta.

\subsection*{Suwerenne chmury}

Koncepcja suwerennej chmury odnosi się do oferty usług chmurowych, która gwarantuje, że dane pozostają pod prawną kontrolą jurysdykcji europejskiej. Jedyną drogą do prawdziwej suwerenności danych jest wybór dostawcy, który nie ma żadnych prawnych powiązań z jurysdykcją USA, czyli europejskich dostawców usług chmurowych całkowicie niezależnych od amerykańskich korporacji.

\subsection*{Przetwarzanie poufne}

Szyfrowanie danych w spoczynku i~podczas przesyłania to wymóg stawiany przed przetwarzaniem poufnym (ang.\ \textit{confidential computing}). Jego bazę stanowią \textit{Zaufane Środowiska Wykonawcze} (TEE), sprzętowo izolowane środowiska, w których dane pozostają zaszyfrowane w czasie aktywnego przetwarzania, niedostępne ani dla hiperwizora, ani systemu operacyjnego, ani tym bardziej dla personelu dostawcy chmury \autocite{ijaibdcms2024tee}.

TEE chronią przed nieautoryzowanym i~pozajurysdykcyjnym dostępem do danych podczas ich przechowywania, transmisji oraz podczas obliczeń \autocite{ijaibdcms2024tee}.

Warto podkreślić, że takie podejście wpisuje się bezpośrednio w oczekiwania polskiego i~europejskiego nadzoru finansowego. Komisja Nadzoru Finansowego już w Komunikacie Chmurowym z 2020 roku wymagała od podmiotów nadzorowanych, by informacje prawnie chronione były szyfrowane zawsze w spoczynku (ang. \textit{at rest}) oraz w trakcie transferu (ang. \textit{in transit}) \autocite{knf2020chmura}, a kontrola nad kluczami
kryptograficznymi pozostawała po stronie instytucji finansowej, nie dostawcy chmury.
Obowiązujące od stycznia 2025~r. rozporządzenie DORA idzie o krok dalej, nakładając wymóg kompleksowej polityki szyfrowania obejmującej dane \textit{at rest}, \textit{in transit}, a także w trakcie użycia (ang. \textit{in use}), co jest bezpośrednim odniesieniem do scenariusza, w którym przetwarzanie poufne staje się nie tyle opcją, ile techniczną odpowiedzią na regulacyjne oczekiwania \autocite{dora2022}\autocite{dora_rts2024}.

Z perspektywy jurysdykcyjnej, szczególnie wynikającej z ryzyk CLOUD Act, poufne obliczenia umożliwiają techniczne egzekwowanie reguł dotyczących lokalizacji przetwarzania, operacje mogą być wykonywane wyłącznie w autoryzowanych geograficznie lokalizacjach ({powinowactwo lokalizacji}) i~na dedykowanych hostach ({powinowactwo hosta}) \autocite{ijaibdcms2024tee}. Co ważne, choć sama technologia nie uchyla prawnych obowiązków dostawcy wynikających z amerykańskiej ustawy CLOUD Act, to utrudnia realizację ewentualnych nakazów dostępu, a przede wszystkim eliminuje ryzyko cichego dostępu
operacyjnego ze strony dostawcy.


\section{Schrems I, II i~konsekwencje dla transferu danych poza EOG}

Rozwiązania prawne regulujące zakres transferu danych osobowych między Unią Europejską, a Stanami Zjednoczonymi stanowią jeden z najważniejszych punktów w historii  prawa o ochronie danych. Z perspektywy oceny ryzyka łańcuchów dostaw ICT, stanowi ona również  przykład tego, jak strukturalna luka prawna może  podważyć zgodność z przepisami tysięcy organizacji jednocześnie. Sekwencja zdarzeń zainicjowana przez jednego austriackiego studenta prawa doskonale ilustruje kategorię ryzyka, którą system ORION ma wykrywać i~mierzyć.

\subsection{Umowa Safe Harbour (2000–2015)}

Podstawowym instrumentem regulującym transfer danych osobowych między UE a USA była decyzja Komisji 2000/520/WE, powszechnie znana jako \textit{Umowa Safe Harbour}. Przyjęta 26 lipca 2000 roku ustawa ustanowiła mechanizm autocertyfikacji, w ramach którego amerykańskie firmy mogły dobrowolnie zobowiązać się do przestrzegania siedmiu zasad ochrony danych: powiadomienia, wyboru, dalszego przekazywania, bezpieczeństwa, integralności danych, dostępu i~egzekwowania. Po dokonaniu autocertyfikacji uzyskiwały prawo do otrzymywania danych osobowych z państw członkowskich UE bez dalszego uzasadnienia prawnego~\autocite{EuropeanCommission2000SafeHarbour}.

Mechanizm działał na podstawie oświadczeń składanych w Departamencie Handlu USA, a ramy dostępu do danych opierały się wyłącznie na dobrej wierze organizacji uczestniczących i~teoretycznej możliwości egzekwowania zasad przez Federalną Komisję Handlu USA. W szczytowym okresie ponad 4500 amerykańskich firm, w tym Facebook, Google, Microsoft i~Amazon, działało na podstawie certyfikatu Safe Harbour, który stanowił podstawę prawną dla zdecydowanej większości komercyjnych przepływów danych między UE a USA.\autocite{EuropeanCommission2000SafeHarbour}.

\subsection{Maximilian Schrems: Tło i~strategia prawna}

Maximilian Schrems jest austriackim prawnikiem, działaczem na rzecz ochrony prywatności i~założycielem organizacji pozarządowej \textit{noyb} (None of Your Business), zajmującej się egzekwowaniem europejskich przepisów o ochronie danych. W wyniku działania Schremsa sprawa trafiła do Trybunału Sprawiedliwości Unii Europejskiej (TSUE) z wnioskiem o wydanie orzeczenia w trybie prejudycjalnym.

\subsection{Schrems I: Unieważnienie zasad Safe Harbor}

6 października 2015 roku TSUE wydał wyrok w sprawie C-362/14, \textit{Maximillian Schrems przeciwko Komisarzowi Ochrony Danych}, powszechnie znanej jako \textit{Schrems~I}~\autocite{CJEU2015SchremsI}. Trybunał orzekł nieważność decyzji Komisji 2000/520/WE ze skutkiem natychmiastowym. Orzeczenie opierało się na dwóch
zasadniczych podstawach. 
Po pierwsze, Trybunał orzekł, że pierwszeństwo przyznane wymogom bezpieczeństwa narodowego, interesu publicznego i~egzekwowania prawa USA nad zasadami Safe Harbor oznaczało, że amerykańskie firmy były systematycznie uprawnione, a w praktyce zmuszone, do odchodzenia od tych zasad, ilekroć wymagało tego prawo amerykańskie. Ramy, które mogą zostać uchylone przez suwerenne prawo państwa odbiorcy, nie mogą stanowić odpowiedniego poziomu ochrony w rozumieniu dyrektywy o ochronie danych.

Po drugie, Trybunał stwierdził, że obywatele UE, których dane zostały przekazane do Stanów Zjednoczonych, nie mieli skutecznego środka odwoławczego ani administracyjnego, aby egzekwować swoje prawa wobec władz USA. Brak jakiegokolwiek mechanizmu umożliwiającego osobom z UE, których dane dotyczą,
uzyskanie dostępu, poprawienie lub usunięcie danych przechowywanych przez amerykańskie agencje wywiadowcze, był nie do pogodzenia z podstawowym prawem do skutecznej ochrony sądowej zapisanym w artykule 47 Karty praw podstawowych UE.

\subsection{Schrems II: Tarcza Prywatności}
W odpowiedzi na wyrok w sprawie Schrems I, Komisja Europejska i~Departament Handlu USA wynegocjowały nowe ramy prawne, Tarczę Prywatności UE–USA, która weszła w życie 12 lipca 2016 roku.

Schrems natychmiast zakwestionował Tarczę Prywatności, argumentując, że nie rozwiązała ona fundamentalnych braków strukturalnych zidentyfikowanych w sprawie Schrems I. 16 lipca 2020 r. TSUE potwierdził tę ocenę w sprawie C-311/18, \textit{Data Protection Commissioner przeciwko Facebook Ireland Limited i~Maximillian Schrems}, zwanej dalej \textit{Schrems~II}~\autocite{CJEU2020SchremsII}. Sąd unieważnił Tarczę Prywatności z zasadniczo tych samych powodów, co Safe Harbor: amerykańskie prawo dotyczące nadzoru (w szczególności art. 702 ustawy o nadzorze nad wywiadem zagranicznym (FISA)) oraz rozporządzenie wykonawcze 12333, nadal zezwalały na masowe gromadzenie danych obywateli UE bez odpowiednich ograniczeń, niezależnego nadzoru lub skutecznych mechanizmów dochodzenia roszczeń.

Sprawa Schrems~II miała jednak szersze znaczenie niż samo unieważnienie Tarczy Prywatności. Trybunał wydał również niezwykle istotne wytyczne dotyczące standardowych klauzul umownych (SCC), orzekając, że eksporterzy i~importerzy danych powołujący się na standardowe klauzule umowne są zobowiązani do oceny, w każdym indywidualnym przypadku, czy prawo kraju przeznaczenia zapewnia odpowiednią ochronę w praktyce, a nie tylko w teorii. Ten obowiązek, Ocena Skutków Przenoszenia (TIA), stał się nowym, obowiązkowym elementem zgodności dla każdej organizacji przekazującej dane osobowe poza EOG i~jest on również istotnym aspektem w modelu ORION.

\subsection{Ramy ochrony prywatności danych (2023–obecnie)}

Po trzech latach intensywnych negocjacji Komisja Europejska przyjęła 10 lipca 2023 roku \textit{Ramy ochrony prywatności danych UE–USA (DPF)} w formie \textit{Decyzji wykonawczej Komisji (UE) 2023/1795}~\autocite{KomisjaEuropejska2023DPF}. DPF wprowadził nowy mechanizm odwoławczy: Sąd ds. Ochrony Danych (DPRC), mający na celu zapewnienie obywatelom UE wiążącego środka odwoławczego w przypadku bezprawnego gromadzenia informacji wywiadowczych. 

DPF stanowi obecnie główną podstawę prawną przekazywania danych osobowych między UE a USA. Jednakże już Maximillian Schrems  zapowiedział dalsze postępowanie sądowe kwestionujące DPF przed sądami europejskimi, argumentując, że DPRC nie stanowi niezależnego organu sądowego w rozumieniu wymaganym przez orzecznictwo TSUE i~że sekcja 702 ustawy FISA pozostaje niezmieniona w swoich najbardziej problematycznych  aspektach. 

\subsection{Implikacje dla struktury ORION}

Wyroki bedące skutkiem działań Schremsa niosą bezpośrednie implikacje dla metodologii ORION. Decyzje sądów wskazują, że ryzyko regulacyjne związane z geograficznym pochodzeniem dostawcy ICT jest problemem realnym, który musi być przez dostawców usług ICT brany pod uwagę.

W modelu ORION ryzyko to jest wyliczane. Dostawcy geograficznie lokalizowani poaza EOG uruchamiają obowiązkowy zestaw pytań audytowych obejmujących aktualną podstawę prawną transferu danych, istnienia \textit{Oceny Skutków Transferu (\textit{TIA})} oraz weryfikacji postanowień umów regulujących wnioski rządowe o dostęp do danych. 
Historia Safe Harbor, Privacy Shield i~ciągła niepewność związana z DPF stanowią zatem empiryczne uzasadnienie dla jednej z kluczowych decyzji projektowych ORION: klasyfikacji pochodzenia geograficznego jako ryzka kluczowego.

\section{Strategia legislacyjna UE: Data Act, EUCS i~CADA}

Europejska odpowiedź na ryzyko jurysdykcyjne ewoluuje z poziomu inicjatyw przemysłowych (Gaia-X) w~stronę wiążących instrumentów prawnych, tworzących zalążek spójnego ekosystemu suwerenności cyfrowej. Trzy wzajemnie uzupełniające
się inicjatywy: \textit{EU Data Act}, \textit{EUCS} oraz \textit{Cloud and AI Development Act} adresują kolejno: prawa dostępu do danych i~ich przenoszalność, certyfikację bezpieczeństwa dostawców chmury oraz budowę europejskiej infrastruktury obliczeniowej.

\subsection*{EU Data Act (Rozporządzenie 2023/2854)}

\textit{Rozporządzenie Parlamentu Europejskiego i~Rady (UE) 2023/2854}, powszechnie znane jako \textit{EU Data Act}, zostało opublikowane w~Dzienniku Urzędowym UE 22~grudnia 2023~r.\ i~stało się w~pełni stosowane od 12~września 2025~r.~\autocite{dataact2023}. Regulacja wprowadza dwa mechanizmy bezpośrednio adresujące ryzyko jurysdykcyjne:

\begin{enumerate}[label=(\roman*)]
    \item Po pierwsze, art.~31 Data Act nakłada na dostawców usług przetwarzania danych obowiązek wdrożenia środków technicznych i~organizacyjnych uniemożliwiających
    organom państw spoza Europejskiego Obszaru Gospodarczego dostęp do danych  przechowywanych w~UE bez ważnej podstawy w~prawie międzynarodowym lub
    w~umowie między UE a~danym państwem trzecim~\autocite{dataact2023}. Przepis ten nie uchyla obowiązków dostawcy wynikających bezpośrednio z~prawa
    państwa jego inkorporacji (takich jak CLOUD Act), lecz nakłada na dostawcę
    obowiązek aktywnego kwestionowania wniosków uznanych za sprzeczne z~prawem
    unijnym i~powiadamiania klienta.

    \item Po drugie, Data Act zakazuje z~dniem 12~stycznia 2027~r.\ pobierania opłat za przełączenie dostawcy (\textit{switching fees}) i~ekstrakcję danych, co~bezpośrednio ogranicza ryzyko uzależnienia od jednego hiperskalera (\textit{vendor lock-in}) identyfikowane przez EBA i~Parlament Europejski \autocite{dataact2023}.
\end{enumerate}

\subsection*{Europejski Schemat Certyfikacji Bezpieczeństwa Chmury (EUCS)}

Europejski Schemat Certyfikacji Bezpieczeństwa Chmury (\textit{European
Cybersecurity Certification Scheme for Cloud Services , EUCS}), opracowywany
przez Agencję Unii Europejskiej ds.\ Cyberbezpieczeństwa (\textit{ENISA}) na podstawie
unijnego \textit{Aktu o~Cyberbezpieczeństwie (Rozporządzenie 2019/881)}, definiuje
trzy poziomy zapewnienia: \textit{Basic}, \textit{Substantial} i~\textit{High}
\autocite{enisa_eucs2022}. Projekt przewidywał również poziom
\textit{High+}, który wprowadzałby wymogi jurysdykcyjne, rezydencję
danych, personel operacyjny z~EOG i~brak wiążących powiązań korporacyjnych
z~podmiotami spoza EOG, jako warunek certyfikacji dla usług krytycznych
\autocite{enisa_eucs2022}.

Brak politycznego porozumienia wśród państw członkowskich co do poziomu
\textit{High+} jest jednym z~głównych czynników blokujących wdrożenie
suwerennej oferty chmurowej dla sektora finansowego i~publicznego i~stanowi
potwierdzenie tezy, że samo umieszczenie serwerów na terytorium UE nie
przesądza o~jurysdykcji prawnej. Z~perspektywy zgodności z~DORA oznacza to,
że instytucje finansowe nie dysponują dotychczas uzgodnionym,
zharmonizowanym poziomem odniesienia dla oceny suwerenności dostawcy ICT
na poziomie unijnym.

\subsection*{Cloud and AI Development Act (CADA)}

W~dniu 3~czerwca 2026~r.\ Komisja Europejska złożyła wniosek legislacyjny
dotyczący \textit{Cloud and AI Development Act} (CADA), COM(2026)~502
\autocite{cada2026}. Inicjatywa ta, zapowiadana wcześniej w~programie prac
Komisji na rok 2026~\autocite{ec_workprogramme2026} i~analizowana przez
Europejskie Biuro Badań Parlamentarnych (EPRS) już w~grudniu 2025~roku
\autocite{eprs_cada2025}, stanowi najdalej idące narzędzie prawne w~arsenale
unijnej suwerenności cyfrowej.

CADA wyrasta bezpośrednio z~rekomendacji raportu Draghiego z~2024~r.,
który wskazał strukturalną zależność UE od nieeuropejskich dostawców chmury
jako ryzyko strategiczne dla konkurencyjności europejskiej
gospodarki~\autocite{draghi2024}. Wniosek legislacyjny opiera się na
trzech filarach:

\begin{enumerate}[label=(\roman*)]
  \item rozbudowie europejskiej infrastruktury obliczeniowej, obejmującej
    tzw.\ \textit{AI Gigafactories}, czyli dużych centrów superobliczeniowych
    dedykowanych trenowaniu modeli sztucznej inteligencji,
  \item ustanowieniu europejskiego systemu oceny i~certyfikacji dostawców
    chmury i~AI, powiązanego z~preferencjami w~zamówieniach publicznych
    dla dostawców spełniających kryteria zgodności,
  \item wzmocnieniu europejskiej niezależności infrastrukturalnej
    przez mechanizmy prawne ograniczające ryzyko pozajurysdykcyjnego dostępu
    do danych.
\end{enumerate}

Podstawą prawną wniosku jest art.~114 Traktatu o~Funkcjonowaniu Unii
Europejskiej (TFUE), a~projekt przewiduje możliwość nałożenia kar finansowych
za niezgodność z~wymaganiami~\autocite{cada2026}.

\paragraph{Definicja suwerenności chmury jako oś konfliktu.}
Kluczowym nierozstrzygniętym zagadnieniem CADA jest definicja
\textit{suwerennej chmury}. Projekt przekształca cyfrową suwerenność
z~politycznej deklaracji w~wiążący wymóg prawny, jednak nie przesądza,
jak ścisłe mają być kryteria kwalifikacji. Zarysowały się dwa
stanowiska. Po pierwsze, {definicja wąska} (popierana m.in.\
przez Francję i~grupę lobby CISPE): suwerenność wymaga pełnej europejskiej
kontroli kapitałowej nad dostawcą, co w~praktyce wyklucza z~zamówień
publicznych podmioty kontrolowane przez korporacje spoza EOG.
Po drugie, {definicja szeroka}: suwerenność oceniana jest
na podstawie faktycznych gwarancji prawno-technicznych (rezydencji danych,
szyfrowania i~braku dostępu operacyjnego podmiotu macierzystego)
niezależnie od struktury właścicielskiej.

Spór ten ma bezpośrednie implikacje rynkowe. Przyjęcie wąskiej definicji
oznaczałoby de facto ograniczenie dostępu do rynku zamówień publicznych
dla Amazon Web Services, Microsoft Azure i~Google Cloud, które łącznie
kontrolują około 70\% europejskiego rynku chmury~\autocite{synergy2025eu}.
CADA działałby wówczas jako nieformalna bariera regulacyjna wobec podmiotów
objętych jurysdykcją CLOUD Act i~FISA~702. Szerokie kryterium pozwoliłoby
natomiast hiperskalersom na zachowanie dostępu do rynku publicznego
przez joint ventures z~europejskimi partnerami lub wdrożenie rozwiązań
\textit{confidential computing} jako środka kompensacyjnego.

\paragraph{Status legislacyjny i~perspektywa czasowa.}
Na dzień złożenia niniejszej pracy CADA znajduje się na wczesnym etapie
procesu legislacyjnego UE, złożony wniosek Komisji musi przejść
przez negocjacje w~Parlamencie Europejskim i~Radzie UE (trilog).
Kluczowe parametry aktu (w~tym wiążąca definicja suwerenności
oraz zakres sankcji) pozostają otwarte na etapie prac
legislacyjnych~\autocite{cada2026}. Przy typowym dla unijnych regulacji
cyfrowych tempie prac, realistyczny termin wejścia w~życie CADA to
nie wcześniej niż 2028--2029~r., z~uwzględnieniem standardowego
okresu \textit{vacatio legis}.

\subsection*{Implikacje dla modelu ORION}

Trzy omówione inicjatywy, tj. Data Act, EUCS i~CADA, tworzą wzajemnie
uzupełniający się system regulacyjny. Data Act reguluje prawa dostępu
i~przenoszenia danych \textit{ex post}, EUCS certyfikuje bezpieczeństwo
dostawców, a~CADA (jako złożony wniosek legislacyjny) ma budować
alternatywną, europejską infrastrukturę \textit{ex ante}. Z~perspektywy
metodologii ORION istotne jest jednak, że żadna z~tych regulacji nie
eliminuje jeszcze w~całości konfliktu jurysdykcyjnego wynikającego
z~CLOUD Act: Data Act nie uchyla prawnych obowiązków dostawcy
wynikających z~prawa USA, EUCS pozostaje nieukończony na poziomie
wymaganym dla usług krytycznych, a~CADA jest na wczesnym etapie
procesu legislacyjnego.

Stanowi to empiryczne uzasadnienie dla jednej z~kluczowych decyzji
projektowych ORION: traktowania jurysdykcji geograficznej dostawcy jako
czynnika ryzyka, którego nie można się pozbyć w~drodze samej konfiguracji
usługi, wymagającego wymiany dostawcy lub zastosowania poufnego przetwarzania
jako warstwy kompensacyjnej.



% -----------------------------------------------
\section{Przykładowe scenariusze konfliktu}

Konflikt wynikający z Cloud Act pojawia się na wielu płaszczyznach. Ogólny ogląd na przechowywanie danych w chmurze należy uszczegółowić, by wskazać, jakie informacje i~w jakich sytuacjach mogą opuszczać EOG. Po pierwsze należy zwrócić uwagę, w jaki sposób dane są przechowywane.

Przekazywane mogą być całe kopie maszyn wirtualnych, bez względu na jakim systemie operacyjnym działają. W takiej sytuacji nie ma znaczenia czy podmiot korzysta z rozwiązań Open Source (Linux vs Windows), nie ma również znaczenia jakie oprogramowanie jest na danym systemie zainstalowane oraz kto jest jego producentem, np. komercyjny silnik bazy danych Oracle, Microsoft SQL, etc. vs PostgreSQL, MySQL. Ryzyko związane z CLOUD Act zależy nie od silnika bazy danych (MySQL, Oracle, PostgreSQL), lecz od lokalizacji infrastruktury i~jej operatora. Sama zmiana silnika bazy danych nic nie zmienia, decydująca jest jurysdykcja nad infrastrukturą hostingową.

W przypadku systemów operacyjnych użytkownik narażony jest na dodatkowy kanał ryzyka, równoległy do przesłania danych z maszyn wirtualnych, czyli informacji zebranych przez system, przykładowo mechanizm telemetrii dostępny nie tylko w systemie Windows\autocite{microsoft_telemetry_docs}, ale również macOS\autocite{apple_telemetry_privacy_guides} oraz zostanie dołączona do dystrybucji Linux Fedora\autocite{fedora_wiki_telemetry}, rozwijanej przez RedHat / IBM.

Telemetria Windows nie przesyła wprost zawartości baz danych, jednak może przekazywać informacje takie jak nazwy procesów i~ścieżki plików, zdarzenia audytu logowania (np. SQL Server Audit), czy błędy połączeń
bazodanowych. W kontekście postępowania informacje te mogą być wystarczające do ukierunkowania żądania CLOUD Act
wobec Google (jak i~każdego innego hiperscalera podlegającego pod CLOUD Act) o wydanie pełnych danych z dysku VM.

\subsection{Mechanizm telemetrii w systemie Windows}

Systemy operacyjne z rodziny Windows zbierają dane diagnostyczne
i behawioralne tzw. \textit{telemetrię},  obejmujące informacje o:
\begin{itemize}[noitemsep]
  \item używanych aplikacjach i~ich awariach,
  \item aktywności sieciowej i~połączeniach,
  \item konfiguracji sprzętu,
  \item zdarzeniach systemowych rejestrowanych w Event Viewer,
  \item lokalizacji użytkownika (przy odpowiednim poziomie uprawnień).
\end{itemize}

Kluczową usługą odpowiedzialną za transmisję telemetrii jest
\texttt{DiagTrack} (\textit{Connected User Experiences and Telemetry}),
która regularnie przesyła dane do serwerów Microsoft zlokalizowanych w USA.
Dane te, trafiając do infrastruktury podmiotu objętego jurysdykcją
amerykańską, wchodzą automatycznie w zakres CLOUD Act.


% \subsubsection{Rekomendacje dotyczące telemetrii}

% W środowiskach podlegających NIS2, DORA lub RODO zaleca się:
% \begin{enumerate}[noitemsep]
%   \item Wyłączenie lub ograniczenie telemetrii do poziomu
%   \texttt{Security} lub \texttt{Basic} za pomocą GPO lub edytora rejestru.
%   \item Zablokowanie ruchu egress do endpointów telemetrycznych
%   Microsoft na poziomie firewalla (\texttt{vortex.data.microsoft.com},
%   \texttt{settings-win.data.microsoft.com} i~inne).
%   \item Unikanie kont Microsoft powiązanych z systemem operacyjnym.
%   \item Rozważenie migracji do systemów operacyjnych typu open-source
%   (Linux) w środowiskach przetwarzających dane wrażliwe.
% \end{enumerate}

% -----------------------------------------------


\subsection{Scenariusz kumulacji ryzyk: Windows Server na Google Cloud, Microsoft Azure, AWS}

Warto przeprowadzić eksperyment myślowy ze scenariuszem, w który  uruchamiono system operacyjny na maszynie wirtualnej w centrum danych. Szczególnie wysokim profilem ryzyka charakteryzuje się konfiguracja polegająca na uruchomieniu Windows Server jako maszyny wirtualnej w Google Cloud Platform (lub innej podległej pod jurysdykcję CLOUD Act) z zainstalowanym silnikiem bazy danych.
Wskazanie Google Cloud Platform oraz Microsoft Windows są wizualizacją przykładu i~wynikają z popularności usług.
\vspace{0.3cm}
\noindent Scenariusz ten łączy jednocześnie trzy warstwy ryzyka.
 
\begin{enumerate}[label=(\roman*)]
  \item \textit{Warstwa 1, Chmura a CLOUD Act.}
  Hostingującą firmą jest podmiot amerykański, z tego wynika, że każda maszyna wirtualna w tym centrum danych,
  niezależnie od regionu (w tym \texttt{europe-central2} w Warszawie), 
  podlega CLOUD Act. Organy USA mogą żądać danych od dostawcy bezpośrednio,
  bez informowania właściciela.

  \item \textit{Warstwa 2, Windows Server a telemetria.}
  Windows Server uruchomiony jako VM na chmurze generuje telemetrię przesyłaną
  do serwerów Microsoft. Oznacza to {dwa niezależne podmioty
  objęte CLOUD Act}: Google (infrastruktura VM) oraz Microsoft (telemetria
  systemowa). Możliwe jest równoległe żądanie skierowane do obu firm.

  \item \textit{Warstwa 3, Dostęp Google do danych bazy danych.}
  Google jako hypervisor ma techniczną możliwość dostępu do dysków
  Persistent Disk przypisanych do maszyny wirtualnej. Przy domyślnym
  szyfrowaniu zarządzanym przez Google KMS, firma może odczytać zawartość
  dysku w tym pliki danych bazy snapshoty, logi i~backupy.
\end{enumerate}

\subsection{Macierz wektorów dostępu}

\begin{table}[h!]
\centering
\caption{Wektory dostępu do danych bazy przy konfiguracji Windows Server / GCP}
\label{tab:vectors}
\renewcommand{\arraystretch}{1.4}
\begin{tabular}{>{\raggedright}p{4.5cm} >{\raggedright}p{2.5cm}
>{\raggedright\arraybackslash}p{5.5cm}}
\toprule
\textbf{Wektor dostępu} & \textbf{Podmiot} & \textbf{Co może być ujawnione} \\
\midrule
Nakaz CLOUD Act do Google & Google LLC & Pliki danych bazy, snapshoty
dysków, logi, backupy GCS \\
Telemetria Windows & Microsoft Corp & Metadane systemu, Event Log,
ścieżki plików \\
Cloud SQL (jeśli użyte) & Google LLC & Bezpośredni dostęp do tabel
przez DBaaS API \\
Klucze Google KMS (domyślne) & Google LLC & Możliwość odszyfrowania
dysków VM \\
\bottomrule
\end{tabular}
\end{table}

\subsection{Rekomendacje dla środowisk regulowanych (NIS2 / DORA / RODO)}

Na podstawie przeprowadzonej analizy formułuje się następujące rekomendacje
dla organizacji przetwarzających dane wrażliwe lub dane operacyjne
podlegające dyrektywie NIS2 i~rozporządzeniu DORA:

\begin{enumerate}[label=(\roman*)]
  \item \textit{CMEK z zewnętrznym HSM} (\textit{Customer-Managed
  Encryption Keys}), klucze do dysków VM przechowywane poza
  infrastrukturą dostawcy chmury eliminują techniczną możliwość
  odszyfrowania przez Google lub AWS.

  \item \textit{BYOK dla baz danych} (\textit{Bring Your Own Key}), niezależne szyfrowanie danych
  w bazie (np. TDE \textit{Transparent Data Encryption}) z kluczem
  zarządzanym przez klienta, poza środowiskiem chmurowym.

  \item \textit{Wyłączenie telemetrii Windows} na poziomie GPO
  (\texttt{Security level}) oraz zablokowanie ruchu do endpointów
  telemetrycznych na poziomie firewalla egress.

  \item \textit{Wybór europejskiego dostawcy infrastruktury w usłudze} spoza jurysdykcji
  USA (w pełni zlokalizowane w EOG) jako infrastruktury
  dla systemów krytycznych.

  \item \textit{Uwzględnienie CLOUD Act w rejestrze ryzyk IT}, jako
  ryzyko regulacyjne kategorii \textit{prawne / compliance},
  z przypisaniem prawdopodobieństwa i~wpływu zgodnie z metodologią
  przyjętą w organizacji.

  \item \textit{Audyt umów z dostawcami chmurowymi} pod kątem klauzul
  dotyczących żądań rządowych oraz regularny monitoring
  \textit{Government Request Transparency Reports} publikowanych przez
  Google, Microsoft i~Amazon.
\end{enumerate}

% -----------------------------------------------
\subsection{Wnioski}

CLOUD Act stanowi realne i~wielowarstwowe zagrożenie dla suwerenności
danych europejskich organizacji. Zagrożenie to nie ogranicza się do
danych przechowywanych bezpośrednio w usługach chmurowych (DBaaS, SaaS),
lecz obejmuje również metadane generowane przez systemy operacyjne
(telemetria Windows) oraz dane na dyskach wirtualnych maszyn IaaS
zarządzanych przez podmioty podlegające prawu USA.

Konfiguracja łącząca Windows Server z Google Cloud Platform jako
infrastrukturą hostingową bazy danych tworzy najwyższy profil
ryzyka, dwa niezależne podmioty objęte CLOUD Act uzyskują potencjalny dostęp do danych przez różne kanały techniczne.

Skuteczna mitigacja tego ryzyka wymaga działań na poziomie infrastrukturalnym (wybór dostawcy spoza jurysdykcji USA), kryptograficznym
(CMEK, BYOK) oraz operacyjnym (wyłączenie telemetrii, audyt umów).
Środki te powinny być traktowane jako elementy systemu zarządzania ryzykiem IT, spójne z wymogami dyrektywy NIS2 i~rozporządzenia DORA.


\chapter{Ryzyka w praktyce: zagrożenia, konsekwencje i dobre praktyki}

Zrozumienie architektury środowiska organizacji oraz pełna inwentaryzacja usług jest kluczem do prawidłowego zarządzania ryzykiem ICT. Badania przeprowadzone przez firmę \textit{Swimlane} wykazały, że~71\% organizacji nie jest w~stanie całościowo dokonać audytu wszystkich systemów zgodnie z~wewnętrznymi i~zewnętrznymi standardami \autocite{swimlane2025grc}.\\
Badania wskazują jako przyczyny:

\begin{enumerate}[label=(\roman*)]
    \item \textit{Fragmentację narzędzi}: 92\% respondentów korzysta
          z~co najmniej trzech różnych narzędzi do zbierania dowodów
          audytowych, co powoduje powielanie pracy i~niespójność danych;
    \item \textit{Uzależnienie od procesów manualnych}: średnio jedynie
          39\% procesu zbierania dowodów jest zautomatyzowane; ponad
          połowa organizacji (54\%) poświęca ponad 5~godzin tygodniowo
          na ręczne zadania compliance;
    \item \textit{Podatność na błędy w~danych}: 62\% organizacji
          przyznaje, że proces zbierania dowodów audytowych jest
          przynajmniej sporadycznie obarczony błędem;
    \item \textit{Brak współpracy między zespołami}:
          90\% organizacji wskazuje, że słaba koordynacja między 
          zespołami  podważa gotowość do audytu;
    \item \textit{Rosnącą złożoność regulacyjną}: 96\% organizacji
          twierdzi, że dotrzymanie kroku rosnącej liczbie regulacji
          branżowych stanowi istotne wyzwanie.
\end{enumerate}


Jakość prowadzonych audytów bywa podważana przez samych badanych. Brak kompetencji osób wypełniających kwestionariusze audytowe nie gwarantuje poprawności udzielanych odpowiedzi.
Badanie \textit{The Next Generation Cybersecurity Auditor}
przeprowadzone na próbie 151~tys. audytorów IT z firm Big Four (\textit{Deloitte, Ernst \& Young, KPMG, PricewaterhouseCoopers}) wykazało, że poziom wiedzy proceduralnej (praktycznej) audytorów był o~17\% niższy niż poziom ich wiedzy teoretycznej, a poziom pewności siebie respondentów nie był skorelowany z ich rzeczywistymi
wynikami, co wskazuje na efekt Dunninga-Krugera w~środowisku audytorskim \autocite{nextgen2022auditor}.


Dodatkowym problemem obniżającym jakość audytów jest zjawisko nierzetelnego wypełniania kwestionariuszy przez respondentów. Badania z zakresu ekonomii behawioralnej wskazują, że nieuczciwość w~systemach samoraportowania jest powszechna, Friesen i~Gangadharan wykazali eksperymentalnie, że~respondenci częściej podają prawdziwe
informacje w~reżimach raportowania obowiązkowego niż dobrowolnego, co sugeruje, iż~bez odpowiednio zaprojektowanych mechanizmów egzekwowania trudno osiągnąć satysfakcjonujący poziom zgodności
odpowiedzi ze~stanem faktycznym \autocite{friesen2013designing}. Potwierdzają to badania Peer i~in., którzy wykazali, że~zobowiązanie do złożenia przysięgi uczciwości zmniejsza skłonność do nieuczciwości (wyrażona w miarze statystycznej ,,iloraz szans'' (ang.  \textit{Odds Ration, OR})): \textit{OR} = 0{,}79 [0{,}72; 0{,}88], jednak sama relacja interpersonalna między wypełniającym a~audytorem nie wywiera istotnego wpływu na~rzetelność odpowiedzi
(\textit{OR} = 1{,}08 [0{,}97; 1{,}20]) \autocite{peer2023commitment}.
Oznacza to, że~dowody pozyskane poprzez nieusystematyzowane, dobrowolne kwestionariusze audytowe mogą być obarczone istotną słabością metodologiczną, a~ich wartość jako podstawy oceny ryzyka ICT powinna być traktowana z~należytą ostrożnością.

\section{Shadow IT}

Trudnym do wykrycia problemów z bezpieczeństwem jest zjawisko zwane jako \textit{Shadow IT}. 
Definiuje się je, jako wykorzystanie przez pracowników narzędzi informatycznych, 
aplikacji, urządzeń lub usług bez wiedzy i zgody działu IT organizacji  \autocite{ibm_shadowit}. W wyniku restrykcyjnych polityk  bezpieczeństwa, użytkownicy sięgają po nieautoryzowane rozwiązania informatyczne w celu usprawnienia codziennej pracy, zjawisko to nasilił wzrost popularności  modelu pracy zdalnej \autocite{jumpcloud2024}.

Raport Capterra z 2023 roku wskazuje, że 76\% małych i średnich przedsiębiorstw identyfikuje \textit{Shadow IT} jako umiarkowane lub poważne zagrożenie dla cyberbezpieczeństwa \autocite{capterra2023}. Badania JumpCloud pokazują, że 
w 2023 roku 41\% pracowników korporacyjnych korzystało z technologii pozostających 
poza nadzorem działu IT \autocite{jumpcloud2024}, a średnio firma dysponuje 975 
nieznanymi usługami chmurowymi, podczas gdy jedynie 108 jest przez nią oficjalnie 
monitorowanych \autocite{josys2024}.

Niepokojący jest wpływ \textit{Shadow IT} na bezpieczeństwo organizacji. Badanie 
Kaspersky z 2023 roku wykazało, że 85\% firm doświadczyło incydentów cyberbezpieczeństwa, 
przy czym w co najmniej 11\% przypadków rola \textit{Shadow IT} była bezpośrednio 
udokumentowana \autocite{kaspersky2023}. Szacuje się, że prawie połowa wszystkich 
cyberataków ma swoje źródło właśnie \textit{w Shadow IT}. Wykrycie \textit{Shadow IT} jest szczególnie trudne, ponieważ zasoby te z definicji funkcjonują 
poza zasięgiem standardowych narzędzi monitorowania sieci i systemów bezpieczeństwa \autocite{isaca2024}.

Istnienie \textit{Shadow IT} w organizacji stanowi sygnał, że istniejące polityki i narzędzia IT nie zaspokajają w pełni potrzeb 
pracowników. Pracownicy wykorzystujący w swojej pracy  nieautoryzowane rozwiązania pokazują, że istnieją luki w oficjalnej ofercie narzędziowej organizacji.  Sutuacja ta powinna  skłoniać  działy IT do werfyikacji swoich polityk oraz systemów szkoleń\autocite{grip2025}.


\section{Shadow AI}

Dynamiczny rozwój narzędzi generatywnej sztucznej inteligencji w~latach 2022-2025 wytworzył nową kategorię zjawiska Shadow IT, określaną jako \textit{Shadow AI}. Definiuje się je jako wykorzystanie przez pracowników systemów sztucznej inteligencji, w~szczególności narzędzi generatywnych, bez wiedzy, zgody i~oceny bezpieczeństwa ze~strony działu IT organizacji \autocite{gartner2025crossborder}.

Shadow AI wyróżnia się od klasycznego Shadow IT jedną właściwością istotnie zwiększającą profil ryzyka: użycie narzędzia AI nieodłącznie wiąże się z~przesyłaniem danych wejściowych do infrastruktury zewnętrznego dostawcy. Klasyczne Shadow IT może działać lokalnie, bez transferu danych poza organizację. Shadow AI natomiast z~definicji eksportuje dane wejściowe każdorazowo, gdy pracownik sformułuje zapytanie do modelu.

Skala zjawiska jest znaczna. Raport LayerX Security wykazał, że 50\% pracowników przyznało się do wklejania wrażliwych danych firmowych do narzędzi generatywnej AI, z~czego 18\% deklarowało udostępnianie danych szczególnie wrażliwych \autocite{layerx_2025}. Analiza Cyberhaven Labs, obejmująca 1,6~mln pracowników, wykazała, że 11\% danych wklejanych do ChatGPT stanowią informacje poufne, a~przeciętna firma notuje kilkaset takich zdarzeń tygodniowo \autocite{cyberhaven_chatgpt}. Raport Microsoft \textit{Work Trend Index 2024} wskazał, że 78\% pracowników korzystających z~AI w~pracy robi to za pomocą własnych, prywatnych kont w~narzędziach niezatwierdzonych przez pracodawcę \autocite{microsoft2024worktrend}.

Shadow AI generuje trzy odrębne kategorie ryzyka regulacyjnego.

\begin{enumerate}[label=(\roman*)]
      \item \textit{Ryzyko naruszenia RODO.}
Dane przesyłane do modeli generatywnej AI trafiają na serwery dostawców zlokalizowane poza Europejskim Obszarem Gospodarczym. Stanowi to transfer danych osobowych wymagający ważnej podstawy prawnej w~rozumieniu rozdziału~V RODO, takiej jak standardowe klauzule umowne (SCC) lub przynależność dostawcy do Ram Ochrony Prywatności Danych (DPF). W~przypadku Shadow AI transfer ten odbywa się bez wiedzy administratora, co wyklucza spełnienie wymogu rozliczalności z~art.~5 ust.~2 RODO: pracownik, bez stosownego umocowania, dokonuje transferu danych, za~który organizacja ponosi pełną odpowiedzialność.

\item \textit{Ryzyko na gruncie AI Act.}
Rozporządzenie (UE) 2024/1689, zwane \textit{AI Act}, które weszło w~życie 1~sierpnia 2024~roku, nakłada na podmioty \textit{stosujące} systemy AI obowiązki uzależnione od poziomu ryzyka danego systemu \autocite{euaiact2024}.
Korzystanie przez pracownika z~niezatwierdzonego modelu AI do wspomagania decyzji kredytowych, rekrutacyjnych lub diagnostyki medycznej może oznaczać, że organizacja faktycznie stosuje system AI wysokiego ryzyka w~rozumieniu art.~6 i~załącznika~III AI~Act bez przeprowadzenia wymaganej oceny zgodności i~bez zapewnienia nadzoru ludzkiego. Odpowiedzialność spoczywa na organizacji niezależnie od tego, że wdrożenie nastąpiło bez wiedzy zarządu.

\item \textit{Ryzyko jakości i~wiarygodności.}
Modele generatywnej AI mogą generować błędne informacje, określane jako \textit{halucynacje modelu}. Wbudowanie wyników działania niezweryfikowanego modelu w~procesy decyzyjne, na~przykład w~analizy prawne, rekomendacje medyczne lub wyceny finansowe, bez procedury weryfikacji tworzy ryzyko operacyjne, za~które odpowiada organizacja. Warunki korzystania z~większości narzędzi generatywnej AI wprost wyłączają odpowiedzialność dostawcy za~skutki biznesowe zastosowania generowanych treści.
\end{enumerate}

Podobnie jak Shadow IT, Shadow AI jest diagnozą organizacyjną: pracownicy sięgają po niezatwierdzone modele AI, ponieważ zatwierdzone rozwiązania nie spełniają ich potrzeb lub procedura akceptacji jest zbyt długa. Odpowiedzią nie jest wyłącznie zakaz, lecz wdrożenie zatwierdzonego środowiska AI uzupełnionego polityką dopuszczalnego użycia, która definiuje, jakich danych nie wolno wprowadzać do żadnego zewnętrznego modelu.


\section{Edukacja jako fundament ochrony danych osobowych}

Problem \textit{shadow IT} oraz niewłaściwego posługiwania się danymi osobowymi jest w~znacznej mierze wynikiem niewystarczającego poziomu świadomości i~wiedzy pracowników.
Badania przeprowadzone przez Stanford University i~firmę Tessian wskazują, że 88\%~naruszeń bezpieczeństwa informacji wynika z~błędu ludzkiego, a~nie
z~działania cyberprzestępców~\autocite{isoqar_blad_ludzki}. Raport Verizon \textit{Data Breach Investigations Report~2024} wskazuje,
że czynnik ludzki, zdefinowany jako  nieświadome błędy lub brak kompetencji, pozostaje głównym źródłem naruszeń bezpieczeństwa danych ~\autocite{verizon_dbir_2024}.
Raport \textit{Cyberportret polskiego biznesu} z~2024~roku przedstawił dane dla Polski: ponad połowa pracowników w~polskich firmach nie przeszła żadnego
szkolenia z~zakresu cyberbezpieczeństwa i~ochrony danych w~ciągu ostatnich pięciu lat~\autocite{brandsit_szkolenia_2025}.

Skala i~charakter najczęstszych incydentów potwierdzają diagnozę. W~2024~roku zespół CERT~Polska odnotował wzrost liczby cyberzagrożeń, liczba zgłoszeń w jednym roku przekroczyła 600~tysięcy, najczęstszym incydentem były ataki phishinggowe, czyli uzyskanie korzyści przez przestępców poprzez manipulacje zachowania ofiery, a~nie na przełamaniu
zabezpieczeń technicznych~\autocite{cert_polska_2024}. W~badaniu przeprowadzonym przez Keepnet Labs w~2026~roku 45\%~pracowników przyznało, że ich
organizacja nie zapewniła szkoleń z~zakresu bezpieczeństwa\autocite{keepnet_2026}.

W~kontekście polskiego rynku sygnałem pozytywnym jest fakt, że cyberbezpieczeństwo stało się najszybciej rosnącym segmentem rynku IT w~Polsce, a w~latach 2024-2025 
niemal dwie trzecie respondentów wskazało je na pierwszym miejscu wśród inwestycji IT, 43\%~badanych uznało je za najważniejszy pojedynczy obszar~\autocite{pap_cyberbezpieczenstwo_2026}. Wzrost wydatków koncentruje się na rozwiązaniach technicznych, a~nie na edukacji pracowników, co~przy dominacji błędu ludzkiego jako źródła incydentów stanowi strategiczną asymetrię w~alokacji zasobów~\autocite{brandsit_szkolenia_2025}.

\subsection*{Wymiar regulacyjny obowiązku szkolenia}

Edukacja pracowników nie jest jedynie kwestią dobrej praktyki lub woli pracodawcy. Aktualnie jest obowiązkiem, wynikającym aktów prawnych.
RODO nakłada na administratora obowiązek zapewnienia, by osoby upoważnione do przetwarzania danych osobowych były odpowiednio przeszkolone (art.~32 ust.~4 RODO),
a~brak szkoleń jako środka techniczno-organizacyjnego stanowi  podstawę do nałożenia kary przez UODO~\autocite{isoqar_blad_ludzki}. Rozporządzenie DORA
zobowiązuje instytucje finansowe do prowadzenia programów podnoszenia świadomości w~zakresie ryzyka ICT, w~tym kadry zarządzającej~\autocite{nflo_dora}.
Norma ISO/IEC~27001 wprost wymaga wdrożenia udokumentowanych programów szkoleń i~uświadamiania jako kontroli organizacyjnej~\autocite{isoqar_blad_ludzki}.
Brak lub niewystarczający program szkoleniowy staje się zatem nie tylko źródłem
ryzyka operacyjnego, lecz również jest przyczyną i źródłem niezgodność z~obowiązującym prawem.

\section{Często występujące ryzyka związane z~nieprawidłowym przetwarzaniem danych}

Niewłaściwe i~ryzykowne zachowania pracowników można zidentyfikować i~wskazać główne
ich źródła. Poniżej przedstawiono najistotniejsze kategorie zachowań, udokumentowane
przez raporty branżowe, organy nadzoru oraz literaturę naukową.

\begin{enumerate}[label=(\roman*)]

\item \textit{Przekazywanie danych do narzędzi generatywnej AI.}
Pracownicy przekazują dane wrażliwe, treści umów, numery PESEL,
historię transakcji,  do narzędzi takich jak ChatGPT, Copilot czy Gemini,
w~celu ich podsumowania lub analizy. Dane te trafiają bezpośrednio na serwery
dostawców zlokalizowane w~Stanach Zjednoczonych, poza kontrolą organizacji.
Raport LayerX Security \textit{Enterprise AI and SaaS Data Security Report~2025}
wykazał, że \textit{50\% pracowników} przyznało się do wklejania wrażliwych
danych firmowych do narzędzi generatywnej AI, przy czym 18\% z~nich
deklarowało udostępnianie danych szczególnie wrażliwych, w~tym danych produkcyjnych
i~zastrzeżonych~\autocite{layerx_2025}. Równolegle analiza Cyberhaven Labs,
obejmująca 1,6~mln pracowników, wykazała, że 11\% danych
wklejanych przez pracowników do ChatGPT stanowią informacje poufne, a~przeciętna
firma notuje kilkaset takich zdarzeń tygodniowo\autocite{cyberhaven_chatgpt}.

\item \textit{Używanie prywatnych kont w~usługach chmurowych do transferu dokumentów służbowych.}
Pracownicy przesyłają dokumenty zawierające dane osobowe na prywatne konta Google Drive,
Dropbox czy Microsoft OneDrive, omijając firmowe procedury DLP
(\textit{Data Loss Prevention}). Według raportu WifiTalents \textit{Shadow IT Statistics~2025},
58\% pracowników korzysta z~prywatnych rozwiązań do przechowywania plików
bez zgody działu IT, a 47\% omija kontrole bezpieczeństwa
przy użyciu narzędzi shadow IT~\autocite{wifitalents_shadow_2025}. Dane te,
gdy zostaną umieszczone poza infrastrukturą organizacji, stają się niedostępne w celach nadzorczych, co uniemożliwia wykazanie zgodności
z~zasadą rozliczalności wynikającą z~art.~5 ust.~2 RODO.

\item \textit{Przesyłanie danych przez niezatwierdzone komunikatory i~dyski chmurowe.}
Pracownicy korzystają z~komunikatorów takich jak WhatsApp lub~Messenger.
Choć komunikacja w~aplikacjach takich jak WhatsApp jest szyfrowana \textit{end-to-end}, czyli wyłącznie między urządzeniami uczestników, nie gwarantuje to zgodności z~RODO, ponieważ metadane, kopie zapasowe i~historia czatów mogą być przechowywane przez dostawcę poza EOG~\autocite{dataprotectionpeople_human_error}. Badanie WifiTalents
wykazało, że 37\% pracowników używa niezatwierdzonych komunikatorów
do komunikacji służbowej~\autocite{wifitalents_shadow_2025}.

\item \textit{Błędne utożsamianie pseudonimizacji z~anonimizacją.}
W~dokumentach przekazywanych do analizy lub testów pozostają inicjały, numery PESEL
lub numery rachunków bankowych, które pracownicy nierzadko traktują jako
wystarczające zanonimizowanie danych. Dane pseudonimowe nadal
identyfikują osobę fizyczną w~rozumieniu art.~4 pkt~1 RODO i~nie wyłączają
stosowania przepisów rozporządzenia~\autocite{aepd_test_environments}.
Preuveneers i~in.\ wykazali, że modele głębokiego uczenia
są w~stanie z~dokładnością 73,4\% zidentyfikować gospodarstwa domowe
na podstawie pozornie pseudonimowych danych pomiarowych w~zbiorze liczącym
ponad 5000 podmiotów~\autocite{preuveneers2024repseudo}.

\item \textit{Używanie narzędzi do transkrypcji spotkań bez oceny ryzyka.}
Narzędzia takie jak Otter.ai, Fireflies.ai czy rozszerzenia oparte na ChatGPT
automatycznie dołączają do posiedzeń, nagrywają rozmowy i~przesyłają transkrypty
na serwery zewnętrznych dostawców, często zlokalizowane w~USA~\autocite{legal_io_transcription}.
Uczestnicy spotkań mogą nie być świadomi obecności utomatu nagrywającego,
co~stanowi naruszenie zasady przejrzystości z~art.~5 ust.~1 lit.~a RODO.
Co istotne, warunki korzystania wielu  narzędzi transkrypcyjnych, zezwalają na
wykorzystanie nagrań i transkrypcji do trenowania modeli AI, co oznacza, że treści posiedzeń
z~danymi klientów mogą zasilać bazy treningowe dostawcy bez podstawy prawnej~\autocite{linkedin_transcription_risk}.

\item \textit{Przesyłanie danych wrażliwych mailem}.
Przekazywane danych wrażliwych pocztą elektroniczną do zewnętrznych adresatów
bez uprzedniego zawarcia umowy powierzenia przetwarzania danych (DPA) oraz odpowiedniego zabezpieczenia załącznika (bezpiecznycm hasłem), wymaganej przez art.~28 RODO.
Badanie MailIntelligence wykazało, że 65\% incydentów utraty danych
jest powodowanych przez wiadomości e-mail skierowane do niewłaściwego adresata
lub przesłane bez odpowiednich zabezpieczeń~\autocite{mailintelligence_2024}.
Takie przekazanie danych  stanowi naruszenie RODO, niezależnie od tego, czy doszło do faktycznego wycieku.

\item \textit{Kopiowanie danych produkcyjnych do środowisk testowych bez anonimizacji.}
Dane wrażliwe z~systemów produkcyjnych są przenoszone do środowisk deweloperskich
lub testowych w~niezmienionej formie, co narusza art.~25 RODO (zasada
\textit{privacy by design and by default}) oraz zasady minimalizacji i~ograniczenia celu~\autocite{aepd_test_environments}.
Agencja ochrony danych Hiszpanii (AEPD) wprost wskazuje, że używanie rzeczywistych
danych osobowych w~środowiskach testowych jest dopuszczalne jedynie gdy jest bezwzględnie konieczne i~zabezpieczone zgodnie z~art.~32 RODO~\autocite{aepd_test_environments}.
W~praktyce konieczność ta rzadko zachodzi, dostępne są narzędzia do generowania
danych syntetycznych, a~rezygnacja z~ich stosowania wynika z~wygody, nie z~merytorycznej potrzeby.

\item \textit{Przekazywanie danych pacjentów firmom IT bez podstawy prawnej.}
Podczas zgłaszania usterek w~oprogramowaniu medycznym pracownicy przekazują
firmom serwisowym fragmenty dokumentacji pacjentów lub całe zrzuty baz danych.
Najwyższa Izba Kontroli w~raporcie \textit{RODO w~szpitalu} z~2019~roku wykryła
ten proceder w~11 z~kontrolowanych szpital, gdzie udostępniono dane podmiotom zewnętrznym
bez umowy powierzenia ani innej podstawy prawnej~\autocite{nik_rodo_szpital}.

\item \textit{Przechowywanie dokumentów tożsamości klientów bez podstawy prawnej.}
Instytucje finansowe i~inne podmioty skanują i~przechowują kopie dokumentów tożsamości w~celach, mimo że przepisy prawa nie nakładają takiego obowiązku i~wystarczająca jest jedynie weryfikacja tożsamości bez utrwalania kopii dokumentu. Prezes UODO nałożył na \textit{ING Bank Śląski} karę w~wysokości 18,4~mln~zł
za masowe i~nieuzasadnione skanowanie oraz przechowywanie dokumentów tożsamości klientów,
co~było najwyższą karą w~historii polskiego organu nadzorczego w~roku 2025~\autocite{gazeta_prawna_kary_2025}.

\item \textit{Wysyłanie dokumentów kadrowych do niewłaściwych adresatów.}
Świadectwa pracy, paski płac, aneksy do umów i~inne dokumenty HR zawierające dane
osobowe kierowane są pocztą elektroniczną lub tradycyjną do błędnych odbiorców wskutek
pomyłek przy wpisywaniu adresu lub autouzupełniania skrzynki pocztowej.
Błędne skierowanie korespondencji jest jednym z najczęściej zgłaszanych
naruszeń do UODO, stanowi 18\% wszystkich raportowanych incydentów
bezpieczeństwa danych~\autocite{dataprotectionpeople_human_error, odo24_naruszenia}.
Każde takie zdarzenie wymaga oceny ryzyka i~potencjalnego zgłoszenia do organu nadzoru
w~ciągu 72~godzin od wykrycia.

\item \textit{Ignorowanie alertów DLP i~poszukiwanie obejść.}
Pracownicy traktują ostrzeżenia systemów DLP jako fałszywe alarmy i~aktywnie
szukają sposobów ich ominięcia, kopiując dane na prywatne pendrive'y,
wysyłając pliki przez prywatne konto e-mail lub korzystając z~platform
wymiany plików (WeTransfer, Smash). Badanie WifiTalents wykazało,
że 47\% pracowników omija kontrole bezpieczeństwa,
a~organizacje dysponują średnio 975 nieznanych usług chmurowych,
z~których monitorują jedynie ok.\ 108~\autocite{wifitalents_shadow_2025, fidelis_shadow_2026}.
Gartner szacuje, że na typowego pracownika korporacyjnego przypada średnio
13 aplikacji shadow IT używanych równolegle~\autocite{wifitalents_shadow_2025}.

\end{enumerate}








% --------------
% Slownik:

% KMS/HSM + BYOK/HYOK 
% Polska, jako państwo członkowskie Unii Europejskiej, od lat uczestniczy w budowaniu europejskiego systemu ochrony danych i cyberbezpieczeństwa. Punktem wyjścia dla krajowych regulacji w tym zakresie stało się wdrożenie unijnej dyrektywy NIS z 2016 roku, która zobowiązała państwa członkowskie do ustanowienia krajowych ram cyberbezpieczeństwa. W odpowiedzi na te wymagania, Polska uchwaliła ustawę o krajowym systemie cyberbezpieczeństwa, przyjętą przez Sejm RP dnia 5 lipca 2018 roku i podpisaną przez Prezydenta RP dnia 1 sierpnia 2018 r., która weszła w życie 28 sierpnia 2018 r.

% Dokładnie co ona nakłada ta ustawa?

% Isotną ustawą w krajowym prawodastwie jest Ustawa z dnia 10 maja 2018 r. o ochronie danych osobowych 

% Odpowiedzialność administracyjna RODO
% Administracyjne kary pieniężne (art. 83 RODO):
% • Do 10 mln EUR lub 2\% całkowitego rocznego obrotu za mniejsze przewinienia (np. brak rejestru 
% czynności przetwarzania, brak współpracy z organem).
% • Do 20 mln EUR lub 4\% obrotu za poważne naruszenia (np. złamanie podstawowych zasad 
% przetwarzania, naruszenie praw osób, nielegalny transfer danych do państw trzecich).
% Środki naprawcze (art. 58 RODO):
% • Ostrzeżenia i upomnienia.
% • Nakazanie spełnienia żądań osoby, której dane dotyczą.
% • Nakazanie dostosowania przetwarzania do przepisów RODO.
% • Czasowe lub całkowite ograniczenie przetwarzania, w tym zakaz przetwarzania danych.
% • Nakazanie usunięcia danych lub powiadomienia o naruszeniu

% Ustawa z dnia 16 kwietnia 1993 r. o zwalczaniu nieuczciwej konkurencji - napsiz jak ona sie ma do mojej monografii

% Skup sie też mocno nad tym:
% Ustawa z dnia 5 lipca 2018 r. o krajowym systemie cyberbezpieczeństwa / Ustawa o zmianie ustawy o krajowym systemie cyberbezpieczeństwa oraz niektórych innych ustaw uchwalona dn. 23.01.2026 (przekazana do podpisu prezydenta RP)


% Nie można pomijać odpowiedzialności z tytułu zapisów w Kodeksie Cywilnym, gdzie ... rozwin

% Na osoby odpowiedzialne jest rowniez kodeks karny - rozwin

% RODO / GDPR
% Ai Act
% NIS2

% \section{Fintech: Analiza DORA, PSD2/3 oraz wytyczne KNF}
% DORA
% NIS2 ->KSC
% RODO

% Akty krajowe

% Kodeks cywilny

% Kodeks karny

% Ustawa z dnia 5 lipca 2018 r. o krajowym systemie cyberbezpieczeństwa / Ustawa o zmianie ustawy o krajowym systemie cyberbezpieczeństwa oraz niektórych innych ustaw uchwalona dn. 23.01.2026 (przekazana do podpisu prezydenta RP)

% Ustawa z dnia 25 czerwca 2025 r. o zmianie niektórych  ustaw w związku z zapewnieniem operacyjnej odporności  cyfrowej sektora finansowego oraz emitowaniem europejskich zielonych obligacji (


% Ustawa z dnia 10 maja 2018 r. o ochronie danych osobowych 

% Ustawa z dnia 16 kwietnia 1993 r. o zwalczaniu nieuczciwej konkurencji

% Ustawa z dnia 6 listopada 2008 r. o prawach pacjenta i Rzeczniku Praw Pacjenta

% Ustawa z dnia 28 kwietnia 2011 r. o systemie informacji w ochronie zdrowia

% Aktu unijne

% Rozporządzenie RODO (ochrona danych osobowych)
% Rozporządzenie DORA (cyberodporność sektora  finansowego)
% Dyrektywa NIS2 (krajowy system cyberbezpieczeństwa)
% Dyrektywa CER (odporność podmiotów krytycznych)
% Data Act (bezpieczeństwo danych z IoT


% \section{Sektor medyczny}

% \section{Sektor finansowy}
% DORA

% \section{Zasada Privacy by DEsign z RODO}
% \section{Zasada TPP (Third Party Providers): Analiza łańcucha dostaw}
% \section{Suwerenność danych i strategia niezależności UE}
% \section{Wspólny mianownik: Ciągłość działania i TPRM}


\chapter{ORION - Ocena Ryzyka Informatycznego Organizacji Nadzorowanych}

\section{Luka badawcza}
\label{sec:luka}

Istniejące metodologie obejmują zarządzanie bezpieczeństwem, zarządzanie ryzykiem i~kontrole organizacyjne w znacznym stopniu w kierunku analizy systemów informatycznych. Są to standardy rynkowe, takie jak metodologie zdefiniowane w  ISO/IEC 27001, NIST Cybersecurity Framework, COBIT. Jednak żadna z nich nie spełnia jednocześnie  zasadniczych wymagań stawianych przez współczesne środowiska regulacyjne. Po pierwsze, brakuje im wielostronnej perspektywy obejmującej cały łańcuch dostaw ICT, od nadzorowanego podmiotu, poprzez jego bezpośrednich dostawców, po podprocesory działające na wielu poziomach. Po drugie, nie zawierają one opartego na dowodach mechanizmu punktacji, który rozróżniałby zweryfikowaną zgodność od zwykłego zapewnienia deklaratywnego rozróżnienia, które jest fundamentalne dla przeciwdziałania dobrze udokumentowanemu zjawisku \textit{teatru zgodności}. Po trzecie, nie zapewniają one wieloaspektowej struktury analitycznej, która formalnie przypisuje pytania audytowe do odpowiedniego podmiotu w organizacji, tj. radcy prawnego, biznesu, architekta oprogramowania lub administratora infrastruktury, w oparciu o charakter pytania, a nie wygodę organizacji. Po czwarte, nie oferują modułowej rozszerzalności sektorowej, która zachowuje uniwersalny rdzeń metodologiczny, jednocześnie aktywując perspektywy i~zestawy pytań specyficzne dla danej dziedziny dla regulowanych branż, takich jak finanse czy opieka zdrowotna.

\renewcommand{\arraystretch}{1.4}
% Use longtable instead of table + tabular
% Note: Total widths set to approx 16.5cm to fit A4
\begin{longtable}{p{4.2cm}p{2.2cm}p{2.2cm}p{2.2cm}p{2.4cm}p{1.8cm}}
\caption{Analiza porównawcza frameworków stosowanych w polskiej praktyce audytowej względem wymagań ORION} \label{tab:porownanie-frameworkow} \\
\toprule
\textbf{Wymaganie} 
& \textbf{ISO/IEC 27001} 
& \textbf{NIST CSF} 
& \textbf{COBIT 2019} 
& \textbf{Narzędzia GRC} 
& \textbf{ORION} \\
\midrule
\endfirsthead

% This header appears on subsequent pages
\multicolumn{6}{c}%
{{\bfseries \tablename\ \thetable{} -- kontynuacja z poprzedniej strony}} \\
\toprule
\textbf{Wymaganie} 
& \textbf{ISO/IEC 27001} 
& \textbf{NIST CSF} 
& \textbf{COBIT 2019} 
& \textbf{Narzędzia GRC} 
& \textbf{ORION} \\
\midrule
\endhead

\midrule
\multicolumn{6}{r}{{Kontynuacja na następnej stronie...}} \\
\endfoot

\bottomrule
\endlastfoot

\textbf{Pokrycie całego łańcucha dostaw ICT} 
& Częściowe\newline(tylko załącznik dot.\ dostawców) 
& Częściowe\newline(profil łańcucha dostaw) 
& Częściowe\newline(zarządzanie dostawcami) 
& Częściowe\newline(zależne od narzędzia) 
& \textbf{Pełne} \\
\midrule

\textbf{Ocena oparta na dowodach} 
& Nie 
& Nie 
& Nie 
& Częściowe\newline(zależne od konfiguracji) 
& \textbf{Tak} \\
\midrule

\textbf{Wieloperspektywiczne przypisanie pytań}\newline(prawna / biznesowa / techniczna) 
& Nie 
& Nie 
& Częściowe\newline(tylko macierz RACI) 
& Nie 
& \textbf{Tak} \\
\midrule

\textbf{Formalny model grafowy zależności ICT} 
& Nie 
& Nie 
& Nie 
& Nie 
& \textbf{Tak} \\
\midrule

\textbf{Propagacja ryzyka przez warstwy łańcucha dostaw} 
& Nie 
& Nie 
& Nie 
& Nie 
& \textbf{Tak} \\
\midrule

\textbf{Scoring zgodności z regulacjami}\newline(DORA / NIS2 / RODO) 
& Częściowe\newline(istnieją mapowania) 
& Częściowe\newline(istnieją mapowania) 
& Częściowe\newline(istnieją mapowania) 
& Częściowe\newline(zależne od konfiguracji) 
& \textbf{Tak} \\
\midrule

\textbf{Rozróżnienie wymogów obowiązkowych od dobrych praktyk} 
& Nie 
& Nie 
& Nie 
& Częściowe 
& \textbf{Tak} \\
\midrule

\textbf{Modułowa rozszerzalność sektorowa}\newline(finanse, zdrowie, energetyka) 
& Nie 
& Częściowe\newline(profile branżowe) 
& Nie 
& Częściowe\newline(płatne moduły) 
& \textbf{Tak} \\
\midrule

\textbf{Penalizacja niezweryfikowanych deklaracji} 
& Nie 
& Nie 
& Nie 
& Nie 
& \textbf{Tak} \\
\midrule

\textbf{Identyfikacja dostawców kluczowych} 
& Częściowe 
& Częściowe 
& Częściowe 
& Częściowe 
& \textbf{Tak} \\
\midrule

\textbf{Ocena planu wyjścia (exit plan)} 
& Nie 
& Nie 
& Częściowe 
& Częściowe 
& \textbf{Tak} \\
\midrule

\textbf{Transparentna, replikowalna metodyka} 
& Tak 
& Tak 
& Tak 
& Nie\newline(metodyka zastrzeżona) 
& \textbf{Tak} \\
\midrule

\textbf{Formalny model matematyczny} 
& Nie 
& Nie 
& Nie 
& Nie 
& \textbf{Tak} \\
\midrule

\textbf{Klasyfikacja ryzyka wg pochodzenia geograficznego}\newline(EOG / USA / pozostałe) 
& Nie 
& Nie 
& Nie 
& Nie 
& \textbf{Tak} \\

\end{longtable}

\begin{flushleft}
\footnotesize
\textit{Legenda:} \textbf{Tak} - wymaganie w pełni adresowane; 
\textbf{Częściowe} - wymaganie adresowane w ograniczonym zakresie lub wymagające dodatkowej konfiguracji; 
\textbf{Nie} - wymaganie nieadresowane.\\
\textit{Uwaga:} Porównanie obejmuje frameworki powszechnie stosowane w polskiej praktyce audytowej 
i compliance. Frameworki NIST SP~800-161 oraz ISACA CMMI, choć obecne w literaturze akademickiej, 
nie są szeroko stosowane jako samodzielne metodyki przez polskie organizacje 
i zostały pominięte w niniejszym zestawieniu.\\
\textit{Skróty:} CSF - \textit{Cybersecurity Framework}; 
GRC - \textit{Governance, Risk and Compliance} (ład, ryzyko i~zgodność).

\end{flushleft}



Mimo intensywności regulacyjnej prowadzonej przez regulatorów krajowych oraz międzynarodowych, brakuje usystematyzowanej metodyki, która pozwoliłaby organizacjom w sposób transparentny, co bardzo istotne powtarzalny i~oparty na obiektywnych dowodach ocenić  stan zgodności wdrożonych systemów informatycznych wraz z ich łańcuchem dostawców. Dobrze przeprowadzony audyt systemów tworzy raport czytelny i~zrozumiały dla wszystkich interesariuszy i~faktycznie wskazuje na niezgodności.
Istniejące narzędzia i~frameworki, norma ISO/IEC 27001,  NIST Cybersecurity Framework, model COBIT czy metodyki ISACA, dostarczają  wskazówek w zakresie zarządzania bezpieczeństwem informacji, jednak żaden z nich nie adresuje jednocześnie: 
\begin{itemize}
  \item perspektywy wielopodmiotowej obejmującej cały łańcuch ICT, 
  \item oceny opartej na weryfikowalnych dowodach jako mechanizmu przeciwdziałającego deklaratywnej zgodności,
  \item wieloperspektywicznej analizy uwzględniającej punkt widzenia prawnika, menedżera biznesowego i~inżyniera IT, \item rozszerzalności sektorowej pozwalającej na adaptację do specyfiki regulacyjnej różnych branż przy zachowaniu wspólnego rdzenia metodycznego.
\end{itemize}

Rezultatem jest rozdrobniona praktyka, w której oceny prawne, techniczne i~biznesowe są przeprowadzane w izolacji, podwykonawcy pozostają niesprawdzani, a wnioski z audytu opierają się na deklaracjach, a nie na weryfikowalnych dowodach.
Audyty przeprowadzane według istniejących metodyk często kończą się certyfikatami potwierdzającymi zgodność procesową, które nie odzwierciedlają faktycznego stanu bezpieczeństwa wdrożeń. Zjawisko to, zwane  \textit{compliance theater}, wynika ze słabości systemów oceny opartych na deklaracjach.  Organizacja może, z pełną znajomością swoich procedur oświadczyć, że wdrożyła wymagane środki ochrony, podczas gdy stan faktyczny jest odmienny, dokumentacja jest nieaktualna, mechanizmy kontrolne nieskuteczne, a dostawcy zewnętrzni nie podlegają żadnej weryfikacji.

\section{Postulaty i~cele ORION}

Metodologia ORION (\textit{Ocena Ryzyka Nadzorowanych Systemów Informatycznych}) nie ma na celu zastępowania metodologii audytowania zgodnych z ISO czy NIST. Wypełnia istotną lukę
w analizie wdrożeń systemów informatycznych z punktu widzenia zgodności z regulacjami. Jest to system oparty na dowodach, ponieważ, jak wykazują badania przytoczone w rozdziale
poświęconym walidacji odpowiedzi, organizacje nagminnie deklarują zgodność, której nie są w stanie udokumentować, a ta funkcja jest kluczowa.

Audyt prowadzony w ORION jest wykonywany ogólnie dla przedsiębiorstwa jako całości, analizując jego zgodność prawną z wymogami regulacyjnymi, a następnie,  zgodnie z przejściem grafu ICT, przeprowadzany jest na wszystkich systemach i~ich dostawcach. Identyfikowani są wszyscy
poddostawcy w łańcuchu ICT i~wszyscy podlegają badaniu. Wynikiem działania systemu jest określenie, który dostawca jest kluczowy, oraz wyznaczenie poziomu zgodności z prawem i~rekomendacji naprawczych z przypisanymi priorytetami.

ORION wprowadza cztery kluczowe innowacje metodyczne względem istniejących frameworków.

Pierwszą z nich jest zasada wieloperspektywiczności. ORION przyjmuje, że każdy element składowy infrastruktury ICT (system informatyczny, dostawca, wymóg, etc), jest rozumiany odmiennie przez prawnika, menedżera biznesowego, architekta oprogramowania czy administratora infrastruktury. Odmienność ta nie jest wadą, którą należy eliminować, jest naturalną cechą wynikającą z różnych wyzwań i~zadań na poszczególnych stanowiskach, którą metodyka musi uwzględniać. Perspektywa nie jest wyłącznie kategorią opisową, determinuje, kto w organizacji jest właściwą osobą do udzielenia odpowiedzi na konkretne pytanie audytowe, co eliminuje jeden z fundamentalnych błędów istniejącej praktyki: odpowiadanie na pytania prawne przez inżynierów i~na pytania techniczne przez prawników.

Drugą innowacją jest ocena oparta na dowodach. Każda odpowiedź na pytanie audytowe podlega weryfikacji przez pryzmat dostarczonego dowodu materialnego. Metodyka rozróżnia trzy poziomy jakości dowodu: brak dowodu, deklarację pisemną oraz dowód potwierdzony materialnie: log systemowy, certyfikat, wynik testu penetracyjnego, protokół z próby odtworzenia systemu.  Dodatkowo, dowód pisemny jest traktowany w zależności od jego sporządzenia, inaczej system bierze pod uwagę oświadczenie pisemne czytelnie podpisane, inaczej oświadczenie pisemne będące wynikiem anonimowego dokumentu. Brak dowodu nie zmienia odpowiedzi twierdzącej na przeczącą, lecz obniża jej efektywną wartość w modelu scoringowym w sposób kontrolowany i~konfigurowalny. Mechanizm ten precyzyjnie odzwierciedla rzeczywistość audytu: organizacja może  stosować właściwe praktyki bez ich pełnej dokumentacji, ale twierdzenie bez weryfikacji ma mniejszą wartość poznawczą niż twierdzenie udokumentowane. Jednoczenie konieczność przygotowania dowodu, co więcej, osobistego podpisania się pod dowodem, sktutjuje wysoką rzetelnością odpowiedzi. Jest to bezpośrednia odpowiedź na compliance theater.

Trzecią innowacją jest grafowe modelowanie łańcucha ICT z propagacją ryzyka. Audyt w ORION nie jest przeprowadzany wyłącznie dla podmiotu nadzorowanego traktowanego jako niepodzielna całość, obejmuje wszystkich dostawców i~podddostawców ICT, reprezentowanych jako wierzchołki skierowanego grafu zależności. Mechanizm systematycznego przejścia grafu zapewnia pełną identyfikację podmiotów w łańcuchu dostaw, które w tradycyjnych audytach pozostają poza polem widzenia. Algorytm propagacji ryzyka umożliwia z kolei obliczenie skumulowanego wpływu niezgodności na poziomie podddostawcy na ryzyko całego podmiotu nadzorowanego ,co odpowiada na pytanie, które organizacje zadają coraz częściej: czy mój dostawca jest bezpieczny, a jeśli nie, jak bardzo to mnie dotyczy?

Czwartą innowacją jest modułowa rozszerzalność sektorowa. System ORION jest zaprojektowany jako uniwersalne narzędzie z możliwością rozszerzenia na kolejne sektory. Rdzeń metodyczny, ontologia pojęć, model grafowy, formuły scoringu, struktura pytań i~raportowania, jest niezależny od branży i~może być stosowany wobec dowolnego podmiotu przetwarzającego dane w systemach ICT. Rozszerzenia domenowe, aktywowane w zależności od sektora audytowanego podmiotu, wprowadzają dodatkowe perspektywy, grupy pytań i~wymogi regulacyjne specyficzne dla danej branży. Oznacza to w praktyce, że ten system ORION może obsłużyć audyt banku, szpitala i~operatora energetycznego, różniąc się zestawem aktywnych pytań i~perspektyw, lecz zachowując identyczny model matematyczny, ontologię i~strukturę raportu. Ta cecha ORION ma istotne znaczenie nie tylko praktyczne, lecz także naukowe: pozwala na prowadzenie badań porównawczych między sektorami z użyciem jednolitej metodyki.



\section{Identyfikacja podmiotu audytowanego}
\label{sec:kim-jestem}

Przed przystąpieniem do właściwego audytu systemów informatycznych
metodologia ORION wymaga przeprowadzenia wstępnej klasyfikacji
podmiotu audytowanego. Celem tej fazy jest ustalenie, jakim
reżimem regulacyjnym podlega organizacja, w szczególności,
czy jest ona podmiotem kluczowym lub ważnym w rozumieniu dyrektywy NIS2~\autocite{dyrektywa_nis2}, podmiotem objętym zakresem stosowania rozporządzenia DORA~\autocite{dora2022}, czy też
jednostką podlegającą wyłącznie wymogom ogólnym, takim jak RODO~\autocite{rodo}. Wynik tej klasyfikacji determinuje, które moduły pytań zostaną aktywowane w dalszych etapach
audytu, jakie mnożniki wagowe zostaną przypisane do poszczególnych wymogów oraz które elementy, takie jak
plan ciągłości działania czy exit plan, mają charakter obowiązkowy, a które stanowią dobrowolne zalecenie.


Klasyfikacja podmiotu jest przeprowadzana metodą walidacji krzyżowej, gdzie zagadnienie analizowane jest niezależnie w perspektywie
zgodności, perspektywie użytkownika oraz perspektywie technicznej, a wynik jest ustalany na podstawie konsensusu odpowiedzi.
Wynikiem klasyfikacji jest przypisanie podmiotu do jednej z czterech kategorii regulacyjnych: \texttt{FINANSOWY} (podmioty objęte rozporządzeniem DORA), \texttt{KLUCZOWY} (podmioty kluczowe w rozumieniu dyrektywy NIS2 i~ustawy KSC), \texttt{WAŻNY} (podmioty ważne w rozumieniu NIS2 i~KSC) lub \texttt{POZOSTAŁE} (podmioty podlegające wyłącznie wymogom ogólnym, w~szczególności RODO). Kategoria determinuje zakres obowiązkowych elementów audytu oraz wartość mnożnika regulacyjnego $w_{reg}$, których szczegółową charakterystykę przedstawia tabela~\ref{tab:klasyfikacja-regulacyjna}.
% ============================================================
\section{Mechanizm walidacji jakości dowodów i~odpowiedzi}

\subsection{Struktura pytań}

Pytania stawiane w audycie mają z definicji wartość binarną. Oznacza to, że ankietowany moze odpowiedzieć na pytanie tylko i~wyłącznie \texttt{TAK} lub \texttt{NIE}. Każdemu pytaniu przypisana jest \texttt{ODPOWIEDŹ OCZEKIWANA},
rozumiana jako wartość wskazująca na stan zgodności z wymaganiem zgodne z semantyką treści pytania.

Działanie wartości oczekiwanej idealnie zilustruje przykład dwóch pytań:
\begin{center}
  Czy dwa plus dwa jest cztery?
\end{center}
Oraz pytanie drugie:
\begin{center}
  Czy dwa plus dwa różne od czterech?
\end{center}
W przypadku pierwszym odpowiedzią oczekiwaną jest \texttt{TAK}, w przypadku drugim  \texttt{NIE}. Współczynnik wagowy przypisywany jest wyłącznie wtedy, gdy odpowiedź jest zgodna z \texttt{ODPOWIEDZIĄ OCZEKIWANĄ}.

Zrezygnowano z opcji \texttt{NIE WIEM / BRAK DANYCH}. W sytuacji braku znajomości odpowiedzi oczekuje się, że ankietowany podejmie wysiłek uzyskania informacji. Brak pewności w odpowiedzi powinien
być rozstrzygany na niekorzyść audytu i~wynika to z konieczności ochrony organizacji.

Decyzja o rezygnacji z opcji \texttt{NIE WIEM / BRAK DANYCH} jest uzasadniona zarówno metodologicznie, jak i~z perspektywy zarządzania ryzykiem. Badania nad konstrukcją kwestionariuszy wykazują, że obecność
opcji neutralnej skłania respondentów do udzielania odpowiedzi zastępczych zamiast podejmowania wysiłku poznawczego niezbędnego do udzielenia odpowiedzi merytorycznej, zjawisko to określane jest w literaturze jako \textit{satisficing}~\autocite{Krosnick2002DontKnow}.
Kwestionariusze z wymuszonym wyborem binarnym generują dane wyższej jakości i~lepiej różnicują rzeczywiste stany wiedzy respondentów niż skale z opcją neutralną~\autocite{Dolnicar2011ForcedBinary}. 
W kontekście audytu bezpieczeństwa informatycznego brak wiedzy o stanie systemu jest sam w sobie wskaźnikiem ryzyka. Organizacja, która nie jest w stanie potwierdzić zgodności z wymaganiem, nie spełnia go z punktu
widzenia ochrony~\autocite{Waters2021DontKnow}. Przyjęcie zasady rozstrzygania wątpliwości na niekorzyść audytu odpowiada powszechnie stosowanej w zarządzaniu ryzykiem zasadzie, zgodnie z którą niepewność nie redukuje ryzyka, lecz je potęguje.

Dodatkowo, wprowadzona została opcja \texttt{flagi blokującej} $\text{FLAG}_{block}(\mathcal{S}) = \texttt{TRUE}$
 do pytań, których niespełnienie (odpowiedź inna niż \texttt{OCZEKIWANA}) oznacza przypisanie wartości zerowej dla całej sekcji. Są to pytania kluczowe, które muszą być bezwględnie spełnione.


\subsection{Model dowodowy}
\label{sec:dowody}


Każda odpowiedź na pytanie audytowe w ORION musi być opatrzona dowodem. Dowód nie zmienia wartości
odpowiedzi z \texttt{TAK} na \texttt{NIE} ani odwrotnie, lecz modyfikuje wagę odpowiedzi w modelu scoringowym.
Logika ta odzwierciedla cel audytu,  organizacja może faktycznie stosować właściwe praktyki bez ich pełnej i~faktycznej dokumentacji istnieje ryzyko, że dowód jest wadliwiwy.

Równie istotny jest aspekt behawioralny: badania z zakresu psychologii decyzji wskazują, że osoby podpisujące oświadczenia własnym imieniem i~nazwiskiem
kłamią istotnie rzadziej niż autorzy anonimowych dokumentów~\autocite{friesen2013designing}. Wymóg osobistego
podpisania dowodu jest zatem nie tylko formalny, lecz pełni funkcję mechanizmu antydefraudacyjnego wbudowanego w strukturę audytu.


\subsection{Walidacja odpowiedzi}

Walidacja odpowiedzi jest kluczowa, do prawidłowego przeprowadzenia procesu audytowania. Wskazują na to wyniki badań, które wskazują, że awersja do kłamstwa pojawia się głównie przy obowiązkowym raportowaniu \autocite{friesen2013designing}.

Osoby wypełniające dane ograniczają kłamstwo w wypełnianych ankietach, gdy są zobligowane do złożenia oświadczenia imiennego, potwierdzającego zgodne z prawdą wypełnienie ankiety. Badania wskazują, że współczynnik nieprawdziwych odpowiedzi wynosi \textit{g = 0.24}\autocite{peer2023commitment}. Współczynnik \textit{g} (zwany \textit{Hedges' g}) to standaryzowana miara wielkości efektu,  wyraża jak duża jest różnica między grupami w jednostkach odchylenia standardowego.


\subsection{Skala jakości dowodów}

Metodologia ORION wymaga przeprowadzenia odpowiedniej walidacji dowodu i~jakość dowodu, co bezpośrednio wpływa na wagę odpowiedzi.
ORION wyróżnia cztery poziomy jakości dowodu,zdefiniowane w tabeli \ref{tab:dowody}.
oznaczane symbolem $d \in \{0;\ 0{,}5;\ 0{,}8;\ 1{,}0\}$.


\begin{table}[H]
\label{sec:skala-dowodow}
\centering
\caption{Skala jakości dowodów w metodologii ORION}
\label{tab:dowody}
\renewcommand{\arraystretch}{1.6}
\begin{tabular}{p{1.5cm}p{2.5cm}p{4.5cm}p{4.5cm}}
\toprule
\textbf{Poziom} & \textbf{Symbol $d$}
& \textbf{Definicja}
& \textbf{Przykłady} \\
\midrule
Brak dowodu & $d = 0$
& Odpowiedź udzielona wyłącznie ustnie lub pisemnie
  bez jakiegokolwiek dokumentu potwierdzającego;
  niemożliwa do weryfikacji przez audytora.
& Oświadczenie słowne podczas wywiadu,
  odpowiedź wpisana bezpośrednio w formularz
  bez załącznika \\
\midrule
Dowód słaby & $d = 0{,}4$
& Dokument pisemny bez wskazania autora lub
  anonimowy; dokument ogólny niepowiązany
  wprost z audytowanym systemem.
& Regulamin ogólny dostawcy, polityka
  bezpieczeństwa bez daty i~podpisu,
  screen z systemu bez kontekstu \\
\midrule
Dowód średni & $d = 0{,}7$
& Oświadczenie pisemne z czytelnym podpisem
  osoby fizycznej i~datą; dokument wskazujący
  jednoznacznie na audytowany system.
& Podpisana deklaracja zgodności,
  umowa powierzenia z klauzulą TIA,
  certyfikat ISO z numerem seryjnym \\
\midrule
Dowód silny & $d = 1{,}0$
& Dowód materialny niemożliwy do sfałszowania
  bez pozostawienia śladu technicznego;
  log systemowy, wynik testu automatycznego,
  lub oświadczenie podpisane z poświadczeniem
  tożsamości wskazujące na pytanie audytowe.
& Log z systemu SIEM, wynik testu
  penetracyjnego, protokół z próby odtworzenia
  z BCP, oświadczenie potwierdzone podpisem osobistym. \\
\bottomrule
\end{tabular}
\end{table}

\subsection{Wpływ dowodu na wartość odpowiedzi}
\label{sec:wplyw-dowodu}

Efektywna wartość odpowiedzi na pytanie audytowe
jest iloczynem odpowiedzi binarnej i~jakości dowodu:

\begin{equation}
V_{eff}(q_j, \pi_i) = a_{ij} \cdot d_{ij}
\label{eq:veff}
\end{equation}

\noindent gdzie $a_{ij} \in \{0, 1\}$ jest odpowiedzią
perspektywy $\pi_i$ na pytanie $q_j$, a
$d_{ij} \in \{0;\ 0{,}4;\ 0{,}7;\ 1{,}0\}$
jest jakością dostarczonego dowodu.

Konsekwencje praktyczne wzoru~\eqref{eq:veff}
są następujące:
\begin{itemize}
  \item odpowiedź TAK bez dowodu ($a=1$, $d=0$)
  daje $V_{eff} = 0$ jest traktowana identycznie jak odpowiedź NIE,
  \item odpowiedź TAK z dowodem silnym ($a=1$, $d=1{,}0$)
  daje $V_{eff} = 1{,}0$  pełna wartość,
  \item odpowiedź NIE z dowodem silnym ($a=0$, $d=1{,}0$) daje $V_{eff} = 0$ wskazuje na niezgodność między odpowiedzią, a dowodem.
\end{itemize}

Mechanizm ten precyzyjnie odzwierciedla cel audytu: nie interesuje nas, co organizacja mówi, że robi, interesuje nas, co jest w stanie udowodnić.



\subsection{Walidacja krzyżowa oraz definicja konsensusu}
\label{sec:crossvalidation}
Walidacja krzyżowa (\textit{cross-validation}) jest kluczowym mechanizmem metodologii ORION, zapewniającym rzetelność klasyfikacji atrybutów i~odpowiedzi na pytania audytowe. To samo zagadnienie przedstawione jest użytkownikom z co najmniej dwóch perspektyw z trzech, a wynik jest ustalany na podstawie porównania odpowiedzi.
Mechanizm ten realizuje dwa cele jednocześnie: po pierwsze, wykrywa rozbieżności między perspektywami, które są
same w sobie wartościowymi sygnałami audytowymi; po drugie, zwiększa odporność systemu na celowe lub niecelowe podawanie nieprawdziwych informacji przez jeden z aktorów.


\begin{definition}[Formalna definicja konsensusu]
\label{sec:konsensus}
Niech $Q$ będzie zagadnieniem audytowym zadanym
perspektywom ze zbioru $\Pi_Q \subseteq \Pi$,
gdzie $|\Pi_Q| \geq 2$. Każda perspektywa
$\pi_i \in \Pi_Q$ udziela odpowiedzi
$a_i \in \{0, 1\}$ (NIE / TAK). Wynik
cross-validation zagadnienia $Q$ definiuje się jako:

\begin{equation}
\text{CV}(Q) =
\begin{cases}
1 & \text{jeśli } \left|\{i : a_i = 1\}\right|
    > \left|\{i : a_i = 0\}\right| \\
0 & \text{jeśli } \left|\{i : a_i = 0\}\right|
    > \left|\{i : a_i = 1\}\right| \\
\varnothing & \text{jeśli remis lub brak odpowiedzi}
\end{cases}
\label{eq:cv}
\end{equation}

Gdy $\text{CV}(Q) = \varnothing$, zagadnienie jest
oznaczane flagą \textsc{Rozbieżność} i~nie może
zostać zatwierdzone bez interwencji audytora.
W takim przypadku efektywna wartość zagadnienia
jest penalizowana:

\begin{equation}
V_{eff}(Q) = \min_{i}(a_i) \cdot \beta_{cv}
\label{eq:veff-rozbiez}
\end{equation}

\noindent gdzie $\beta_{cv} \in [0,1]$ jest
konfigurowalnym współczynnikiem kary za brak
konsensusu (wartość domyślna: $\beta_{cv} = 0{,}5$).

\end{definition}
Innymi słowy: brak zgody między perspektywami
jest traktowany jako połowicznie negatywna odpowiedź,
nawet jeśli jeden z aktorów odpowiedział twierdząco.

Gdy konsensus jest osiągnięty i~wszystkie perspektywy
dostarczyły dowody najwyższej jakości, efektywna wartość
jest dodatkowo premiowana:

\begin{equation}
V_{eff}(Q) = \text{CV}(Q) \cdot
\Bigl(1 + \gamma \cdot
\mathbb{1}\bigl[\forall\, i~: d_i = d_{max}\bigr]\Bigr)
\label{eq:veff-premia}
\end{equation}

\noindent gdzie $\gamma = 0{,}1$ jest premią za
pełny konsensus z kompletem dowodów najwyższej jakości,
a $d_{max}$ jest najwyższym poziomem dowodu
w skali opisanej w sekcji~\ref{sec:dowody}.

\subsection{Interpretacja rozbieżności}
\label{sec:interpretacja-rozbiez}

Rozbieżność między perspektywami nie jest błędem
procesu audytowego, jest wartościową informacją
merytoryczną. Tabela~\ref{tab:rozbiez-interpretacja}
przedstawia najczęstsze wzorce rozbieżności
i ich standardową interpretację w ORION.

\begin{table}[H]
\centering
\caption{Wzorce rozbieżności cross-validation
         i~ich interpretacja audytowa}
\label{tab:rozbiez-interpretacja}
\renewcommand{\arraystretch}{1.5}
\begin{tabular}{p{3.5cm}p{3.5cm}p{6cm}}
\toprule
\textbf{Wzorzec rozbieżności}
& \textbf{Typowa przyczyna}
& \textbf{Interpretacja i~zalecane działanie} \\
\midrule
Zgodność: TAK\newline Użytkownik: NIE\newline Techniczna: NIE
& Wymóg zdefiniowany w przepisach, lecz nieimplementowany
& Compliance theater: organizacja deklaruje zgodność formalną
  bez rzeczywistego wdrożenia. Flaga P1. \\
\midrule
Zgodność: NIE\newline Użytkownik: TAK\newline Techniczna: TAK
& Praktyka wdrożona bez formalnej podstawy prawnej
& Brak świadomości regulacyjnej. Ryzyko braku dokumentacji
  wymaganej przez regulatora. Rekomendacja formalizacji. \\
\midrule
Zgodność: TAK\newline Użytkownik: TAK\newline Techniczna: NIE
& Wymaganie znane, polityka istnieje, wdrożenie niezrealizowane
& Luka techniczna. Wymóg zdefiniowany lecz nieimplementowany
  na poziomie systemów. Flaga P2. \\
\midrule
Zgodność: NIE\newline Użytkownik: NIE\newline Techniczna: TAK
& System techniczny bez świadomości biznesowej i~prawnej
& Shadow IT lub wdrożenie bez akceptacji zarządczej.
  Wymaga formalnej inwentaryzacji i~decyzji. \\
\bottomrule
\end{tabular}
\end{table}



\section{Definicje i~aktorzy}

\subsection{Perspektywa}
\label{perspektywa-pytań}
Perspektywa punktu widzenia systemu określa interesariuszy uczestniczących w ocenie oraz sposób ich rozumienia analizowanych zagadnień. Każde pytanie ankietowe jest przypisane do co najmniej dwóch perspektyw, o ile jest to możliwe, do wszystkich trzech. Umożliwia to prowadzenie analizy porównawczej odpowiedzi na to samo kryterium, udzielonych niezależnie przez różnych interesariuszy. Rozbieżności między perspektywami stanowią sygnał do dalszego wyjaśnienia i~są równie wartościowe analitycznie jak odpowiedzi zgodne.

Zdefiniowane są trzy perspektywy:
\begin{enumerate}[label=(\roman*)]
  \item \textit{Perspektywa zgodności}, reprezentuje punkt widzenia osób odpowiedzialnych za przestrzeganie regulacji, takich jak prawnik, inspektor ochrony danych osobowych (IOD/DPO) czy specjalista ds. compliance. Celem tej perspektywy jest ocena kryterium pod kątem zgodności z obowiązującymi przepisami prawa, regulacjami rynkowymi oraz wewnętrznymi politykami organizacji (np. RODO, NIS2, DORA).

  \item  \textit{Perspektywa użytkownika}, reprezentuje jednostki organizacyjne firmy aktywnie korzystające z systemu w codziennej pracy, takie jak właściciel biznesowy systemu, kierownik działu lub kluczowy użytkownik końcowy. Celem tej perspektywy jest uzyskanie odpowiedzi odzwierciedlającej rzeczywisty sposób wykorzystania systemu, nie jego założenia projektowe, lecz faktyczny stan i~pola eksploatacji.

  \item \textit{Perspektywa techniczna}, reprezentuje punkt widzenia osób odpowiedzialnych za budowę, utrzymanie i~bezpieczeństwo systemu, takich jak architekt systemów, administrator lub inżynier. Celem tej perspektywy jest ocena, czy implementacja i~wdrożenie systemu zostały przeprowadzone zgodnie z obowiązującymi standardami technicznymi i~najlepszymi praktykami branżowymi.
\end{enumerate}


Przyjęte perspektywy determinują strukturę pytań ankietowych, w których pojedyncze kryterium analizowane jest przez pryzmat trzech różnych ujęć oraz wpływają na sposób definiowania obiektów. Ilustrując to podejście przykładem, zadajmy sobie pytanie:
\begin{center}
\vspace{1em} % Dodaje odstęp pionowy 
\textit{Jaka jest definicja systemu CRM?}
\end{center}


Na tak postawione pytanie zostaną udzielone trzy komplementarne odpowiedzi:

\begin{enumerate}[label=(\roman*)]
    \item Z punktu widzenia \textit{użytkownika}, system CRM służy do rejestrowania wszyskich kontaktów z klientami oraz realizacji procesów sprzedażowych.
    \item Z punktu widzenia \textit{zgodności}, system oraz jego fizyczna lokalizacja muszą gwarantować zgodność z polityką RODO, zapewniając organizacji pełną wiedzę o miejscu przetwarzania danych osobowych.
    \item Z punktu widzenia \textit{technicznego}, jest to rozwiązanie informatyczne, które ma zapewniać dostępność zasobów dla uprawnionych użytkowników poprzez określone kanały komunikacyjne.
\end{enumerate}


\subsection{Obiekt}

Centralnym pojęciem metodyki ORION jest \texttt{obiekt}, reprezentujący  byt w systemie ICT organizacji, który może być przedmiotem  oceny zgodności regulacyjnej. Pojęcie obiektu jest szerokie: obejmuje zarówno systemy informatyczne działające 
na infrastrukturze własnej organizacji, jak i~zewnętrznych dostawców usług, platformy chmurowe oraz poddostawców odkrytych w toku rozszerzania grafu ICT.

\begin{definition}[Obiekt]
Obiekt jest formalną jednostką audytową systemu ORION, 
reprezentującą byt w ekosystemie ICT organizacji, który:
\begin{enumerate}[label=(\roman*)]
    \item posiada tożsamość możliwą do jednoznacznej identyfikacji w grafie zależności ICT,
    \item może być przypisany do co najmniej jednego właściciela  lub podmiotu odpowiedzialnego,
    \item podlega co najmniej jednemu wymogowi regulacyjnemu wynikającemu z obowiązujących przepisów prawa,
    \item może być oceniany z co najmniej dwóch perspektyw audytowych zdefiniowanych w zbiorze $\Pi$.
\end{enumerate}
Formalnie, obiekt jest wierzchołkiem $v \in V$ grafu ICT 
$G = (V, E, W)$, posiadającym zbiór atrybutów  klasyfikacyjnych $\mathcal{A}(v)$ determinujących zakres i~intensywność przeprowadzanego audytu.
\end{definition}

W ORION obiektem jest zarówno aplikacja webowa obsługująca klientów banku, dostawca chmury publicznej, na której ta aplikacja działa, jak i~usługi firmy
 świadczącej usługi zarządzania kluczami kryptograficznymi na rzecz tego dostawcy chmury. Każdy z tych 
bytów jest odrębnym wierzchołkiem grafu, odrębnym podmiotem  audytu i~odrębnym źródłem ryzyka, choć wszystkie wymione podmioty uczestniczą w tym samym łańcuchu przetwarzania danych.

\subsubsection{Typy obiektów}

ORION wyróżnia cztery typy obiektów, które determinują domyślny zestaw aktywowanych sekcji pytań oraz rolę wierzchołka w grafie ICT.
Każdy obiekt posiada właściwości, czyli pola przyjmujące definiowane przez audytora wartości:

\begin{table}[H]
\centering
\caption{Typy obiektów w systemie ORION}
\label{tab:typy-artefaktow}
\renewcommand{\arraystretch}{1.5}
\begin{tabular}{p{3cm}p{7cm}p{5cm}}
\toprule
\textbf{Typ} & \textbf{Definicja} & \textbf{Przykłady} \\
\midrule
\texttt{PODMIOT} 
\newline Właściwości: nazwa
& Korzeń grafu ICT ($v_0$); audytowana organizacja jako całość. 
  Nie jest bezpośrednio oceniany pytaniami szczegółowymi, 
  jego wynik wynika z agregacji wszystkich pozostałych obiektów.
& Bank, szpital, operator infrastruktury krytycznej \\

\midrule
\texttt{SYSTEM\_INFORMATYCZNY}
\newline Właściwości: nazwa, dostawca, właściciel biznesowy
&  Oprogramowanie służące do gromadzenia, przetwarzania i~przesyłania informacji.
& System kadrowo-płacowy, baza danych pacjentów, 
  portal intranetowy, API płatności \\
\midrule

\texttt{DOSTAWCA} 
\newline Właściwości: nazwa, kraj pochodzenia, poddostawca
& Zewnętrzny podmiot prawny świadczący usługi ICT na rzecz 
  organizacji na podstawie umowy; odkryty na pierwszym 
  poziomie grafu.
& Dostawca ERP, operator chmury publicznej, 
  firma hostingowa, integrator systemów \\
\midrule
\texttt{PODDOSTAWCA}
\newline Właściwości: nazwa, kraj pochodzenia
& Zewnętrzny podmiot prawny świadczący usługi ICT 
  na rzecz dostawcy lub innego poddostawcy; 
  odkryty przez mechanizm triggera TIA lub deklarację 
  dostawcy wyższego poziomu.
& Google LLC jako operator infrastruktury dostawcy ERP, 
  Cloudflare jako CDN aplikacji webowej, 
  firma zarządzająca kluczami kryptograficznymi \\
\bottomrule
\end{tabular}
\end{table}

\subsection{Atrybuty obiektów}


Każdy obiektów jest scharakteryzowany przez zbiór atrybutów 
klasyfikacyjnych $\mathcal{A}(v)$, które pełnią dwie funkcje. 
Po pierwsze, determinują, które sekcje pytań audytowych 
zostaną aktywowane dla danego obiektów, na przykład 
system wewnętrzny działający wyłącznie on-premise nie 
otrzyma pytań dotyczących transferu 
danych do chmury. Po drugie, wyznaczają wagi stosowane 
w modelu scoringowym, decydując o tym, jak mocno wynik 
danego obiektu wpływa na wynik całego podmiotu.

\subsection{Atrybut: dostępność systemu}

Atrybuty obiektów są ustalane w fazie wstępnej audytu za pomocą pomocniczej ankiety, 
przed wygenerowaniem kwestionariusza szczegółowego. 
Widoczność systemu i~jego krytyczność, jest wyznaczana przez mechanizm  cross-validation z udziałem co najmniej dwóch perspektyw.


\begin{definition}[Poziom dostępności systemu]
Poziom dostępności obiektu $v$ jest wartością ze zbioru:
\[
\textit{dostępność}(v) \in 
\{\texttt{PRYWATNY},\ \texttt{PUBLICZNY}\}
\]
wyznaczaną przez mechanizm cross-validation na podstawie 
odpowiedzi co najmniej dwóch z trzech perspektyw rdzeniowych: 
\texttt{ZGODNOŚCI}, \texttt{UŻYTKOWNIKA} i~\texttt{TECHNICZNEJ}. 
\end{definition}


\begin{table}[H]
\centering
\caption{Poziomy dostępności obiektu, kryteria klasyfikacji
         i~mnożniki wagowe}
\label{tab:dostepnosc}
\renewcommand{\arraystretch}{1.5}
\begin{tabular}{p{2.8cm}p{1.5cm}p{4.5cm}p{4cm}}
\hline
\textbf{Poziom} & \textbf{$w_{krit}$} 
& \textbf{Kryteria klasyfikacji} 
& \textbf{Przykład} \\
\toprule
\texttt{PRYWATNY} & $1{,}0$ 
& System informatyczny działający na infrastrukturze 
  własnej organizacji lub zarządzanej wyłącznie przez nią; 
  bez bezpośredniego dostępu z sieci publicznej.
& System kadrowo-płacowy on-premise, baza danych pacjentów, 
  wewnętrzny portal intranetowy \\
\midrule
\texttt{PUBLICZNY} & $2{,}5$ 
& System informatyczny dostępny z sieci publicznej lub 
  świadczący usługi na rzecz podmiotów zewnętrznych; 
  obejmuje systemy SaaS, API i~aplikacje webowe.
& Portal klienta banku, system e-rejestracji, API płatności \\

\bottomrule
\end{tabular}
\end{table}

Klasyfikacja jest przeprowadzana przez cztery kryteria, z których każde jest zadawane wszystkim trzem perspektywom.
Logika wynikowa jest progowa: nawet jedna odpowiedź twierdząca z dowolnej perspektywy na dowolne kryterium
skutkuje klasyfikacją \texttt{PUBLICZNY}. Uzasadnienie tej asymetrii jest regulacyjne: jedna niepotwierdzona
ścieżka dostępu zewnętrznego jest wystarczającym powodem do traktowania systemu jako publicznie dostępnego
i stosowania względem niego pełnego zestawu wymogów bezpieczeństwa dla systemów publicznych.

\subsection{Atrybut: poziom krytyczności}

Poziom krytyczności jest atrybutem klasyfikacyjnym obiektu 
określającym, jak poważne skutki operacyjne, prawne i~społeczne 
pociągnęłaby za sobą niedostępność, naruszenie integralności 
lub utrata poufności danego systemu lub dostawcy. Atrybut ten 
jest wyznaczany w fazie wstępnej audytu metodą cross-validation, to samo zagadnienie jest zadawane niezależnie trzem perspektywom audytowym, a wynik klasyfikacji jest ustalany na podstawie 
konsensusu odpowiedzi.

\begin{definition}[Poziom krytyczności obiektu]
Poziom krytyczności obiektu $v$ jest wartością ze zbioru:
\[
\textit{krytyczność}(v) \in 
\{\texttt{KRYTYCZNY},\ \texttt{WAŻNY},\ \texttt{STANDARDOWY}\}
\]
wyznaczaną przez mechanizm cross-validation na podstawie 
odpowiedzi co najmniej dwóch z trzech perspektyw rdzeniowych: 
\texttt{ZGODNOŚCI}, \texttt{UŻYTKOWNIKA} i~\texttt{TECHNICZNEJ}. 
Poziom krytyczności determinuje mnożnik wagowy $w_{krit}(v)$ 
stosowany w równaniu agregacji wyników na poziomie 
podmiotu~\eqref{eq:Spodmiot_rozszerzone}.
\end{definition}

\begin{table}[H]
\centering
\caption{Poziomy krytyczności obiektu, kryteria klasyfikacji 
         i~mnożniki wagowe}
\label{tab:krytycznosc}
\renewcommand{\arraystretch}{1.5}
\begin{tabular}{p{2.8cm}p{1.5cm}p{4.5cm}p{4cm}}
\hline
\textbf{Poziom} & \textbf{$w_{krit}$} 
& \textbf{Kryteria klasyfikacji} 
& \textbf{Podstawa regulacyjna} \\
\toprule
\texttt{KRYTYCZNY} & $2{,}0$ 
& Niedostępność przez ponad 4~h powoduje bezpośrednie straty 
  operacyjne, naruszenie przepisów prawa lub zagrożenie 
  bezpieczeństwa osób; brak możliwości zastąpienia; 
  system wspiera funkcję kluczową w rozumieniu NIS2 lub DORA
& NIS2 Art.~3; DORA Art.~2; 
  ustawa o krajowym systemie cyberbezpieczeństwa \\
\midrule
\texttt{WAŻNY} & $1{,}5$ 
& Niedostępność poważnie utrudnia działalność, lecz jej 
  nie wstrzymuje; istnieje procedura zastępcza; 
  system wspiera funkcję ważną w rozumieniu NIS2
& NIS2 Art.~3 ust.~2; 
  ISO/IEC~22301 \\
\midrule
\texttt{STANDARDOWY} & $1{,}0$ 
& Niedostępność powoduje utrudnienia pomocnicze; 
  istnieje możliwość pracy bez systemu przez 
  co najmniej 48~h bez istotnych strat
& --- \\
\bottomrule
\end{tabular}
\end{table}

Klasyfikacja krytyczności jest przeprowadzana przez sześć kryteriów (K1--K6) zadawanych wszystkim trzem perspektywom. Wynik końcowy wyznaczany jest progowo: suma wszystkich odpowiedzi twierdzących ze wszystkich perspektyw i~wszystkich kryteriów jest normalizowana do skali klasyfikacyjnej zgodnie z~ankietą przedstawioną w~sekcji~\ref{sec:krok23}.

\subsection{Atrybut: klasyfikacja utrzymania w chmurze}
\label{sec:klasyfikacja-chmura}
Lokalizacja miejsca utrzymania usługi ma kluczowy wpływ na warunki jej świadczenia. Oprogramowanie działające \textit{on-premise}, czyli w siedzibie organizacji, nie jest narażone na ryzyko przesyłu danych do zewnętrznych dostawców. Istotna jest również prawna jurysdykcja podmiotu, który hostuje dane. Fakt przechowywania danych w chmurze będzie implikować konieczność dokładnego audytu łańcucha dostawców i~wykonania oceny skutków transferu (\textit{TIA}), wymaganej przez RODO w sytuacji, gdy dane przesyłane są z Europejskiego Obszaru Gospodarczego (\textit{EOG}) do tzw. "państwa trzeciego" (np. USA, Indie, Chiny).

\subsubsection*{Klasyfikacja utrzymania  w chmurze}
Metodologia ORION definiuje dwa typy lokalizacji danych, a ich atrybut wpływa na aktywacje sekcji pytań dotyczących rozwiązań chmurowych oraz transferu danych.


\begin{table}[H]
\centering
\caption{Klasyfikacja utrzymania w chmurze w ORION}
\label{tab:typy-chmura}
\renewcommand{\arraystretch}{1.4}
\begin{tabular}{p{5cm}p{8cm}}
\toprule
\textbf{Typ utrzymania} & \textbf{Opis}\\
\midrule
\texttt{CHMURA\_K3} & Dane przechowywane są poza organizacją u zewnętrznego podmiotu (SaaS, IaaS/PaaS, model hybrydowy, hosting, kolokacja), którego lokalizacja lub jurysdykcja nie gwarantuje przechowywania danych w EOG. Obejmuje dostawców podlegających prawu państw trzecich, np. Amazon Web Services, Google Cloud, Microsoft Azure.\\
\midrule
\texttt{CHMURA\_EOG} & Dane przechowywane są poza organizacją u zewnętrznego podmiotu (SaaS, IaaS/PaaS, model hybrydowy, hosting, kolokacja), którego lokalizacja i~jurysdykcja gwarantuje przechowywanie danych wyłącznie w EOG, np. NASK SA, Exea Data Center, Beyond.pl, OVH (region PL/EU).\\
\midrule
\texttt{ON\_PREMISE} & Dane przechowywane są wyłącznie na serwerach i~infrastrukturze będącej własnością organizacji, fizycznie zlokalizowanej w jej siedzibie lub własnym centrum danych.\\
\bottomrule
\end{tabular}
\end{table}

\subsection{Atrybut: klasyfikacja atrybutu typów przetwarzanych danych}
\label{sec:klasyfikacja-danych}

Dotychczasowe modele oceny ryzyka ICT nie uwzględniają charakter  przetwarzanych danych jako samodzielny wymiar ryzyka z formalnym  przełożeniem na wynik scoringowy. W metodologii ORION charakter danych jest traktowany jako atrybut pierwszorzędny, który wpływa bezpośrednio  na wagę artefaktu w agregacji końcowej. Logika jest prosta: naruszenie bezpieczeństwa systemu przetwarzającego tajemnicę lekarską lub dane finansowe klientów pociąga za sobą nieporównywalnie poważniejsze  skutki prawne, finansowe i~społeczne niż naruszenie systemu  obsługującego wyłącznie dane wewnętrzne organizacji.

\subsubsection{Klasyfikacja typów danych}

Metodologia ORION wyróżnia cztery typy przetwarzanych danych, uporządkowane  rosnąco według wrażliwości prawnej i~społecznej. Każdemu typowi  przypisany jest mnożnik wagowy $w_{data}$, który modyfikuje 
wpływ artefaktu na wynik podmiotu zgodnie 
z równaniem~\eqref{eq:Spodmiot_rozszerzone}.

\begin{table}[H]
\centering
\caption{Klasyfikacja typów danych w ORION i~odpowiadające mnożniki wagowe}
\label{tab:typy-danych}
\renewcommand{\arraystretch}{1.4}
\begin{tabular}{p{2.5cm}p{1.5cm}p{1.5cm}p{3.5cm}p{4.5cm}}
\toprule
\textbf{Typ danych} & \textbf{Symbol} & \textbf{$w_{data}$} 
& \textbf{Podstawa prawna} & \textbf{Przykłady} \\
\midrule
\texttt{BRAK} & $T_0$ & $1{,}0$ 
& --- 
& Dane konfiguracyjne, logi techniczne bez PII \\
\midrule
\texttt{OSOBOWE} & $T_1$ & $1{,}3$ 
& RODO Art.~4 ust.~1 
& Imię, adres, adres e-mail, numer IP \\
\midrule
\texttt{WRAŻLIWE} & $T_2$ & $1{,}6$ 
& RODO Art.~9 ust.~1 
& Stan zdrowia, dane biometryczne, 
  przekonania religijne, orientacja seksualna \\
\midrule
\texttt{TAJEMNICA} & $T_3$ & $2{,}0$ 
& Ustawa sektorowa\newline(bankowa, lekarska, adwokacka) 
& Tajemnica bankowa, tajemnica lekarska, 
  tajemnica zawodowa, dane niejawne \\
\bottomrule
\end{tabular}
\end{table}

\subsection{Atrybut: kraj pochodzenia}
\label{sec:kraj-pochodzenia}
Atrybut ten określa kraj pochodzenia podmiotu w celu jednoznacznego ustalenia jego statusu prawnego, tj. czy jest to podmiot mający siedzibę na terenie Europejskiego Obszaru Gospodarczego (EOG), czy też kwalifikuje się on jako podmiot z kraju trzeciego w myśl obowiązujących przepisów prawa (w szczególności w kontekście regulacji dotyczących ochrony danych osobowych RODO oraz cyberbezpieczeństwa).

Wartość tego atrybutu ma kluczowe znaczenie dla logiki warunkowego uruchamiania sekcji pytań podczas procesu audytowego. W przypadku zidentyfikowania podmiotu pochodzącego z kraju spoza EOG, system automatycznie rozszerza zakres audytu o dodatkowe wymagania. Wynika to z podwyższonego ryzyka operacyjnego i~prawnego (np. w obszarze transferu danych lub łańcucha dostaw). W konsekwencji na organizację audytującą nakładany jest obowiązek przeprowadzenia pogłębionej weryfikacji, obejmującej m.in. szczegółowe badanie mechanizmów zabezpieczeń, umów transferowych oraz analizę wpływu na ryzyko (np. TIA).


\begin{table}[H]
\centering % Opcjonalnie: wyśrodkowanie tabeli na stronie
\caption{Dopuszczalne wartości pochodzenia przedsiębiorstw }
\label{tab:country-values}

\resizebox{\textwidth}{!}{%
\begin{tabular}{p{5cm} p{11cm}}
\toprule
\textbf{Kod kraju} & \textbf{Opis i~zakres odpowiedzialności} \\
\midrule
\texttt{KOD ISO} & Kod kraju pochodzenia, np. PL, DE, US \\
\bottomrule
\end{tabular}%
}
\end{table}

\section{Plan ciągłości dzialania}

Plan ciągłości działania jest procedurą wymaganą od podmiotów kluczowych zarówno na gruncie ustawy o KSC\autocite{uksc2026}, jak i~rozporządzenia DORA\autocite{dora2022}.

\begin{definition}[Plan ciągłości działania (BCP)]
Zbiór udokumentowanych procedur, które określają działania, jakie podmiot musi podjąć w celu odzyskania zdolności do wykonywania operacji oraz zapewnienia ciągłości procesów krytycznych na akceptowalnym poziomie po wystąpieniu incydentu zakłócającego.
\end{definition}

Podstawą planu ciągłości działania jest określenie dwóch wskaźników odporności. Pierwszy, RTO (\textit{Recovery Time Objective}) określa maksymalny dopuszczalny czas, przez jaki system może pozostawać niedostępny po wystąpieniu awarii, im niższy, tym wyższe wymagania wobec infrastruktury. Drugi, RPO (Recovery\textit{ Point Objective}) wyznacza  maksymalną akceptowalną utratę danych, wyrażoną jako punkt w czasie, do którego dane muszą zostać odtworzone, ilustrując przykładem, RPO wynoszące 15 minut oznacza, że żadna transakcja starsza niż kwadrans nie może zostać bezpowrotnie utracona.

Plan musi przewidywać konkretne scenariusze postępowania na wypadek różnych klas awarii systemów IT, ataku ransomware czy niedostępności kluczowego personelu. Równolegle obowiązują procedury wykrywania, klasyfikacji i~eskalacji incydentów, w tym raportowania do organów nadzorczych, takich jak KNF, w ściśle określonych ramach czasowych wynikających z obiwązującego prawa, a ich definicje szczegółowo określa DORA i~KSC.

Niezbędne jest opracowanie polityki wytwarzania oraz przywracania kopii zapasowych (DRP). Kopie zapasowe muszą być fizycznie lub logicznie odizolowane od głównej infrastuktury produkcyjnej, co chroni jest przed zaszyfrowaniem w przypadku ataku ransomware. Plan odtwarzania precyzuje harmonogram tworzenia kopii, lokalizację ich przechowywania, odpowiedzialnych za proces oraz procedury testowego odtwarzania danych w kontrolowanych warunkach.

Bardzo ważnym elementem plan ciągłości działania jest strategia komunikacji na czas kryzysu, który określa, kto, kiedy i~w jaki sposób informuje klientów, pracowników, kontrahentów oraz nadzór o wystąpieniu awarii i~przewidywanym czasie jej usunięcia. 

Zgodnie z wymogami DORA, instytucja finansowa nie może ograniczyć planowania ciągłości wyłącznie do własnych zasobów, obowiązkiem jest zapewnienie, że dostawcy usług chmurowych i~IT również dysponują udokumentowanymi planami ciągłości i~są w stanie dotrzymać uzgodnionych poziomów usług (\textit{SLA}) w sytuacji kryzysowej. Wymóg ten powinien być wprost zawarty w umowach z dostawcami, wraz z prawem do audytu.

Plan ciągłości działania podlega obowiązkowym testom ich skuteczności, które obejmują symulacje cyberataków, przełączanie ruchu na ośrodek zapasowy oraz testy odtwarzania danych z kopii zapasowych. Wyniki każdego testu powinny skutkować aktualizacją dokumentacji.

Odpowiedzialność za zatwierdzenie i~skuteczność planów ciągłości działąnia, spoczywa na zarządzie instytucji. Od strony technicznej plan musi eliminować Single Point of Failure (\textit{SPOF}) poprzez zapewnienie redundantnych systemów, łączy komunikacyjnych i~zasilania, tak aby awaria pojedynczego komponentu nie prowadziła do niedostępności całej usługi.

W metodoligii ORION, plan ciągłości działania jest wymagany dla podmiotów określonych w KSC oraz DORA. Jego brak lub nieprawidłowe zdefiniowanie w organizacjach, gdzie jest on obowiązkowy wskazują binarnie na uchybienia w działaniu organizacji.

\begin{table}[H]
\centering
\caption{Mnożniki wagowe BCP wg typu organizacji}
\label{tab:bcp-wagi}
\renewcommand{\arraystretch}{1.4}
\begin{tabular}{p{4.5cm}p{1.5cm}p{2.5cm}p{4.5cm}}
\toprule
\textbf{Typ organizacji} & \textbf{Symbol} & \textbf{$w_{reg}$} & Obowiązkowe\\
\midrule
\texttt{FINANSOWY} & $D_0$ & $1{,}0$   & tak\\
\midrule
\texttt{KLUCZOWY} & $D_1$ & $1{,}0$   & tak\\
\midrule
\texttt{WAŻNY} & $D_2$ & $1{,}0$   & tak\\
\midrule
\texttt{POZOSTAŁE} & $D_3$ & $0{,}5$ & nie \\
\bottomrule
\end{tabular}
\end{table}

\section{Exit plan}

Exit plan (plan wyjścia) to obowiązkowy element zarządzania dostawcami IT w banku, wymagany przez rozporządzenie DORA oraz ustawę KSC.

\begin{definition}[Exit plan]
Strategia wyjścia z usługi pozwalająca organizacji wyłączyć system, zakończyć współpracę z dotychczasowym dostawcą oraz ponownie uruchomić usługę u innego dostawcy przy zachowaniu ciągłości działania lub minimalizacji przestoju.
\end{definition}

Plan wyjścia stanowi pisemnie udokumentowaną procedurę umożliwiającą organizacji zakończenie współpracy z dostawcą ICT i~przeniesienie usług do alternatywnego środowiska, bez przestoju operacyjnego i~bez utraty danych. Definicja ma zagwarantować ciągłości działania, czyli zapewnić minimalny czas niedostępności usługi dla jej użytkowników. Równie istotna jest ochrona przed zjawiskiem \textit{vendor lock-in}, czyli uzależnienia od jednego dostawcy

Wedługi definicji DORA, plan musi identyfikować musi wskazywać rzeczywistego dostawcę zastępczego lub opisywać scenariusz powrotu do infrastruktury własnej. Co istotne, umowa z dostawcą musi zawierać klauzulę zobowiązującą go do udostępnienia wszystkich danych w formacie umożliwiającym ich import do innego systemu

Dostawca ICT jest zobowiązany do świadczenia usług przez określony czas po wypowiedzeniu umowy, tak aby organizacja przeprowadziła migrację danych. Co istotne, wsparcie to nie może mieć charakteru pasywnego, dostawca ma obowiązek aktywnej pomocy technicznej przy przenoszeniu procesów, danych i~konfiguracji do nowego środowiska. Tego rodzaju zobowiązania muszą być precyzyjnie zapisane w umowie, z określonymi terminami i~zakresem odpowiedzialności.

Exit plan traci swoją wartość, jeśli nie jest regularnie weryfikowany. DORA wymaga przeprowadzania testów planu co najmniej raz w roku lub każdorazowo po istotnych zmianach w architekturze systemu (Art. 28 ust. 4\autocite{dora2022}), również ustawa KSC (Art. 8\autocite{uksc2026}) wymaga opracowania exit planu.  Całość procesu musi być zaplanowana tak, aby klient instytucji w ogóle nie odczuł zmiany dostawcy.

Ważma jest forma przygotowania dokumentacji. DORA, w art 28 ust. 4\autocite{dora2022} wymaga oddzielnego dokumentu, natomiast KSC wymaga realizuje cele exit planu poprzez klauzule umowne z dostawcą.

W metodoligii ORION, exit plan jest wymagany dla podmiotów określonych w KSC oraz DORA. Podobnie jak plan ciągłości działania,  brak lub nieprawidłowo zdefiniowany exit plan, gdu jest obowiązkowy, wskazuje binarnie na uchybienia w działaniu organizacji.

\begin{table}[H]
\centering
\caption{Mnożniki wagowe exit plan wg typu organizacji}
\label{tab:exitplan-wagi}
\renewcommand{\arraystretch}{1.4}
\begin{tabular}{p{4.5cm}p{1.5cm}p{2.5cm}p{4.5cm}}
\toprule
\textbf{Typ organizacji} & \textbf{Symbol} & \textbf{$w_{reg}$} & Obowiązkowe\\
\midrule
\texttt{FINANSOWY} & $D_0$ & $1{,}0$   & tak\\
\midrule
\texttt{KLUCZOWY} & $D_1$ & $1{,}0$   & tak\\
\midrule
\texttt{WAŻNY} & $D_2$ & $1{,}0$   & nie\\
\midrule
\texttt{POZOSTAŁE} & $D_3$ & $0{,}5$ & nie \\
\bottomrule

\end{tabular}
\end{table}




% \section{Analiza}

% Nąstepnie kazde pytanie zadane do kazdego systemu (artefaktu):

% Wymóg regulacyjny - czyli w czy wartość wynika z ustawy i~mysi być spełniona, czy też jest dobrą praktyką, więc opcjonalnie.
% Punkt widzenia artefaktu - prawne, biznesowe, techniczne w zakresie oprogramowania, techniczne w zakresie wdrożenia. 
% Dotyczy to samego pytania, ktore pomoze na wejsciu ocenic kto powien na pytanie odpowiedziec.

% i~odpowiedzi przypisywane są do binarnych wyników.

% Dowód - czy przeprowadzony zostal dowod. Nie jest on wymagany, ale jego brak obniza wartosc. Jest to wartosc wylisotwana: 0, 0.5 (jest deklaracja), 1 - potwierdzeni dowodem, ktora okresla jakosc dowodu. Waga dowodu - jesli inna niz 1 to rosnie ryzyka mocniej, zaproponuj jak.

% Waga ryzyka -  okreslnie specjalnej wagi ryzyka, domyslnie 1. Globalna waga dla pytania.

% Następnie kazda odpowiedz grupowana jest w grupy i~powstaje procentowy wynik na podstawie wktora wartosci uwzglednaijac wage. 



% Wnioskowanie:

% Ryzyko

% Poziom zgodności

% Luka zgodności

% Każde pojęcie:

% definicja słowna

% definicja formalna

% relacje z innymi

% \textbf{zrobie zgodnosc binarną}



\newcommand{\generujKryteria}[2]{%
\vspace{0.5cm}\\
Pytania w sekcji:\\\\
  \csvreader[
    separator=semicolon,
    no head
  ]{#1}{}{%%
    \ifcsvstrcmp{\csvcoli}{#2}{%
      \noindent
      Obszar: \textbf{\csvcolii} \\
      Podstawa prawna: \textbf{\csvcoliii} \\
      Typ pytania: \textbf{\csvcolviii}, blokujące: \textbf{\csvcolix}

      \begin{itemize}
        \item \textbf{Perspektywa zgodności:} \csvcoliv
        \item \textbf{Perspektywa użytkownika:} \csvcolvi
        \item \textbf{Perspektywa techniczna:} \csvcolv
      \end{itemize}
    }{}%
  }
}

% ============================================================
\chapter{Przeprowadzenie audytu ORION}
\label{chap:audyt}
% ============================================================

Niniejszy rozdział opisuje pełny algorytm przeprowadzenia audytu zgodności systemów informatycznych w metodologii ORION.
Audyt jest procesem iteracyjnym, co oznacza, że po zakończeniu jednego etapu audytor przechodzi do kolejnego, składającym się z pięciu faz następujących po sobie w ściśle określonej kolejności. 
Każda faza produkuje dane wejściowe dla fazy następnej, a wynik całości, czyli ocena zgodności podmiotu z aktualnym stanem regulacyjnym z punktu widzenia ICT,  jest wynikiem wszystkich odpowiedzi na każdym etapie.
Schemat ogólny procesu przedstawia rysunek~\ref{fig:fazy-audytu}.

\begin{figure}[H]
\centering
\begin{tikzpicture}[
  node distance=1.6cm,
  box/.style={rectangle, rounded corners=4pt,
              draw=black, fill=gray!10,
              minimum width=10cm, minimum height=1cm,
              text centered, font=\small},
  arrow/.style={-Stealth, thick}
]
\node[box] (f1) {\textbf{Faza 1} \quad Klasyfikacja podmiotu};
\node[box, below of=f1] (f2)
  {\textbf{Faza 2} \quad Inwentaryzacja i~klasyfikacja obiektów};
\node[box, below of=f2] (f3)
  {\textbf{Faza 3} \quad Kwestionariusz szczegółowy i~zbieranie dowodów};
\node[box, below of=f3] (f4)
  {\textbf{Faza 4} \quad Obliczenie scoringu i~propagacja ryzyka};
\node[box, below of=f4] (f5)
  {\textbf{Faza 5} \quad Raport końcowy i~rekomendacje};
\draw[arrow] (f1) -- (f2);
\draw[arrow] (f2) -- (f3);
\draw[arrow] (f3) -- (f4);
\draw[arrow] (f4) -- (f5);
\end{tikzpicture}
\caption{Pięć faz audytu ORION}
\label{fig:fazy-audytu}
\end{figure}

% ============================================================
\section{Faza 1: Klasyfikacja podmiotu}
\label{sec:faza1}
% ============================================================

Audyt rozpoczyna się od ustalenia tożsamości regulacyjnej audytowanego podmiotu. Wynik tej fazy determinuje, które zestawy pytań zostaną zaprezentowane w dalszych etapach, jakie wymogi mają charakter obligatoryjny, a jakie fakultatywne oraz jakie progi klasyfikacji dostawców kluczowych będą stosowane. 
Faza pierwsza składa się z dwóch kroków wykonywanych wyłącznie raz, przed przystąpieniem do analizy jakiegokolwiek obiektu i~dotyczy globalnie całej organizacji.

\subsection{Krok 1.1: Ustalenie reżimu regulacyjnego}
\label{sec:krok11}

Dla każdego kryterium poniżej odpowiedz \texttt{TAK} lub \texttt{NIE}. Klasyfikacja wyznaczana jest na podstawie konsensusu odpowiedzi z~co najmniej dwóch perspektyw.

\bigskip

\noindent\textbf{Kryterium R1: Podleganie reżimowi regulacyjnemu}

\begin{itemize}
  \item \textbf{Perspektywa zgodności} (DPO / Radca prawny): Czy podmiot prowadzi działalność wymienioną w~załączniku~I lub~II dyrektywy NIS2 albo w~art.~2 rozporządzenia DORA i~przekracza progi wielkości określone przez te akty?
  \item \textbf{Perspektywa użytkownika} (Zarząd / Właściciel): Czy organizacja świadczy usługi finansowe, prowadzi infrastrukturę krytyczną lub przetwarza dane na zlecenie podmiotów regulowanych?
  \item \textbf{Perspektywa techniczna} (CTO / Architekt): Czy systemy IT organizacji były kiedykolwiek objęte obowiązkowym audytem regulacyjnym lub wymogiem zgłaszania incydentów do organu nadzorczego?
\end{itemize}

Na podstawie konsensusu odpowiedzi podmiot otrzymuje
jedną z następujących klasyfikacji regulacyjnych,
które określają obowiązkowość kluczowych elementów audytu:

\begin{table}[ht] % Zmieniono z [H] na [ht] dla lepszej kompatybilności
\centering
\caption{Klasyfikacja regulacyjna podmiotu i~jej wpływ na zakres audytu}
\label{tab:klasyfikacja-regulacyjna}
\renewcommand{\arraystretch}{1.5}
% Usunięto resizebox na rzecz precyzyjnych szerokości p{}
\begin{tabular}{p{2.5cm}p{3cm}cccc}
\toprule
\textbf{Klasyfikacja} & \textbf{Podstawa} & \textbf{BCP} & \textbf{Exit plan} & \textbf{TIA} & $w_{reg}$ \\ \midrule
\texttt{FINANSOWY} & DORA Art.~2 & Tak & Tak & Tak & $1{,}5$ \\ \midrule
\texttt{KLUCZOWY} & NIS2 / KSC Art.~5 & Tak & Tak & Tak & $1{,}3$ \\ \midrule
\texttt{WAŻNY} & NIS2 / KSC Art.~6 & Tak & Nie & Tak & $1{,}1$ \\ \midrule
\texttt{POZOSTAŁE} & RODO / ogólne & Nie & Nie & Warunkowy & $1{,}0$ \\ \bottomrule
\end{tabular}
\end{table}

% ============================================================
\section{Faza 2: Inwentaryzacja i~klasyfikacja obiektów}
\label{sec:faza2}
% ============================================================

Celem fazy drugiej jest zbudowanie grafu ICT podmiotu,  pełnego rejestru wszystkich obiektów podlegających audytowi
wraz z ich atrybutami klasyfikacyjnymi. Graf ICT jest strukturą dynamiczną: jego wierzchołki są odkrywane
iteracyjnie, a mechanizm triggera TIA opisany w sekcji~\ref{sec:trigger-tia} może rozszerzać graf o nowe obiekty w trakcie fazy trzeciej. Faza druga stanowi zatem pierwszą, wstępną wersję grafu, która ulega uzupełnieniu w toku audytu szczegółowego.

\subsection{Krok 2.1: Inwentaryzacja systemów}
\label{sec:krok21}

Audytor przeprowadza inwentaryzację wszystkich systemów informatycznych eksploatowanych przez podmiot. Jako punkt wyjścia służą: rejestr systemów prowadzony przez dział IT, wykaz aplikacji z licencji oprogramowania (pobranych z systemów księgowych oraz archiwów licencji), rejestr umów z dostawcami ICT, wyniki poprzednich audytów lub testów penetracyjnych. Każdy zidentyfikowany system jest rejestrowany jako obiekt w grafie ICT i~przypisany do właściciela biznesowego.
Należy również przeskanować systemy informatyczne w celu zaintwentaryzowania zainstalowanego oprogramowania, na rynku dostępne są do tego narzędzia typu \textit{open source}.

Kluczowe w inwentaryzacji jest podejście wieloperspektywiczne: inwentaryzacja przeprowadzana wyłącznie przez dział IT często pomija systemy zakupione i~utrzymywane bezpośrednio przez jednostki biznesowe (\textit{shadow IT}), podczas gdy wywiad z właścicielami biznesowymi i~prawnikami ujawnia dodatkowe systemy i~usługi SaaS nieprzechodzące przez standardowy proces zakupów IT.
Wynikiem operacji identyfikacji będzie lista obiektów typu \texttt{SYSTEM\_INFORMATYCZNY}, która będzie dołączona do grafu ICT.

\subsection{Krok 2.2: Identyfikacja bezpośrednich dostawców}
\label{sec:krok22}

Dla każdego zidentyfikowanego systemu audytor identyfikuje bezpośrednich dostawców ICT: podmioty zewnętrzne świadczące
usługi hostingu, utrzymania, wsparcia lub integracji.
Każdy dostawca jest rejestrowany jako odrębny obiekt typu \texttt{DOSTAWCA} w grafie ICT i~połączony krawędzią skierowaną z systemem, który obsługuje. Na tym etapie nie są jeszcze identyfikowani poddostawcy, ich identyfikacja następuje w fazie trzeciej przez mechanizm triggera TIA.

\subsection{Krok 2.3: Przypisanie atrybutów klasyfikacyjnych}
\label{sec:krok23}

Dla każdego obiektu zidentyfikowanego w kroku \ref{sec:krok21} audytor przeprowadza klasyfikację wstępną, ustalając wartości atrybutów:
\textit{dostępność}(v),
\textit{krytyczność}(v),
\textit{utrzymanie\_chmura}(v),
\textit{typ\_danych\_wrażliwych}(v), 


Każdy atrybut jest wyznaczany metodą cross-validation opisaną w sekcji~\ref{sec:crossvalidation}. 


\subsection{Ankieta klasyfikacji atrybuty dostępności }

Poniższe kryteria służą do określenia, z jaka jest dostępność systemu z punktu widzenia prywatności rozwiązania.  Dla każdego kryterium odpowiedz \texttt{TAK} lub \texttt{NIE}. Klasyfikacja jest wyznaczana na ilości odpowiedzi zgodnych z odpowiedzią \textbf{oczekiwaną} jako poprawna.

\bigskip

\textbf{Kryterium D1: Dostęp spoza sieci organizacji}

\begin{itemize}
\item \textbf{Perspektywa zgodności:} Czy umowa lub regulamin systemu przewiduje dostęp dla podmiotów spoza struktury prawnej organizacji?
\item \textbf{Perspektywa użytkownika:} Czy z systemu korzystają osoby lub firmy spoza naszej organizacji (klienci, partnerzy, dostawcy)?
\item \textbf{Perspektywa techniczna:} Czy system jest dostępny spoza sieci LAN/VPN organizacji (publiczny IP, internet)?
\end{itemize}

% ============================================================

\textbf{Kryterium D2: Użytkownicy zewnętrzni}

\begin{itemize}
\item \textbf{Perspektywa zgodności:} Czy system wymaga odrębnych umów powierzenia danych lub regulaminów dla jego użytkowników?
\item \textbf{Perspektywa użytkownika:} Czy system służy do komunikacji lub wymiany informacji z klientami bądź kontrahentami?
\item \textbf{Perspektywa techniczna:} Czy system posiada konta lub sesje dla użytkowników spoza domeny organizacji?
\end{itemize}

% ============================================================

\textbf{Kryterium D3: Integracje zewnętrzne}

\begin{itemize}
\item \textbf{Perspektywa zgodności:} Czy system przesyła dane do podmiotów trzecich na podstawie umów lub przepisów prawa?
\item \textbf{Perspektywa użytkownika:} Czy system automatycznie wymienia dane z zewnętrznymi dostawcami, bankami lub platformami?
\item \textbf{Perspektywa techniczna:} Czy system wywołuje zewnętrzne API lub odbiera dane z systemów spoza infrastruktury organizacji?
\end{itemize}

% ============================================================

\textbf{Kryterium D4: Publiczny adres lub domena}

\begin{itemize}
\item \textbf{Perspektywa zgodności:} Czy system jest zarejestrowany lub dostępny pod adresem publicznie widocznym w internecie?
\item \textbf{Perspektywa użytkownika:} Czy system ma adres www lub można do niego dotrzeć bez łączenia się z siecią wewnętrzną firmy?
\item \textbf{Perspektywa techniczna:} Czy system posiada publiczny rekord DNS lub publiczny adres IP?
\end{itemize}


\bigskip
\textbf{Wynik końcowy:}

\begin{table}[ht]
\begin{tabular}{p{10cm} p{4cm}}
\toprule
Suma TAK & Klasyfikacja \\
\midrule
Zero odpowiedzi twierdzących     & System \textbf{prywatny}    \\
\midrule
Minimum jedna twierdząca     & System \textbf{publiczny}     \\
\bottomrule
\end{tabular}
\end{table}



% ============================================================
%  Ankieta system krytyczny, ważny czy zwykły?
% ============================================================

\subsection{Ankieta klasyfikacja atrybutu poziomu krytyczności}
Poniższe kryteria służą do określenia, z jakim poziomem krytyczności operuje system. Dla każdego kryterium odpowiedz \texttt{TAK} lub \texttt{NIE}. Klasyfikacja jest wyznaczana na ilości odpowiedzi twierdzących.

\bigskip
\noindent\textbf{Kryterium K1: Wpływ finansowy i~operacyjny
niedostępności systemu}
\begin{itemize}
\item \textbf{Perspektywa zgodności:} Czy niedostępność
systemu może powodować roszczenia kontraktowe lub
kary umowne?
\item \textbf{Perspektywa użytkownika:} Czy kilkugodzinny
przestój zatrzymuje sprzedaż, obsługę klienta lub
kluczowy proces biznesowy?
\item \textbf{Perspektywa techniczna:} Czy system ma
zdefiniowane RTO poniżej 4~godzin lub SLA z karami
za niedostępność?
\end{itemize}

\noindent\textbf{Kryterium K2: Przetwarzanie danych osobowych
lub wrażliwych}
\begin{itemize}
\item \textbf{Perspektywa zgodności:} Czy system przetwarza
dane osobowe lub szczególne kategorie danych
w rozumieniu art.~9 RODO?
\item \textbf{Perspektywa użytkownika:} Czy w systemie
przechowywane są dane klientów, pracowników lub inne
informacje poufne?
\item \textbf{Perspektywa techniczna:} Czy baza danych
zawiera pola z danymi PII
(np.\ imię, PESEL, adres, numer konta)?
\end{itemize}

\noindent\textbf{Kryterium K3: Podleganie regulacjom zewnętrznym}
\begin{itemize}
\item \textbf{Perspektywa zgodności:} Czy system podlega
regulacjom sektorowym
(np.\ DORA, NIS2, ustawa KSC)?
\item \textbf{Perspektywa użytkownika:} Czy system był
przedmiotem pytań audytorów lub regulatorów
podczas kontroli?
\item \textbf{Perspektywa techniczna:} Czy system znajduje
się w wykazie systemów objętych audytem zgodności
lub testami penetracyjnymi?
\end{itemize}

\noindent\textbf{Kryterium K4: Zależności innych systemów
produkcyjnych}
\begin{itemize}
\item \textbf{Perspektywa zgodności:} Czy awaria systemu
może uniemożliwić realizację zobowiązań wobec klientów
lub partnerów?
\item \textbf{Perspektywa użytkownika:} Czy inne działy
lub systemy firmy nie mogą działać, gdy system jest
niedostępny?
\item \textbf{Perspektywa techniczna:} Czy inne systemy
produkcyjne mają bezpośrednie zależności od tego systemu?
\end{itemize}

\noindent\textbf{Kryterium K5: Ryzyko wycieku lub utraty danych}
\begin{itemize}
\item \textbf{Perspektywa zgodności:} Czy wyciek danych
z systemu wymagałby zgłoszenia incydentu do UODO
lub regulatora?
\item \textbf{Perspektywa użytkownika:} Czy ujawnienie
danych mogłoby zaszkodzić reputacji firmy lub spowodować
utratę klientów?
\item \textbf{Perspektywa techniczna:} Czy system przechowuje
dane bez szyfrowania lub bez możliwości pseudonimizacji?
\end{itemize}

\noindent\textbf{Kryterium K6: Obsługa procesów podstawowej
działalności (core business)}
\begin{itemize}
\item \textbf{Perspektywa zgodności:} Czy system jest
wskazany w dokumentacji BCM/BCP jako niezbędny
do kontynuacji działalności?
\item \textbf{Perspektywa użytkownika:} Czy system obsługuje
główny produkt lub usługę firmy?
\item \textbf{Perspektywa techniczna:} Czy system jest
sklasyfikowany jako produkcyjny tier-1
w architekturze IT?
\end{itemize}

W celu uniknięcia zjawiska sztucznego zawyżania krytyczności (gdzie niemal każdy system przetwarzający dane osobowe stawał się krytyczny), w metodologii ORION zastosowano zasadę dominacji operacyjnej. System nie może zostać uznany za krytyczny tylko na podstawie faktu, że podlega regulacjom (K3) lub przetwarza dane (K2), o ile jego niedostępność nie paraliżuje podstawowej działalności organizacji (K1, K6).

\begin{table}[H]
\centering
\caption{Wynik klasyfikacji krytyczności na podstawie
         sumy odpowiedzi twierdzących}
\renewcommand{\arraystretch}{1.3}
\begin{tabular}{p{3cm} p{4cm} p{5cm}}
\toprule
\textbf{Suma TAK} & \textbf{Klasyfikacja} & \textbf{Warunek dodatkowy} \\
\midrule
0--2  & System \textbf{standardowy} & --- \\
\midrule
3--4  & System \textbf{ważny}       & --- \\
\midrule
5--6  & System \textbf{krytyczny}   & Wymagane K1\,=\,TAK lub K6\,=\,TAK
                                      (zasada dominacji operacyjnej) \\
\bottomrule
\end{tabular}
\end{table}



% ============================================================
% ANKIETA 2: System w chmurze
\subsection{Ankieta klasyfikacji atrybutu utrzymania w chmurze}
\label{sec:ankieta-chmura}
Poniższe kryteria służą do określenia, czy system jest rozwiązaniem on-premise czy chmurowym . Dla każdego kryterium odpowiedz \texttt{TAK} lub \texttt{NIE}. Klasyfikacja jest wyznaczana na ilości odpowiedzi twierdzących.


\bigskip

\textbf{Kryterium L1: Model wdrożenia infrastruktury}

\begin{itemize}
\item \textbf{Perspektywa zgodności:} Czy umowa lub dokumentacja systemu wskazuje na dostawcę chmury publicznej, prywatnej lub hybrydowej (np. AWS, Azure, GCP, OVH)?
\item \textbf{Perspektywa użytkownika:} Czy system działa w internecie lub w środowisku zarządzanym przez zewnętrznego dostawcę, a nie na serwerach w siedzibie firmy?
\item \textbf{Perspektywa techniczna:} Czy serwery lub maszyny wirtualne systemu są uruchomione poza fizycznym data center organizacji?
\end{itemize}

% ============================================================

\textbf{Kryterium L2: Odpowiedzialność operacyjna}

\begin{itemize}
\item \textbf{Perspektywa zgodności:} Czy odpowiedzialność za dostępność i~utrzymanie infrastruktury systemu spoczywa na zewnętrznym dostawcy zgodnie z umową SLA?
\item \textbf{Perspektywa użytkownika:} Czy w przypadku awarii serwera zgłaszamy problem do zewnętrznej firmy, a nie do naszego działu IT?
\item \textbf{Perspektywa techniczna:} Czy organizacja nie ma bezpośredniego dostępu administracyjnego do warstwy sprzętowej lub hiperwizora systemu?
\end{itemize}

% ============================================================

\textbf{Kryterium L3: Umowa outsourcingowa lub chmurowa}

\begin{itemize}
\item \textbf{Perspektywa zgodności:} Czy z dostawcą systemu lub infrastruktury podpisano umowę outsourcingu IT, umowę SaaS/PaaS/IaaS lub umowę powierzenia przetwarzania danych?
\item \textbf{Perspektywa użytkownika:} Czy dostęp do systemu lub jego funkcji jest uzależniony od aktywnej subskrypcji lub kontraktu z zewnętrznym dostawcą?
\item \textbf{Perspektywa techniczna:} Czy system jest wdrożony w modelu usługowym (SaaS, PaaS lub IaaS), w którym zasoby obliczeniowe są rozliczane na podstawie zużycia lub abonamentu?
\end{itemize}

% ============================================================

\textbf{Kryterium L4: Lokalizacja i~kontrola danych}

\begin{itemize}
\item \textbf{Perspektywa zgodności:} Czy dane przetwarzane przez system są przechowywane w lokalizacji określonej przez dostawcę zewnętrznego oraz czy przeprowadzono pełne TIA?
\item \textbf{Perspektywa użytkownika:} Czy jako użytkownik systemu masz pewność, że fizycznie dane przechowywane są na serwerach kraju EOG?
\item \textbf{Perspektywa techniczna:} Czy kopie zapasowe lub dane produkcyjne systemu są składowane na zasobach poza kontrolą administracyjną organizacji?
\end{itemize}

% ============================================================

\textbf{Kryterium L5: Jurysdykcja podmiotu hostującego}

\begin{itemize}
  \item \textbf{Perspektywa zgodności:} Czy dostawca przeprowadził pełne TIA
        wraz z poddostawcami, z którego wynika, że wszystkie podmioty
        zaangażowane w przetwarzanie danych są własnością kapitałową
        przedsiębiorstw z EOG lub nie podlegają jurysdykcji państw trzecich
        (np.\ US CLOUD Act, FISA Section 702)?

  \item \textbf{Perspektywa użytkownika:} Czy użytkownik może potwierdzić,
        że wszystkie usługi i~okna logowania systemu należą do podmiotów
        zarejestrowanych i~działających wyłącznie na podstawie prawa
        europejskiego?

  \item \textbf{Perspektywa techniczna:} Czy certyfikaty SSL/TLS, adresy IP,
        rekordy DNS oraz dane WHOIS infrastruktury systemu wskazują wyłącznie
        na podmioty zarejestrowane w EOG, a żaden z komponentów technicznych
        (CDN, load balancer, usługa uwierzytelniania) nie pochodzi od dostawcy
        spoza EOG lub podlegającego obcemu prawu dostępu do danych?
\end{itemize}

\textbf{Kryterium L6: Wsparcie techniczne}

\begin{itemize}
  \item \textbf{Perspektywa zgodności:} Czy umowa z dostawcą systemu
        jednoznacznie określa, że wsparcie techniczne (helpdesk, administracja,
        dostęp serwisowy) jest świadczone wyłącznie przez osoby lub podmioty
        mające siedzibę w EOG?

\item \textbf{Perspektywa użytkownika:} Czy użytkownik może potwierdzić,
      że wszystkie punkty kontaktowe wsparcia technicznego (w tym adresy
      e-mail, domeny helpdesk, numery telefonów oraz profile agentów
      w systemie zgłoszeń) wskazują na podmioty zarejestrowane
      w EOG (europejskie domeny, kierunkowe numery telefonów,
      europejskie dane rejestrowe dostawcy)?

  \item \textbf{Perspektywa techniczna:} Czy narzędzia zdalnego dostępu
        serwisowego (np.\ VPN, RDP, sesje administracyjne) są skonfigurowane
        w sposób uniemożliwiający połączenie spoza EOG, a logi dostępu
        potwierdzają, że żadna sesja wsparcia technicznego nie była
        inicjowana z adresów IP spoza Europejskiego Obszaru Gospodarczego?
\end{itemize}
\bigskip
\textbf{Wynik końcowy:}

\begin{table}[ht]
\begin{tabular}{p{10cm} p{4cm}}
\toprule
Suma TAK & Klasyfikacja \\
\midrule
Zero odpowiedzi twierdzących  & System \textbf{on-premise}   \\
\midrule
Minimum jedna odpowiedź twierdząca oraz minimum jedna odpowiedź negatywna na pytania L5, L6      & System \textbf{w chmurze krajów trzecich}    \\
\midrule
Minimum trzy odpowiedi twierdząca w tym na pytania L5, L6     & System \textbf{w chmurze EOG}    \\
\bottomrule
\end{tabular}
\end{table}

% ====================================
\subsection{Ankieta klasyfikacji typu danych}

Poniższe kryteria służą do określenia, z jakim typem danych operuje system. Dla każdego kryterium odpowiedz \texttt{TAK} lub \texttt{NIE}. Klasyfikacja jest wyznaczana na podstawie najwyższej kategorii, dla której udzielono co najmniej jednej odpowiedzi twierdzącej.

\bigskip
\noindent\textbf{Kryterium W: Dane szczególnych kategorii (wrażliwe)}
\begin{itemize}
    \item \textbf{Perspektywa zgodności:} Czy system przetwarza dane
    szczególnych kategorii w rozumieniu art.~9 RODO
    (np.\ dane dotyczące zdrowia, wyroków skazujących, przekonań
    religijnych, orientacji seksualnej, przynależności związkowej)?
    \item \textbf{Perspektywa użytkownika:} Czy system zawiera
    dokumentację medyczną, orzeczenia, dane HR
    o niepełnosprawności lub inne informacje uznawane
    za szczególnie chronione?
    \item \textbf{Perspektywa techniczna:} Czy baza danych zawiera
    pola klasyfikowane jako dane wrażliwe (np.\ diagnoza,
    orzeczenie, wyznanie, dane biometryczne)?
\end{itemize}

\bigskip
\noindent\textbf{Kryterium O: Dane osobowe (zwykłe)}
\begin{itemize}
    \item \textbf{Perspektywa zgodności:} Czy system przetwarza
    dane osobowe w rozumieniu art.~4 pkt~1 RODO,
    pozwalające na identyfikację osoby fizycznej
    (np.\ imię i~nazwisko, PESEL, adres e-mail, numer konta)?
    \item \textbf{Perspektywa użytkownika:} Czy w systemie
    przechowywane są dane klientów, pracowników
    lub kontrahentów w jakiejkolwiek formie?
    \item \textbf{Perspektywa techniczna:} Czy baza danych zawiera
    pola, takie jak imię, adres, identyfikator użytkownika, adres IP lub dane logowania?
\end{itemize}

\bigskip
\noindent\textbf{Kryterium T: Tajemnice prawnie chronione}
\begin{itemize}
    \item \textbf{Perspektywa zgodności:} Czy system przetwarza
    informacje objęte tajemnicą zawodową, handlową,
    bankową, skarbową, adwokacką lub inną tajemnicą
    chronioną przepisami prawa?
    \item \textbf{Perspektywa użytkownika:} Czy ujawnienie danych
    z systemu mogłoby naruszyć obowiązki zachowania
    poufności wynikające z umów lub przepisów prawa
    (np.\ NDA, tajemnica przedsiębiorstwa)?
    \item \textbf{Perspektywa techniczna:} Czy system posiada
    oznaczenie klauzulą poufności lub jest objęty
    restrykcjami dostępu wynikającymi z wymogów
    regulacyjnych lub umownych?
\end{itemize}

\begin{table}[H]
\centering
\caption{Wynik klasyfikacji typu danych na podstawie
         udzielonych odpowiedzi twierdzących}
\renewcommand{\arraystretch}{1.3}
\begin{tabular}{p{5cm} p{4cm} p{4cm}}
\toprule
\textbf{Warunek} & \textbf{Typ danych} & \textbf{Reżim ochrony} \\
\midrule
Co najmniej jedno TAK w~kryterium W
    & \textbf{Wrażliwe}
    & art.~9 RODO, DPIA wymagane \\
\midrule    
Co najmniej jedno TAK w~kryterium O
(brak TAK w~W)
    & \textbf{Osobowe}
    & RODO, zgłoszenie do UODO \\
\midrule    
Co najmniej jedno TAK w~kryterium T
(brak TAK w~W i~O)
    & \textbf{Tajemnica}
    & przepisy sektorowe, NDA \\
\midrule    
Wyłącznie NIE we wszystkich kryteriach
    & \textbf{Brak}
    & brak szczególnych wymogów \\
\bottomrule
\end{tabular}
\end{table}

\noindent\textbf{Uwaga:} Klasyfikacja ma charakter kaskadowy ---
wystarczy jedno TAK w~kryterium wyższej kategorii, aby
przypisać wyższy typ danych. Jeśli system spełnia warunki
kilku kryteriów jednocześnie, stosuje się typ danych
o~najwyższym ryzyku.

% ============================================================
\section{Faza 3: Kwestionariusz szczegółowy i~zbieranie dowodów}
\label{sec:faza3}
% ============================================================

Na podstawie atrybutów obiektów i~wartości im przypisanych w fazie drugiej system ORION generuje dla każdego
obiektu indywidualny kwestionariusz szczegółowy, zestaw pytań audytowych pogrupowanych w sekcje tematyczne,
aktywowane warunkowo w zależności od wartości atrybutów. 
Żaden obiekt nie jest pytany o aspekty, które są dla niego nieistotne z powodów technicznych lub regulacyjnych.



\subsection{Warunkowe aktywowanie sekcji pytań}
\label{sec:aktywowanie-sekcji}

Zestaw aktywnych sekcji pytań dla obiektu $v$ jest wyznaczany przez funkcję:

\begin{equation}
\mathcal{S}(v) =
\bigcup_{\substack{\tau \in \mathcal{T} \\ \mathrm{type}(v)=\tau}} \mathcal{S}_{\tau}
\;\cup\;
\bigcup_{\substack{k \in K \\ \Phi_k(\mathcal{A}(v))}} \mathcal{S}_k
\label{eq:sekcje}
\end{equation}

\noindent gdzie:
\begin{itemize}
  \item $\mathcal{T} = \{\texttt{PODMIOT},\, \texttt{SYSTEM\_INFORMATYCZNY},\, \texttt{DOSTAWCA}\}$
        jest zbiorem typów obiektów w grafie ICT,
  \item $\mathcal{S}_{\tau}$ jest zbiorem sekcji aktywowanych wyłącznie
        dla obiektów danego typu $\tau$ (np.\ sekcje
        \texttt{ORGANIZACJA}, \texttt{BEZPIECZENSTWO\_TECHNICZNE}
        i~\texttt{BEZPIECZENSTWO\_FIZYCZNE} są aktywowane
        gdy $\mathrm{type}(v) = \texttt{PODMIOT}$,
        a \texttt{APLIKACJA} gdy $\mathrm{type}(v) = \texttt{SYSTEM\_INFORMATYCZNY}$),
  \item $K$ jest zbiorem indeksów wszystkich sekcji aktywowanych warunkowo,
  \item $\mathcal{S}_k$ jest $k$-tą sekcją aktywowaną warunkowo
        na podstawie atrybutów obiektu,
  \item $\Phi_k\bigl(\mathcal{A}(v)\bigr)$ jest warunkiem logicznym
        aktywacji sekcji $\mathcal{S}_k$ wyznaczanym na podstawie
        wektora atrybutów $\mathcal{A}(v)$ obiektu $v$
        (np.\ klasyfikacji utrzymania w chmurze, typu danych, krytyczności).
\end{itemize}

Pierwszy składnik sumy pokrywa sekcje aktywowane przez \emph{typ} obiektu i~jest niezależny od jego atrybutów. Drugi składnik pokrywa sekcje aktywowane przez \emph{atrybuty} obiektu i~jest niezależny od jego typu. Oba zbiory mogą być rozłączne lub przecinać się, przy czym wynikowy zbiór aktywnych sekcji $\mathcal{S}(v)$ jest ich sumą teoriomnogościową.

\noindent\textbf{Przykład.}
Dla obiektu $v$ z~$\mathrm{type}(v) = \texttt{SYSTEM\_INFORMATYCZNY}$
i~atrybutem $\textit{utrzymanie\_chmura}(v) = \texttt{CHMURA\_K3}$
pierwszy składnik wzoru \eqref{eq:sekcje} daje $\mathcal{S}_{\texttt{SYSTEM\_INFORMATYCZNY}}$,
a drugi (ponieważ warunek $\Phi_{\texttt{CHMURA}}$ jest spełniony) dodaje $\mathcal{S}_{\texttt{CHMURA}}$.
Wynikowy zbiór aktywnych sekcji to
$\mathcal{S}(v) = \mathcal{S}_{\texttt{SYSTEM\_INFORMATYCZNY}} \cup \mathcal{S}_{\texttt{CHMURA}}$.

Pytania zostały podzielone na sekcje zgodnie z klasyfikacją przyjętą w rodziale \ref{perspektywa-pytań} \textit{Perspektywa}. Każda pozycja odnosi się do konkretnego obszaru zagadnień, zawiera podstawę prawną, definicję istotności oraz oznaczenie (flagę binarną) pytania blokującego.


\subsection{Sekcja organizacja}
\noindent Aktywacja: \textit{sekcja aktywna jednorazowo dla całego audytowanego podmiotu.}\\
Pytania w tej sekcji służą do oceny dojrzałości zarządczej organizacji w zakresie cyberbezpieczeństwa i~odporności cyfrowej. Sprawdzają, czy na poziomie strategicznym podmiot jest świadomy ryzyk ICT, czyli czy zarząd formalnie zatwierdził strategię cyberbezpieczeństwa i~ponosi za nią odpowiedzialność prawną, czy przeszedł wymagane szkolenia oraz czy wdrożono ramy zarządzania ryzykiem ICT z rejestrem ryzyk, właścicielami i~planami. Ocenie podlega również gotowość operacyjna: czy organizacja potrafi klasyfikować incydenty, zgłaszać poważne zdarzenia do organów krajowych w ustawowych terminach oraz informować klientów o incydentach według gotowych procedur.
\generujKryteria{pytania.csv}{ORGANIZACJA}


\subsection{Sekcja dane poufne: osobowe, wrażliwe oraz tajemnice przedsiębiorstwa}
\noindent Aktywacja: \textit{sekcja aktywna jednorazowo dla całego audytowanego podmiotu.}\\
Pytania w tej sekcji służą do oceny, czy organizacja prawidłowo zarządza procesem gromadzenia oraz przetwarzania danych poufnych. Weryfikowana jest również gotowość do zgłoszenia naruszenia ochrony danych osobowych, które chroni prawa osób fizycznych i~wymaga odrębnej procedury, odrębnej ścieżki eskalacji i~odrębnych formularzy. Organizacja nieposiadająca oddzielnej procedury zgłoszeń do UODO nie wykazuje gotowości do wywiązania się z fundamentalnego obowiązku w zakresie ochrony danych osobowych.
\generujKryteria{pytania.csv}{DANE_WRAZLIWE}

\subsection{Sekcja bezpieczeństwo techniczne}
\noindent Aktywacja: \textit{sekcja aktywna jednorazowo dla całego audytowanego podmiotu.}\\
Pytania w tej sekcji służą do oceny skuteczności technicznych środków ochrony systemów i~infrastruktury ICT organizacji przed zagrożeniami cybernetycznymi. Sprawdzane jest, czy organizacja wdrożyła mechanizmy ciągłego monitorowania i~wykrywania anomalii oraz czy przeprowadza regularne testy bezpieczeństwa i~testy penetracyjne z udokumentowanymi wynikami i~planem naprawy zidentyfikowanych podatności. Ocenie podlega stosowanie zasady minimalnych uprawnień (least privilege) w odniesieniu zarówno do kont użytkowników, jak i~kont serwisowych aplikacji dostawców, z kwartalnym przeglądem uprawnień i~automatycznym alertowaniem przy próbach eskalacji.
\generujKryteria{pytania.csv}{BEZPIECZENSTWO_TECH}

\subsection{Sekcja bezpieczeństwo fizyczne}
\noindent Aktywacja: \textit{sekcja aktywna jednorazowo dla całego audytowanego podmiotu.}\\
Pytania w tej sekcji służą do oceny środków ochrony fizycznej stosowanych w obiektach, w których zlokalizowana jest infrastruktura ICT organizacji lub jej kluczowych dostawców. Sprawdzane jest, czy dostęp do stref krytycznych (serwerowni i~pomieszczeń technicznych) jest kontrolowany przez elektroniczne systemy uwierzytelnienia (karty dostępu, biometria) z automatycznym logowaniem każdego wejścia, eliminującym możliwość nieautoryzowanego wejścia bez śladu w rejestrze. Ocenie podlega podział obiektu na strefy bezpieczeństwa o różnych poziomach dostępu, monitoring wizyjny CCTV z retencją nagrań zgodną z RODO, redundancja zasilania testowana pod rzeczywistym obciążeniem (UPS i~agregat prądotwórczy) oraz procedury dotyczące urządzeń przenośnych: szyfrowanie całego dysku, rejestr wynoszenia i~zdalna możliwość wyczyszczenia. Weryfikowane jest ponadto, czy nośniki danych przy utylizacji są niszczone fizycznie lub nadpisywane certyfikowaną metodą, a serwisanci zewnętrzni w strefach krytycznych są przez cały czas eskortowani i~ich działania są logowane. Pytania blokujące obejmują brak formalnej polityki bezpieczeństwa fizycznego oraz brak elektronicznej kontroli dostępu do serwerowni.
\generujKryteria{pytania.csv}{BEZPIECZENSTWO_FIZYCZNE}

\subsection{Sekcja system informatyczny}
\noindent Aktywacja: \textit{aktywacja warunkowa, gdy typ obiektu} = \texttt{SYSTEM\_INFORMATYCZNY}\\
Pytania w tej sekcji służą do oceny bezpieczeństwa i~zgodności konkretnego systemu informatycznego rozpatrywanego jako samodzielna jednostka. Zakres oceny jest stosowany odrębnie dla każdego audytowanego systemu, co pozwala zidentyfikować różnice w poziomach zabezpieczeń między poszczególnymi aplikacjami w portfelu ICT organizacji. Sprawdzane jest, czy system figuruje w oficjalnym rejestrze zasobów ICT z przypisaną klasyfikacją krytyczności i~właścicielem, czy wymusza uwierzytelnianie wieloskładnikowe (MFA) dla wszystkich kont, czy stosuje aktualny protokół szyfrowania komunikacji TLS 1.2 lub nowszy, a logi są centralnie agregowane i~przechowywane w niezmienialnym storage przez co najmniej 12 miesięcy. Ocenie podlegają również polityka haseł wymuszana technicznie przez system, zarządzanie podatnościami z udokumentowanymi terminami naprawy oraz, jako kryterium o szczególnym znaczeniu dla ciągłości działania, kopie zapasowe przechowywane w izolowanej, geograficznie odrębnej lokalizacji, z empirycznie weryfikowanymi wskaźnikami RTO i~RPO. Pytania blokujące obejmują brak wpisu w rejestrze CMDB, brak MFA oraz brak backupu.
\generujKryteria{pytania.csv}{APLIKACJA}

\subsection{Sekcja usługi w chmurze}
\noindent Aktywacja: \textit{aktywacja warunkowa, gdy} $\textit{utrzymanie\_chmura}(v) \in \{\texttt{CHMURA\_K3}, \texttt{CHMURA\_EOG}\}$\\
Pytania w tej sekcji służą do oceny zgodności wynikającej ze specyfiki przetwarzania danych w środowisku chmurowym, gdzie organizacja powierza przetwarzanie zewnętrznemu dostawcy usług (CSP). Sprawdzane jest, czy przed uruchomieniem systemu przeprowadzono i~udokumentowano Ocenę Skutków dla Ochrony Danych (DPIA) oraz czy jest ona aktualizowana przy każdej istotnej zmianie systemu lub dostawcy. Ocenie podlega model zarządzania kluczami kryptograficznymi, uniemożliwiając dostawcy samodzielny dostęp do przechowywanych danych. Weryfikowane jest również, czy umowa z dostawcą chmury zapewnia prawo do audytu, bezpieczną i~pełną izolację danych między klientami (multi-tenancy).
\generujKryteria{pytania.csv}{CHMURA}

\subsection{Sekcja Transfer Impact Assessment (TIA)}
\noindent Aktywacja: \textit{aktywacja warunkowa, gdy} $\textit{utrzymanie\_chmura}(v) = \texttt{CHMURA\_K3}$\\
Pytania w tej sekcji służą do oceny zgodności transferów danych osobowych realizowanych poza Europejski Obszar Gospodarczy (EOG). Obowiązek przeprowadzenia TIA wynika z wyroku Trybunału Sprawiedliwości UE w sprawie Schrems II z 2020 r., który uchylił mechanizm Privacy Shield i~nałożył na administratorów obowiązek indywidualnej weryfikacji, czy prawo i~praktyka państwa docelowego nie podważają skuteczności ochrony gwarantowanej przez Standardowe Klauzule Umowne (SCC). Sprawdzane jest, czy dla każdego transferu przeprowadzono TIA uwzględniające ustawodawstwo krajowe państwa odbiorcy (w tym ryzyko eksterytorialnego dostępu wynikające z amerykańskiego CLOUD Act i~FISA 702) a także czy transfer jest obsługiwany jako komplet trzech dokumentów: aktualnej umowy DPA zgodnej z art. 28 RODO, klauzul SCC z 2021 r. oraz dokumentacji TIA. Weryfikowane jest ponadto, czy organizacja posiada pisemne oświadczenie dostawcy o nieprzekazywaniu danych poza EOG, obejmujące backupy, logi i~metadane, oraz czy oświadczenie to jest technicznie weryfikowalne. Kryterium blokującym jest brak ważnej umowy DPA, bez niej jakikolwiek transfer danych osobowych jest z mocy prawa niedopuszczalny.

\paragraph{Trigger identyfikacji poddostawców i~TIA}\mbox{}\\
\label{sec:trigger-tia}

Jednym z najpoważniejszych i~najczęściej pomijanych problemów współczesnej praktyki compliance jest zjawisko \textit{deklaratywnego zatrzymania łańcucha} analizy poddostawców: 
organizacja lub jej dostawca deklaruje zgodność z wymogami regulacyjnymi dotyczącymi miejsca przetwarzania danych, nie weryfikując jednak,
czy deklaracja ta jest prawdziwa w odniesieniu do wszystkich ogniw łańcucha dostaw ICT leżących poniżej.

Klasycznym przykładem jest firma zewnętrzna,
która informuje audytowaną organizację:
\textit{,,Hostujemy Państwa dane wyłącznie
na serwerach zlokalizowanych w Unii Europejskiej''}
i jako potwierdzenie wskazuje na regulamin usługi amerykańskiego dostawcy chmury, wskazujący region lokalizacji danych w europejskim obszarze gospodarczym. Jeśli serwer w EOG obsługiwany jest przez podmiot podlegający CLOUD Act oznacza to, że władze federalne USA mogą żądać dostępu
do przechowywanych na nim danych bez wiedzy i~zgody europejskiego podmiotu danych, a okoliczność ta powinna być objęta procedurą Transfer
Impact Assessment (TIA), jednak w praktyce audytowej jest nagminnie pomijana.

ORION rozwiązuje ten problem przez formalny mechanizm triggera identyfikacji poddostawców, który jest zaimplementowany w ankiecie w rozdziale \ref{sec:ankieta-chmura} poprzez pytania dotyczące kraju pochodzenia podmiotu hostującego oraz lokalizacji przechowywania danych.

\begin{definition}[Trigger identyfikacji poddostawcy]
Trigger identyfikacji poddostawcy jest automatycznym
mechanizmem rozszerzania grafu ICT, aktywowanym gdy
obiekt $v$ spełnia co najmniej jeden warunek
ze zbioru $\mathcal{C}_{TIA}$. Po aktywacji trigger
generuje obowiązek audytowy rozszerzenia grafu $G$
o wierzchołki reprezentujące poddostawców obiektu $v$
oraz przeprowadzenia pełnego audytu ORION
dla każdego nowo odkrytego wierzchołka.
\end{definition}

Zbiór warunków aktywujących trigger:

\begin{multline}
\mathcal{C}_{TIA}(v) =
\bigl(\textit{pochodzenie podmiotu}(v) \neq \texttt{EOG}\bigr)
\;\vee\;
\bigl(\textit{lokalizacja danych}(v) \neq \texttt{EOG}\bigr) \\
\;\vee\;
\bigl(\textit{wsparcie systemu}(v) \neq \texttt{EOG}\bigr)
\label{eq:ctia}
\end{multline}


Gdy $\mathcal{C}_{TIA}(v) = \texttt{TRUE}$,
do kwestionariusza obiektu $v$ dodawana jest sekcja \texttt{TIA}, zawierająca pytania dotyczące ryzyk przekazanie danych poza EOG.
\\
\generujKryteria{pytania.csv}{TIA}

% Brak pozytywnej odpowiedzi na pytania TIA z dowodem jakości co najmniej $d = 0{,}8$ aktywuje flagę blokującą:

% \begin{equation}
% \text{FLAG}_{TIA}(v) =
% \mathcal{C}_{TIA}(v)
% \;\wedge\;
% \Bigl(a_{TIA\text{-}01} = 0
% \;\vee\; d_{TIA\text{-}01} < 0{,}8
% \;\vee\; a_{TIA\text{-}02} = 0\Bigr)
% \label{eq:flag-tia}
% \end{equation}

% Gdy $\text{FLAG}_{TIA}(v) = \texttt{TRUE}$,
% obiekt jest klasyfikowany jako \texttt{NIEZGODNY}
% niezależnie od wyników pozostałych pytań,
% a graf ICT jest obowiązkowo rozszerzany
% o poddostawców obiektu $v$.

\subsection{Sekcja dostawcy ICT}
\noindent Aktywacja: \textit{aktywacja warunkowa, gdy typ obiektu} = \texttt{SYSTEM\_INFORMATYCZNY}\\
Pytania w tej sekcji służą do oceny zarządzania ryzykiem stron trzecich w odniesieniu do zewnętrznych dostawców ICT wspierających oceniany system lub funkcje krytyczne organizacji. Sprawdzane jest, czy przed zawarciem umowy przeprowadzono due diligence obejmujące weryfikację certyfikatów bezpieczeństwa, lokalizacji przetwarzania danych, listy sub-procesorów oraz wyników testów penetracyjnych środowiska dostawcy. Ocenie podlega kompletność wymogów umownych: czy kontrakt zawiera mierzalne wskaźniki SLA z sankcjami za ich niedotrzymanie, prawo do audytu onsite, obowiązek powiadomienia o istotnych zmianach (np. zmianie sub-procesora) z wyprzedzeniem co najmniej 30 dni, klauzulę exit umożliwiającą pełny eksport danych w otwartym formacie bez dodatkowych opłat oraz klauzulę regresu umożliwiającą dochodzenie zwrotu kar regulacyjnych nałożonych z winy dostawcy.
\generujKryteria{pytania.csv}{DOSTAWCY}

\subsection{Sekcja exit plan}
\noindent Aktywacja: \textit{aktywacja warunkowa, gdy typ obiektu} = \texttt{SYSTEM\_INFORMATYCZNY} \textbf{ORAZ} \textit{klasyfikacja podmiotu} $\in \{\texttt{FINANSOWY}, \texttt{KLUCZOWY}\}$\\
Pytania w tej sekcji służą do oceny gotowości organizacji do zakończenia współpracy z dostawcą ICT w sposób kontrolowany, bez utraty danych i~bez zakłócenia ciągłości działania obsługiwanych funkcji krytycznych. Exit plan jest traktowany jako dokument żywy, zatwierdzony przez zarząd i~aktualizowany przy każdej istotnej zmianie systemu lub dostawcy, zawierający konkretne kroki techniczne migracji, a nie jedynie deklaratywny opis intencji. Sprawdzane jest, czy możliwy jest pełny eksport danych w otwartym formacie niezależnym od narzędzi własnościowych dostawcy oraz czy eksport był faktycznie przetestowany poza środowiskiem dostawcy. Ocenie podlega empirycznie zmierzony czas pełnej migracji (RTE), zidentyfikowanie alternatywnego dostawcy zdolnego do przejęcia funkcji oraz posiadanie procedury weryfikacji bezpiecznego i~certyfikowanego usunięcia wszystkich danych z infrastruktury dostawcy po zakończeniu współpracy.
\generujKryteria{pytania.csv}{EXIT_PLAN}

\subsection{Sekcja plan ciągłości działania (BCP)}
\noindent Aktywacja: \textit{aktywacja warunkowa, gdy} $\textit{krytyczność}(v) \in \{\texttt{KRYTYCZNY}, \texttt{WAŻNY}\}$ \textbf{ORAZ} \textit{klasyfikacja podmiotu} $\in \{\texttt{FINANSOWY}, \texttt{KLUCZOWY}\}$\\
Pytania w tej sekcji służą do oceny zdolności organizacji do utrzymania lub sprawnego odtworzenia funkcji krytycznych w przypadku poważnego incydentu ICT, awarii dostawcy lub katastrofy. Sprawdzane jest, czy organizacja posiada formalne, zatwierdzone przez zarząd plany BCP i~DRP zawierające konkretne kroki techniczne i~sekwencje przełączeń oraz czy każdy system krytyczny ma zdefiniowane i~technicznie weryfikowalne wskaźniki RTO i~RPO. Ocenie podlega zakres scenariuszy ujętych w planach: czy obejmują ransomware, utratę centrum danych, niedostępność dostawcy chmury i~awarię kaskadową w łańcuchu dostaw, a dla każdego scenariusza zdefiniowano odrębną ścieżkę odtwarzania.
\generujKryteria{pytania.csv}{PLAN_CIAGLOSCI_DZIALANIA}

\subsection{Sekcja koncentracja ryzyka}
\noindent Aktywacja: \textit{aktywacja warunkowa, gdy obiekt obsługuje} $\geq 3$ \textit{systemy podmiotu}\\
Pytania w tej sekcji służą do oceny narażenia organizacji na ryzyko nadmiernego uzależnienia od ograniczonej liczby zewnętrznych dostawców ICT, grup powiązanych dostawców lub lokalizacji geograficznych, co może skutkować systemowym zakłóceniem funkcji krytycznych w przypadku pojedynczej awarii lub incydentu. Ryzyko koncentracji jest oceniane w czterech wymiarach: podmiotowym: uzależnienie od jednego dostawcy lub grupy kapitałowej; geograficznym: koncentracja przetwarzania w jednej lokalizacji lub jurysdykcji prawnej; technologicznym:  vendor lock-in wynikający z zastrzeżonych formatów danych lub interfejsów API uniemożliwiających migrację; oraz systemowym:  wspólna infrastruktura bazowa leżąca u podstaw pozornie niezależnych dostawców warstwy aplikacyjnej.
\generujKryteria{pytania.csv}{KONCENTRACJA_RYZYKA}

\begin{table}[H]
\centering
\caption{Podsumowanie pytań audytowych ORION i
         warunki aktywacji}
\label{tab:sekcje-aktywacja}
\renewcommand{\arraystretch}{1.5}
\resizebox{\textwidth}{!}{%
\begin{tabular}{p{5cm}p{3cm}p{8cm}}
\toprule
\textbf{Sekcja} & \textbf{Typ} & \textbf{Warunek aktywacji $\Phi_k$} \\
\midrule
Organizacja & Warunkowa & \textit{typ obiektu} = \texttt{PODMIOT} \\
\midrule
Bezpieczeństwo techniczne & Warunkowa & \textit{typ obiektu} = \texttt{PODMIOT} \\
\midrule
Bezpieczeństwo fizyczne & Warunkowa & \textit{typ obiektu} = \texttt{PODMIOT} \\
\midrule
System informatyczny & Warunkowa & \textit{typ obiektu} = \texttt{SYSTEM\_INFORMATYCZNY} \\
\midrule
Usługi w chmurze & Warunkowa & $\textit{utrzymanie\_chmura}(v) \in \{\texttt{CHMURA\_K3}, \texttt{CHMURA\_EOG}\}$ \\
\midrule
Transfer Impact Assessment (TIA) & Warunkowa & $\textit{utrzymanie\_chmura}(v) = \texttt{CHMURA\_K3}$ \\
\midrule
Dane poufne: osobowe, wrażliwe oraz tajemnice przedsiębiorstwa & Warunkowa & $\textit{typ\_danych}(v) \in \{D_1, D_2, D_3\}$ \\
\midrule
Dostawcy ICT & Warunkowa & \textit{typ obiektu} = \texttt{SYSTEM\_INFORMATYCZNY} \\
\midrule
Plan ciągłości działania (BCP) & Warunkowa &
  $\textit{krytyczność}(v) \in \{\texttt{KRYTYCZNY}, \texttt{WAŻNY}\}$
  \textbf{ORAZ} klasyfikacja podmiotu
  $\in \{\texttt{FINANSOWY}, \texttt{KLUCZOWY}\}$ \\
\midrule
Exit plan & Warunkowa &
  \textit{typ obiektu} $= \texttt{SYSTEM\_INFORMATYCZNY}$
  \textbf{ORAZ} klasyfikacja podmiotu
  $\in \{\texttt{FINANSOWY}, \texttt{KLUCZOWY}\}$ \\
\midrule
Koncentracja ryzyka & Warunkowa &
  Obiekt obsługuje $\geq 3$ systemy podmiotu \\
\bottomrule
\end{tabular}}

\end{table}



% ============================================================
\section{Faza 4: Obliczenie scoringu
          i~propagacja ryzyka}
\label{sec:faza4}
% ============================================================

Po zebraniu wszystkich odpowiedzi i~dowodów
dla wszystkich obiektów w grafie ICT system ORION
oblicza wyniki scoringowe na trzech poziomach
agregacji: pytania, obiektu i~podmiotu.

\subsection{Poziom 1: Wynik pytania}
\label{sec:wynik-pytania}

Wynik pojedynczego pytania $q_j$ jest wyznaczany
jako ważona średnia efektywnych wartości odpowiedzi
ze wszystkich perspektyw, które pytanie to otrzymały:

\begin{equation}
S(q_j) =
\frac{\displaystyle\sum_{\pi_i \in \Pi_{q_j}}
      V_{eff}(q_j, \pi_i) \cdot w_{\pi_i}}
     {\displaystyle\sum_{\pi_i \in \Pi_{q_j}} w_{\pi_i}}
\label{eq:wynik-pytania}
\end{equation}

\noindent gdzie $\Pi_{q_j}$ jest zbiorem perspektyw,
którym zadano pytanie $q_j$, a $w_{\pi_i}$
jest wagą perspektywy. W rdzeniu metodycznym
wszystkie perspektywy mają równą wagę
$w_{\pi_i} = 1$, co upraszcza wzór do
 średniej arytmetycznej efektywnych
wartości odpowiedzi. Wagi perspektyw mogą być
różnicowane w rozszerzeniach domenowych,
na przykład w module finansowym DORA perspektywa
zgodności może otrzymać wyższą wagę w pytaniach
dotyczących obowiązków sprawozdawczych.

\subsection{Poziom 2: Wynik sekcji}
\label{sec:wynik-sekcji}

Wynik sekcji tematycznej $\mathcal{S}_k$ jest
wyznaczany jako ważona średnia wyników wszystkich
pytań należących do tej sekcji, z uwzględnieniem
obligatoryjności pytania.

Wynik sekcji $\mathcal{S}_k$ jest wyznaczany dwuetapowo.
Najpierw sprawdzany jest warunek blokujący:
jeśli $\text{FLAG}_{block}(\mathcal{S}_k) = \texttt{TRUE}$,
wynik sekcji wynosi zero niezależnie od pozostałych odpowiedzi:

\begin{equation}
S(\mathcal{S}_k) = 0
\qquad
\text{jeśli } \text{FLAG}_{block}(\mathcal{S}_k) = \texttt{TRUE}
\label{eq:wynik-sekcji-blok}
\end{equation}

\noindent W przeciwnym razie wynik sekcji jest wyznaczany
jako ważona średnia wyników pytań:

\begin{equation}
S(\mathcal{S}_k) =
\frac{%
  \sum\limits_{q_j \in \mathcal{S}_k} S(q_j) \cdot \omega(q_j)%
}{%
  \sum\limits_{q_j \in \mathcal{S}_k} \omega(q_j)%
}
\label{eq:wynik-sekcji-avg}
\end{equation}


\noindent gdzie $\omega(q_j)$ jest wagą pytania
wynikającą z jego obligatoryjności: 

\begin{equation}
\omega(q_j) =
\begin{cases}
2{,}0 & \text{jeśli pytanie jest obligatoryjne
               (\textsc{OBOWIĄZKOWE})} \\
1{,}0 & \text{jeśli pytanie jest zalecane
               (\textsc{ZALECANE})} \\
0{,}5 & \text{jeśli pytanie jest informacyjne
               (\textsc{INFORMACYJNE})}
\end{cases}
\label{eq:omega}
\end{equation}

Pytania oznaczone jako \textsc{OBOWIĄZKOWE} wynikają
bezpośrednio z przepisów prawa i~mają podwojoną wagę
względem pytań zalecanych. Ich wynik $S(q_j) = 0$, to znaczy brak zgodności bez żadnego dowodu,
jest dodatkowo oznaczany flagą audytową niezależnie
od agregacji sekcji.


\subsection{Poziom 3: Wynik obiektu}
\label{sec:wynik-obiektu}


Wynik obiektu $v$ jest wyznaczany jako
\textit{średnia arytmetyczna} wyników wszystkich
aktywnych sekcji. Wszystkie sekcje pytań audytowych
są traktowane jako równoważne pod względem znaczenia
regulacyjnego ($\rho_k = 1{,}0$ dla każdej sekcji
$\mathcal{S}_k$):

\begin{equation}
S_{total}(v) =
\frac{1}{|\mathcal{S}(v)|}
\sum_{\mathcal{S}_k \in \mathcal{S}(v)}
S(\mathcal{S}_k)
\label{eq:wynik-obiektu}
\end{equation}

Flaga blokująca przy pytaniu z odpowiedzią
\texttt{NIE} ustawia wartość sekcji na zero,
co bezpośrednio obniża $S_{total}(v)$
niezależnie od wyników pozostałych sekcji:

\begin{equation}
\text{FLAG}_{block}(\mathcal{S}_k) = \texttt{TRUE}
\iff
\exists\, q_j \in \mathcal{S}_k :
\kappa(q_j) = \texttt{BLOKUJĄCE}
\;\wedge\;
a_j = 0
\label{eq:flag-block-sekcja}
\end{equation}

\noindent Wynik $S_{total}(v) \in [0, 1]$ jest
interpretowany na pięciostopniowej skali zgodności:

\begin{table}[H]
\centering
\caption{Skala interpretacji wyniku obiektu}
\label{tab:skala-wyniku}
\renewcommand{\arraystretch}{1.4}
\begin{tabular}{p{3cm}p{3cm}p{7cm}}
\toprule
\textbf{Wynik $S_{total}$} & \textbf{Klasyfikacja}
& \textbf{Interpretacja} \\
\midrule
$[0{,}85;\ 1{,}0]$   & \textsc{ZGODNY}
& Wysoki poziom zgodności; dowody kompletne;
  brak istotnych niezgodności \\
\midrule
$[0{,}70;\ 0{,}85)$  & \textsc{CZĘŚCIOWO ZGODNY}
& Większość wymogów spełniona; luki w dokumentacji
  lub pojedyncze niezgodności drugorzędne \\
\midrule
$[0{,}50;\ 0{,}70)$  & \textsc{WYMAGA POPRAWY}
& Znaczące luki; wymogi obowiązkowe niespełnione
  lub spełnione bez wystarczających dowodów \\
\midrule
$[0{,}25;\ 0{,}50)$  & \textsc{NIEZGODNY}
& Poważne naruszenia; wysokie ryzyko sankcji;
  pilne działania naprawcze \\
\midrule
$[0;\ 0{,}25)$        & \textsc{KRYTYCZNIE NIEZGODNY}
& Brak podstawowych mechanizmów zgodności;
  ryzyko natychmiastowych sankcji regulatora \\
\bottomrule
\end{tabular}
\end{table}

\textit{Uwaga:} Aktywacja $\text{FLAG}_{block}^{obj}(v)$
lub $\text{FLAG}_{TIA}(v)$ nadpisuje klasyfikację
do \textsc{NIEZGODNY} niezależnie od wartości $S_{total}$.
Aktywacja flagi \textsc{BLOKUJĄCY} z pytania obligatoryjnego
nadpisuje klasyfikację do \textsc{KRYTYCZNIE NIEZGODNY}.

\noindent Niezależnie od wartości $S_{total}(v)$,
aktywacja flagi $\text{FLAG}_{TIA}(v) = \texttt{TRUE}$
nadpisuje klasyfikację do \textsc{NIEZGODNY},
a aktywacja flagi \textsc{BLOKUJĄCY} z pytania
obligatoryjnego nadpisuje klasyfikację do
\textsc{KRYTYCZNIE NIEZGODNY}.
Żadna agregacja arytmetyczna nie może ukryć
fundamentalnej niezgodności z przepisem prawa
o charakterze bezwzględnie obowiązującym.



\subsection{Poziom 4: Wynik podmiotu
            z propagacją ryzyka}
\label{sec:wynik-podmiotu}

Wynik podmiotu-korzenia $v_0$ jest wyznaczany
jako ważona średnia wyników wszystkich obiektów
w grafie ICT, z uwzględnieniem łącznej wagi obiektu
$w_{total}(v)$ oraz współczynnika deprecjacji
głębokości $\lambda$:

\begin{equation}
S_{podmiot} =
\frac{\displaystyle\sum_{v \in V \setminus \{v_0\}}
      S_{total}(v)
      \cdot w_{total}(v)
      \cdot \lambda^{\,depth(v)}}
     {\displaystyle\sum_{v \in V \setminus \{v_0\}}
      w_{total}(v)
      \cdot \lambda^{\,depth(v)}}
\label{eq:Spodmiot_rozszerzone}
\end{equation}

\noindent gdzie:
\begin{enumerate}[label=(\roman*)]
  \item $w_{total}(v) = w_{krit}(v) \cdot w_{data}(v)$
  jest łączną wagą obiektu wynikającą z jego
  krytyczności i~klasyfikacji przetwarzanych danych
  zgodnie z tabelą~\ref{tab:krytycznosc},
  \item $depth(v)$ jest głębokością wierzchołka $v$
  w grafie ICT, czyli odległością od korzenia $v_0$
  mierzoną liczbą krawędzi,
  \item $\lambda \in (0, 1]$ jest konfigurowalnym
  współczynnikiem deprecjacji głębokości
  (wartość domyślna: $\lambda = 0{,}8$),
  określającym, jak mocno niezgodności głębszych
  warstw łańcucha dostaw wpływają na wynik podmiotu.
\end{enumerate}

Interpretacja współczynnika $\lambda$ jest następująca:
przy $\lambda = 1{,}0$ niezgodności poddostawcy
trzeciego poziomu wpływają na wynik podmiotu
z taką samą siłą jak niezgodności systemu własnego;
przy $\lambda = 0{,}8$ każdy kolejny poziom głębokości
zmniejsza wpływ o 20\%. Wartość $\lambda$ może być
konfigurowana przez audytora w zależności od
oceny faktycznej ekspozycji podmiotu na ryzyko
łańcucha dostaw.

\subsection{Identyfikacja dostawców kluczowych}
\label{sec:dostawcy-kluczowi}

Dostawca kluczowy (\textit{critical ICT third-party service provider}) to dostawca, którego niezgodność lub niedostępność stanowi zagrożenie dla stabilności podmiotu lub całego systemu. ORION identyfikuje dostawców kluczowych formalnie na podstawie warunku alternatywnego wystarczy spełnienie \textbf{co najmniej jednego} z poniższych kryteriów:

\begin{definition}[Dostawca kluczowy]
Dostawca $v$ jest uznawany za \textbf{kluczowego}
(\textsc{KeyVendor}), jeżeli spełnia co najmniej
jeden z następujących warunków:
\begin{enumerate}[label=(\roman*)]
  \item obsługuje co najmniej jeden system
        o klasyfikacji \textbf{KRYTYCZNY}:
        \[
          \exists\, u \in V :
          (v, u) \in E_{svc}
          \;\wedge\;
          \textit{krytyczność}(u) = \texttt{KRYTYCZNY}
        \]
  \item obsługuje co najmniej $n_{KV}^{*} = 3$
        systemy podmiotu (niezależnie od ich
        klasyfikacji krytyczności):
        \[
          \bigl|\{u \in V : (v,u) \in E_{svc}\}\bigr|
          \geq n_{KV}^{*}
        \]
\end{enumerate}
Formalnie:

\begin{multline}
\textsc{KeyVendor}(v) = \texttt{TRUE}
\iff
\Bigl(
  \exists\, u \in V :
  (v,u) \in E_{svc}
  \wedge
  \textit{krytyczność}(u) = \texttt{KRYTYCZNY}
\Bigr)
\\
\;\vee\;
\Bigl(
  |\{u \in V : (v,u) \in E_{svc}\}|
  \geq n_{KV}^{*}
\Bigr)
\label{eq:keyvendor}
\end{multline}

\noindent gdzie $E_{svc}$ jest zbiorem krawędzi
usługowych w grafie ICT, a $n_{KV}^{*} = 3$
jest minimalną liczbą systemów obsługiwanych
przez dostawcę (wartość domyślna, konfigurowalna
przez audytora).
\end{definition}

Innymi słowy: dostawca jest kluczowy, jeśli
obsługuje choćby jeden system krytyczny 
{niezależnie od liczby obsługiwanych systemów} lub jeśli obsługuje co najmniej trzy systemy podmiotu, niezależnie od ich poziomu krytyczności.
Spełnienie wyłącznie jednego warunku {jest wystarczające}.

Sekcja \texttt{KONCENTRACJA\_RYZYKA} jest aktywowana
\textbf{wyłącznie} wtedy, gdy dostawca obsługuje
co najmniej $n_{KV}^{*} = 3$ systemy podmiotu
(warunek~2 z definicji powyżej), niezależnie
od tego, czy stał się dostawcą kluczowym
na podstawie warunku~1 czy~2:

\begin{equation}
\Phi_{\texttt{KONCENTRACJA\_RYZYKA}}(v) = \texttt{TRUE}
\iff
\bigl|\{u \in V : (v,u) \in E_{svc}\}\bigr|
\geq n_{KV}^{*}
\label{eq:koncentracja-aktywacja}
\end{equation}

\noindent Uzasadnienie tego rozróżnienia jest
następujące: ryzyko koncentracji w rozumieniu
DORA Art.~29 materializuje się przez uzależnienie
od jednego dostawcy obsługującego wiele systemów
jednocześnie. Dostawca obsługujący jeden system
krytyczny jest kluczowy z uwagi na wagę tego systemu,
lecz nie generuje ryzyka koncentracji w sensie
regulacyjnym,  stąd sekcja \texttt{KONCENTRACJA\_RYZYKA} nie jest dla niego aktywowana.

Wynik sekcji \texttt{KONCENTRACJA\_RYZYKA} wchodzi do agregacji wyniku obiektu zgodnie
z równaniem~\eqref{eq:wynik-obiektu}.
Karty dostawców kluczowych są bezpośrednio adresowane do zarządu i~rady nadzorczej,
gdyż to na tych organach spoczywa odpowiedzialność za nadzór nad kluczowymi zależnościami ICT
w rozumieniu DORA Art.~5 i~KSC Art.~8.
% ============================================================
\section{Faza 5: Raport końcowy i~rekomendacje}
\label{sec:faza5}
% ============================================================

Wynikiem audytu ORION jest ustrukturyzowany raport
zawierający cztery elementy: macierz zgodności,
graf ICT z oznaczeniami, listę priorytetyzowanych
rekomendacji oraz kartę dostawców kluczowych.

\subsection{Macierz zgodności}
\label{sec:macierz-zgodnosci}

Macierz zgodności przedstawia wyniki $S_{total}(v)$
dla wszystkich obiektów w grafie ICT, pogrupowane
według sekcji tematycznych. Umożliwia to audytorowi
i zarządowi szybką identyfikację obszarów,
w których poziom zgodności jest najniższy, niezależnie od tego, których konkretnych
systemów lub dostawców te obszary dotyczą.

\begin{table}[H]
\centering
\caption{Przykładowy fragment macierzy zgodności
         (ilustracja struktury)}
\label{tab:macierz-zgodnosci}
\renewcommand{\arraystretch}{1.4}
\resizebox{\textwidth}{!}{%
\begin{tabular}{p{4cm}p{2cm}p{2cm}p{2cm}p{2cm}p{2cm}}
\toprule
\textbf{Sekcja / Obiekt}
& \textbf{System ERP} & \textbf{CRM SaaS}
& \textbf{Dostawca cloud} & \textbf{Poddos-tawca CDN}
& \textbf{Średnia} \\
\midrule
Organizacja
& $0{,}91$ & $0{,}85$ & $0{,}78$ & $0{,}62$ & $0{,}79$ \\
\midrule
Bezpieczeństwo techniczne
& $0{,}88$ & $0{,}72$ & $0{,}65$ & $0{,}41$ & $0{,}67$ \\
\midrule
Transfer Impact Assessment
& --- & $0{,}40$ & $0{,}38$ & $0{,}20$ & $0{,}33$ \\
\midrule
Poddostawcy ICT (TIA)
& --- & $0{,}10$ & $0{,}15$ & $0{,}10$ & $0{,}12$ \\
\midrule
Plan ciągłości działania
& $0{,}80$ & $0{,}75$ & $0{,}50$ & --- & $0{,}68$ \\
\midrule
Exit plan
& --- & $0{,}30$ & $0{,}25$ & --- & $0{,}28$ \\
\midrule
\textbf{Wynik obiektu $S_{total}$}
& $0{,}86$ & $0{,}55$ & $0{,}48$ & $0{,}31$ & \\
\midrule
\textbf{Klasyfikacja}
& \textsc{zgodny}
& \textsc{wymaga poprawy}
& \textsc{niezgodny}
& \textsc{niezgodny} & \\
\bottomrule
\end{tabular}%
}
\end{table}

\noindent Wiersze sekcji \texttt{TIA}
i \texttt{EXIT\_PLAN} w powyższej tabeli ilustrują
kluczową obserwację metodyczną: najniższe wyniki
skupiają się konsekwentnie w sekcjach dotyczących
łańcucha dostaw i~transferu danych, a nie
w obszarach bezpieczeństwa technicznego, które
są tradycyjnie najlepiej monitorowane.
Jest to bezpośrednie potwierdzenie luki badawczej
zidentyfikowanej w sekcji~\ref{sec:luka}.

\subsection{Graf ICT z oznaczeniami}
\label{sec:graf-ict}


Graf ICT generowany jako element raportu końcowego jest wizualną reprezentacją skierowanego grafu ważonego $G = (V, E, W)$ zbudowanego w toku przeprowadzenia audytu. Wierzchołki $V$ reprezentują obiekty audytowe, tj.: systemy informatyczne typu \texttt{SYSTEM\_INFORMATYCZNY} oraz dostawców ICT,  a krawędzie skierowane $E \subseteq V \times V$ wskazują kierunek zależności usługowej: od systemu do podmiotu zewnętrznego, który go obsługuje. Wagi $W$ krawędzi odzwierciedlają poziom krytyczności obsługiwanego systemu zgodnie z klasyfikacją przyjętą w Sekcji~\ref{sec:krok23}.

Graf spełnia podwójną funkcję audytową. Po pierwsze, jest \emph{narzędziem komunikacji wyników do zarządu}, pozwala w kilka sekund zidentyfikować, na którym poziomie łańcucha dostaw ICT koncentruje się ryzyko regulacyjne. Po drugie, jest \emph{dowodem audytowym} dokumentującym kompletność przeprowadzonej analizy i~wszystkie odkryte ogniwa łańcucha, włącznie z poddostawcami ujawnionymi przez mechanizm triggera TIA opisany w Sekcji~\ref{sec:trigger-tia}.

\subsubsection{Notacja wierzchołków}
\label{sec:notacja-wierzcholkow}

Każdy wierzchołek grafu jest identyfikowany przez \emph{kod literowy z indeksem głębokości} $\mathrm{depth}(v)$, rozumianej jako odległość od korzenia $v_0$ mierzona liczbą krawędzi. Przyjęta konwencja nazewnicza jest następująca:

\begin{itemize}
    \item $v_0$: podmiot-korzeń, czyli audytowana organizacja (głębokość~0);
    \item $v_1, v_2, \ldots, v_n$: systemy informatyczne własne podmiotu, zidentyfikowane w Kroku~2.1 (głębokość~1);
    \item $d_1, d_2, \ldots, d_m$: bezpośredni dostawcy ICT, zidentyfikowani w Kroku~2.2 (głębokość~2);
    \item $p_1, p_2, \ldots, p_k$: poddostawcy odkryci przez trigger TIA w toku Fazy~3 (głębokość~3 i~wyżej).
\end{itemize}

W przypadku rozległych grafów dopuszcza się stosowanie notacji złożonej, np.\ $d_{1.1}, d_{1.2}$ dla kolejnych dostawców powiązanych z systemem $v_1$, co ułatwia czytelność przy dużej liczbie obiektów.

\subsubsection{Oznaczenia klasyfikacji zgodności}
\label{sec:oznaczenia-zgodnosci}

Biorąc pod uwagę ograniczenia techniczne związane z technologią druku i~ograniczonej dostępności druku kolorowego,  raport audytowy może być drukowany w wersji czarno-białej, klasyfikacja zgodności obiektu jest kodowana \emph{symbolem wewnątrz wierzchołka}, a nie wyłącznie kolorem. Tabela~\ref{tab:oznaczenia-grafu} przedstawia kompletny słownik oznaczeń.

\begin{table}[ht]
\renewcommand{\arraystretch}{1.4}
\caption{Słownik oznaczeń wierzchołków w grafie ICT}
\label{tab:oznaczenia-grafu}
\begin{tabular}{p{2cm} p{4cm} p{4.5cm} p{3cm}}
\hline
\textbf{Symbol} & \textbf{Klasyfikacja} & \textbf{Warunek} & \textbf{Priorytet naprawczy} \\
\hline
$[\checkmark]$ & ZGODNY & $S_{total} \geq 0{,}85$ & brak \\
$[\sim]$ & CZĘŚCIOWO ZGODNY & $0{,}70 \leq S_{total} < 0{,}85$ & P3 \\
$[!]$   & WYMAGA POPRAWY   & $0{,}50 \leq S_{total} < 0{,}70$ & P2 \\
$[\times]$ & NIEZGODNY & $0{,}25 \leq S_{total} < 0{,}50$ & P1--P2 \\
$[\times\!]$ & KRYTYCZNIE NIEZGODNY & $S_{total} < 0{,}25$ & P0 \\
$[\times!!]$ & Aktywna flaga \texttt{FLAG\_TIA} \newline lub \texttt{BLOKUJĄCY} & niezależnie od $S_{total}$ -- nadpisuje klasyfikację & P0--P1 \\
\hline
\end{tabular}
\end{table}

Aktywacja flagi \texttt{FLAG\_TIA} lub flagi \texttt{BLOKUJĄCY} z pytania obligatoryjnego nadpisuje klasyfikację obiektu niezależnie od wartości $S_{total}$, co jest sygnalizowane przez dodanie symbolu~\texttt{!!} do oznaczenia wierzchołka. Żadna agregacja arytmetyczna nie może ukryć fundamentalnej niezgodności z przepisem prawa o charakterze bezwzględnie obowiązującym.

\subsubsection{Oznaczenia krawędzi}
\label{sec:oznaczenia-krawedzi}

Krawędzie grafu ICT różnicowane są typem linii oraz opcjonalną etykietą:

\begin{enumerate}[label=(\roman*)]
    \item \textit{Linia ciągła} ($\rightarrow$) -- standardowa zależność usługowa; system korzysta z usług dostawcy na podstawie umowy;
    \item \textit{Linia przerywana} ($\dashrightarrow$) -- krawędź odkryta przez trigger TIA; oznacza poddostawcę niewidocznego na etapie Fazy~2, ujawnionego dopiero w toku analizy kwestionariusza szczegółowego;
    \item \textit{Etykieta \texttt{[KV]}} na krawędzi, dostawca kluczowy (\textit{KeyVendor}), spełniający co najmniej jeden warunek z definicji~\ref{def:dostawca-kluczowy}: obsługuje co najmniej jeden system KRYTYCZNY lub obsługuje co najmniej $n_{KV} = 3$ systemy podmiotu.
\end{enumerate}

\subsubsection{Przykładowy graf ICT}
\label{sec:przykladowy-graf}

Poniższy przykład ilustruje strukturę grafu ICT dla przykładowego podmiotu. Podmiot $v_0$ eksploatuje trzy systemy informatyczne:

\begin{enumerate}[label=(\roman*)]
    \item $v_1$ - system ERP, klasyfikacja KRYTYCZNY, wdrożony on-premise, wynik $S_{total} = 0{,}91$ ($[\checkmark]$~ZGODNY);
    \item $v_2$ - system CRM, klasyfikacja WAŻNY, chmura EOG, wynik $S_{total} = 0{,}78$ ($[\sim]$~CZĘŚCIOWO ZGODNY);
    \item $v_3$ - portal kliencki SaaS, klasyfikacja STANDARDOWY, chmura kraju trzeciego (USA), wynik $S_{total} = 0{,}48$ ($[\times]$~NIEZGODNY).
\end{enumerate}

Systemy $v_1$ i~$v_2$ obsługiwane są przez dostawcę $d_1$ (firma cloudowa z EOG, $S_{total} = 0{,}72$, $[\sim]$~CZĘŚCIOWO ZGODNY). System $v_3$ obsługiwany jest przez dostawcę $d_2$ (firma z USA, $S_{total} = 0{,}31$, $[\times!!]$~- \texttt{FLAG\_TIA} aktywna). Dostawca $d_1$ obsługuje dwa systemy podmiotu ($n < n_{KV}$), zatem nie spełnia progu dostawcy kluczowego. Trigger TIA aktywowany dla $d_2$ rozszerza graf o poddostawcę $p_1$ -- sieć CDN zlokalizowaną w USA ($S_{total} = 0{,}20$, $[\times!]$~KRYTYCZNIE NIEZGODNY), obsługiwaną przez dostawcę $p_2$ (podmiot zarejestrowany poza EOG, $S_{total}$ nieustalony na etapie Fazy~2).

\begin{figure}[htbp]
  \centering
  \begin{tikzpicture}[
    podmiot/.style={rectangle, rounded corners=6pt,
      draw=black, thick, fill=gray!25,
      minimum width=3.2cm, minimum height=0.9cm,
      align=center, font=\small\bfseries},
    sys/.style={rectangle, rounded corners=3pt,
      draw=black, thick,
      minimum width=3.6cm, minimum height=1.2cm,
      align=center, font=\scriptsize},
    vendor/.style={ellipse, draw=black, thick,
      minimum width=3.6cm, minimum height=1.2cm,
      align=center, font=\scriptsize},
    cfok/.style={fill=white},
    cfwarn/.style={fill=white, pattern=north east lines, pattern color=gray!50},
    cfwarn2/.style={fill=white, pattern=north west lines, pattern color=gray!50},
    cfniezg/.style={fill=white, pattern=crosshatch, pattern color=gray!60},
    cfcrit/.style={fill=white, pattern=crosshatch dots, pattern color=gray},
    cfunk/.style={fill=white, pattern=vertical lines, pattern color=gray!40},
    flagtia/.style={double, double distance=2pt, ultra thick},
    dep/.style={-{Stealth}, thick},
    tiadep/.style={-{Stealth}, thick, dashed},
    lbl/.style={font=\tiny, inner sep=1pt}
  ]
  % depth 0: PODMIOT
  \node[podmiot] (v0) at (0,0)
    {$v_0$\\[2pt]Podmiot audytowany};
  % depth 1: SYSTEM\_INFORMATYCZNY
  \node[sys, cfok] (v1) at (-4.5,-3)
    {$v_1$ \textbf{ERP}\\KRYTYCZNY, on-premise\\$S\!=\!0{,}91$~$[\checkmark]$};
  \node[sys, cfwarn] (v2) at (0,-3)
    {$v_2$ \textbf{CRM}\\WAŻNY, \texttt{CHMURA\_EOG}\\$S\!=\!0{,}78$~$[\sim]$};
  \node[sys, cfniezg] (v3) at (4.5,-3)
    {$v_3$ \textbf{Portal SaaS}\\STD., \texttt{CHMURA\_K3}\\$S\!=\!0{,}48$~$[\times]$};
  % depth 2: DOSTAWCY
  \node[vendor, cfwarn] (d1) at (-2.2,-6.5)
    {$d_1$ Dostawca EOG\\$S\!=\!0{,}72$~$[\sim]$};
  \node[vendor, cfniezg, flagtia] (d2) at (4.5,-6.5)
    {$d_2$ Firma z USA\\\texttt{FLAG\_TIA}\\$S\!=\!0{,}31$~$[\times!!]$};
  % depth 3: poddostawca TIA
  \node[vendor, cfcrit] (p1) at (4.5,-10)
    {$p_1$ Sieć CDN (USA)\\KR.~NIEZGODNY\\$S\!=\!0{,}20$~$[\times!]$};
  % depth 4: poddostawca TIA
  \node[vendor, cfunk] (p2) at (4.5,-13)
    {$p_2$ Podmiot poza EOG\\$S$ nieustalony};
  % krawędzie: v0 -> systemy
  \draw[dep] (v0) -- (v1);
  \draw[dep] (v0) -- (v2);
  \draw[dep] (v0) -- (v3);
  % krawędzie: systemy -> dostawcy (waga = krytyczność systemu)
  \draw[dep] (v1) -- node[lbl, left] {$w\!=\!3$} (d1);
  \draw[dep] (v2) -- node[lbl, right] {$w\!=\!2$} (d1);
  \draw[dep] (v3) -- node[lbl, right] {$w\!=\!1$} (d2);
  % krawędzie TIA (odkryte przez trigger)
  \draw[tiadep] (d2) -- node[lbl, right, text=red!70!black] {\textsc{tia}} (p1);
  \draw[tiadep] (p1) -- node[lbl, right, text=red!70!black] {\textsc{tia}} (p2);
  % etykiety głębokości
  \foreach \y/\d in {0/0, -3/1, -6.5/2, -10/3, -13/4}
    \node[font=\tiny, text=gray!80, anchor=west] at (7.5,\y) {głęb.~\d};
  \end{tikzpicture}

  \medskip\noindent
  \begin{scriptsize}
  \begin{tabular}{@{}cl@{\hspace{14pt}}cl@{}}
    \tikz[baseline=-0.5ex]\draw[fill=white, draw=black, thick]
      (0,0) rectangle (0.4,0.25); &
    ZGODNY ($S\!\geq\!0{,}85$) &
    \tikz[baseline=-0.5ex]\draw[fill=white, pattern=crosshatch,
      pattern color=black!60, draw=black, thick]
      (0,0) rectangle (0.4,0.25); &
    NIEZGODNY ($0{,}25\!\leq\!S\!<\!0{,}50$) \\[2pt]
    \tikz[baseline=-0.5ex]\draw[fill=white, pattern=north east lines,
      pattern color=black!50, draw=black, thick]
      (0,0) rectangle (0.4,0.25); &
    CZĘŚCIOWO ZGODNY ($0{,}70\!\leq\!S\!<\!0{,}85$) &
    \tikz[baseline=-0.5ex]\draw[fill=white, pattern=crosshatch dots,
      pattern color=black, draw=black, thick]
      (0,0) rectangle (0.4,0.25); &
    KRYTYCZNIE NIEZGODNY ($S\!<\!0{,}25$) \\[2pt]
    \tikz[baseline=-0.5ex]\draw[fill=white, pattern=north west lines,
      pattern color=black!50, draw=black, thick]
      (0,0) rectangle (0.4,0.25); &
    WYMAGA POPRAWY ($0{,}50\!\leq\!S\!<\!0{,}70$) &
    \tikz[baseline=-0.5ex]\draw[double, double distance=2pt, ultra thick]
      (0,0.12) ellipse (0.35 and 0.18); &
    \texttt{FLAG\_TIA} aktywna (podwójna obwódka) \\[2pt]
    \tikz[baseline=-0.5ex]\draw[-{Stealth}, dashed, thick]
      (0,0) -- (0.7,0); &
    krawędź odkryta przez TIA &
    \multicolumn{2}{l}{} \\
  \end{tabular}
  \end{scriptsize}

  \caption{Przykładowy graf ICT podmiotu z oznaczeniami klasyfikacji zgodności.
    Prostokąty oznaczają systemy informatyczne (\texttt{SYSTEM\_INFORMATYCZNY}), elipsy oznaczają dostawców ICT.
    Podwójna obwódka sygnalizuje aktywną flagę \texttt{FLAG\_TIA} lub kryterium blokujące.
    Linia przerywana oznacza krawędź odkrytą przez trigger TIA. Wagi krawędzi odzwierciedlają
    poziom krytyczności obsługiwanego systemu ($w\!=\!3$: KRYTYCZNY, $w\!=\!2$: WAŻNY,
    $w\!=\!1$: STANDARDOWY).}
  \label{fig:przykladowy-graf-ict}
\end{figure}




Wynik podmiotu $S_{podmiot}$ obliczany jest zgodnie z równaniem~\eqref{eq:Spodmiot_rozszerzone}, przy czym wierzchołki odkryte przez trigger TIA (głębokość~3) są objęte współczynnikiem deprecjacji głębokości $\lambda^{\mathrm{depth}(v)} = 0{,}8^3 = 0{,}512$, co oznacza, że niezgodność poddostawcy CDN~($p_1$) wpływa na wynik podmiotu z siłą ok.~51\% siły niezgodności systemu własnego.

\subsubsection{Interpretacja grafu przez zarząd}
\label{sec:interpretacja-grafu}

Graf ICT jest jedynym elementem raportu przeznaczonym do odczytania bez uprzedniego zapoznania się z treścią kwestionariuszy. Zarząd i~rada nadzorcza powinni być w stanie, wyłącznie na podstawie grafu, odpowiedzieć na trzy pytania operacyjne:

\begin{enumerate}[label=(\roman*)]
    \item \textit{Gdzie skupia się ryzyko?} Wierzchołki oznaczone $[\times]$, $[\times!]$ lub $[\times!!]$ wymagają natychmiastowej uwagi, niezależnie od ich głębokości w grafie.
    \item \textit{Który dostawca jest kluczowy?} Krawędzie z etykietą \texttt{[KV]} wskazują podmioty, których awaria lub niezgodność zagraża stabilności całej organizacji.
    \item \textit{Jak głęboko sięga łańcuch?} Liczba poziomów grafu ujawnia, czy organizacja jest świadoma pełnej ekspozycji na ryzyko poddostawców, co jest wymogiem art.~28 DORA i~art.~6 ust.~2 dyrektywy NIS2.
\end{enumerate}

Brak wierzchołków na poziomie głębokości~3 i~wyżej \emph{nie} oznacza braku poddostawców, oznacza brak przeprowadzenia analizy TIA. Audytor jest zobowiązany do odnotowania w raporcie, czy głębokość grafu wynika z faktycznego braku poddostawców podlegających triggerowi, czy też z niewykonania należytej staranności przez audytowany podmiot lub jego dostawców.

\subsection{Lista rekomendacji z priorytetami}
\label{sec:rekomendacje}

Każda niezgodność zidentyfikowana w toku audytu
generuje rekomendację przypisaną do jednego
z czterech poziomów priorytetu:

\begin{table}[H]
\centering
\caption{Priorytety rekomendacji naprawczych w ORION}
\label{tab:priorytety}
\renewcommand{\arraystretch}{1.5}
\begin{tabular}{p{2.5cm}p{3cm}p{3.5cm}p{4cm}}
\toprule
\textbf{Priorytet} & \textbf{Definicja}
& \textbf{Termin działania}
& \textbf{Warunek nadania} \\
\midrule
\textbf{P0} \newline \textsc{Krytyczny}
& Naruszenie przepisu bezwzględnie obowiązującego
  grożące natychmiastową sankcją regulatora
& Natychmiast\newline(max.\ 72~h)
& Flaga \textsc{BLOKUJĄCY}
  lub $S_{total}(v) < 0{,}25$ \\
\midrule
\textbf{P1} \newline \textsc{Wysoki}
& Istotna niezgodność z przepisami prawa
  lub brak TIA przy przetwarzaniu danych osobowych
& 30~dni
& $\text{FLAG}_{TIA} = \texttt{TRUE}$
  lub pytanie obligatoryjne
  z $S(q_j) < 0{,}3$ \\
\midrule
\textbf{P2} \newline \textsc{Średni}
& Niezgodność z zaleceniami regulacyjnymi
  lub luki w dokumentacji
& 90~dni
& $S_{total}(v) \in [0{,}25;\ 0{,}50)$
  lub pytanie zalecane z $d < 0{,}5$ \\
\midrule
\textbf{P3} \newline \textsc{Niski}
& Obszar do usprawnienia;
  brak aktualnej dokumentacji lub dowodów słabych
& 180~dni
& $S_{total}(v) \in [0{,}50;\ 0{,}70)$
  lub odpowiedź z $d = 0{,}5$ \\
\bottomrule
\end{tabular}
\end{table}

\noindent Każda rekomendacja zawiera: identyfikator
niezgodności, numer pytania audytowego i~sekcji,
obiekt którego dotyczy, perspektywę która ujawniła
niezgodność, treść niezgodności sformułowaną
w języku zrozumiałym dla zarządu,
podstawę regulacyjną z numerem artykułu,
sugerowane działanie naprawcze oraz termin
wynikający z priorytetu.

\subsection{Karta dostawców kluczowych}
\label{sec:karta-dostawcow}

Raport końcowy zawiera oddzielną kartę dla każdego
dostawcy spełniającego warunek~\eqref{eq:keyvendor}.
Karta zawiera: wynik $S_{total}$ dostawcy,
listę systemów podmiotu przez niego obsługiwanych,
status TIA i~exit planu, klasyfikację zgodności
ze wskazaniem priorytetów naprawczych oraz
rekomendację dotyczącą zarządzania ryzykiem
koncentracji. Karty dostawców kluczowych
są bezpośrednio adresowane do zarządu i~rady
nadzorczej, gdyż to na tych organach spoczywa
odpowiedzialność za nadzór nad kluczowymi
zależnościami ICT w rozumieniu DORA Art.~5
i KSC Art.~8.


% \section{Pytania audytowe}


% Wywołanie

% Pytania zostały podzielone na sekcje zgodnie z klasyfikacją przyjętą w rodziale \ref{perspektywa-pytań} \textit{Perspektywa}. Każda pozycja odnosi się do konkretnego obszaru zagadnień, zawiera podstawę prawną, definicję istotności oraz oznaczenie (flagę binarną) pytania blokującego.


% \subsection{Sekcja organizacja}
% \generujKryteria{pytania.csv}{ORGANIZACJA}

% \subsection{Sekcja dane poufne: wrażliwe oraz osobowe}
% \generujKryteria{pytania.csv}{DANE_WRAZLIWE}

% \subsection{Sekcja utrzymanie w chmurze}
% \generujKryteria{pytania.csv}{CHMURA}

% \subsection{Sekcja bezpieczeństwo techniczne}
% \generujKryteria{pytania.csv}{BEZPIECZENSTWO_TECH}

% \subsection{Sekcja bezpieczeństwo fizyczne}
% \generujKryteria{pytania.csv}{BEZPIECZENSTWO_FIZYCZNE}

% \subsection{Sekcja system informatyczny}
% \generujKryteria{pytania.csv}{APLIKACJA}

% \subsection{Sekcja Transfer Impact Assessment (TIA)}
% \generujKryteria{pytania.csv}{TIA}

% \subsection{Sekcja dostawcy ICT}
% \generujKryteria{pytania.csv}{DOSTAWCY}

% \subsection{Sekcja exit plan}
% \generujKryteria{pytania.csv}{EXIT_PLAN}

% \subsection{Sekcja plan ciągłości działania (BCP)}
% \generujKryteria{pytania.csv}{PLAN_CIAGLOSCI_DZIALANIA}

% \subsection{Sekcja koncentracja ryzyka}
% \generujKryteria{pytania.csv}{KONCENTRACJA_RYZYKA}




% % ============================================================
% \section{Pełny algorytm audytu ORION}
% \label{sec:algorytm-pelny}
% % ============================================================
% \todo{poprawic, są braki}
% Poniższy algorytm formalnie opisuje kompletny
% przebieg audytu ORION, integrując wszystkie
% elementy opisane w poprzednich sekcjach.

% \begin{algorithm}[H]
% \caption{Pełny algorytm audytu ORION -- część 1/2:
%          Inicjalizacja i~przejście BFS}
% \label{alg:orion-czesc1}
% \begin{algorithmic}[1]

% \Require Podmiot $v_0$,
%          próg dostawcy kluczowego $w_{KV}^{*}$, $n_{KV}^{*}$,
%          współczynnik deprecjacji $\lambda$
% \Ensure  Wyniki $\mathcal{R}$,
%          klasyfikacja podmiotu $S_{podmiot}$,
%          lista dostawców kluczowych $\mathcal{K}$,
%          lista flag $\mathcal{F}$,
%          lista rekomendacji $\mathcal{REC}$

% \medskip
% \Statex \Comment{\textbf{Faza 1: Klasyfikacja podmiotu}}
% \State $\textit{reżim}(v_0) \leftarrow
%        \textsc{KlasyfikacjaRegulatoryjna}(v_0)$
% \State $\textit{zasięg\_geo}(v_0) \leftarrow
%        \textsc{ZakresGeograficzny}(v_0)$

% \medskip
% \Statex \Comment{\textbf{Faza 2: Inicjalizacja grafu ICT}}
% \State $V \leftarrow \{v_0\}$,\quad
%        $E \leftarrow \emptyset$,\quad
%        $Q_{BFS} \leftarrow [v_0]$
% \State $\textit{Visited} \leftarrow \emptyset$,\quad
%        $\mathcal{R} \leftarrow \emptyset$,\quad
%        $\mathcal{K} \leftarrow \emptyset$,\quad
%        $\mathcal{F} \leftarrow \emptyset$,\quad
%        $\mathcal{REC} \leftarrow \emptyset$

% \medskip
% \Statex \Comment{Inwentaryzacja systemów i~dostawców pierwszego poziomu}
% \ForAll{$v \in \textsc{InwentaryzacjaSystemów}(v_0)$}
%     \State $V \leftarrow V \cup \{v\}$
%     \State $E \leftarrow E \cup \{(v_0, v)\}$
%     \State $\textsc{Enqueue}(Q_{BFS}, v)$
% \EndFor

% \medskip
% \Statex \Comment{\textbf{Fazy 3--4: Przejście BFS po grafie ICT}}
% \While{$Q_{BFS} \neq \emptyset$}
%     \State $v \leftarrow \textsc{Dequeue}(Q_{BFS})$
%     \If{$v \in \textit{Visited}$}
%         \textbf{continue}
%     \EndIf
%     \State $\textit{Visited} \leftarrow \textit{Visited} \cup \{v\}$

%     \medskip
%     \Statex \Comment{Klasyfikacja atrybutów obiektu (cross-validation)}
%     \State $\mathcal{A}(v) \leftarrow
%            \textsc{KlasyfikujAtrybuty}(v)$
%     \State $w_{total}(v) \leftarrow
%            w_{krit}(v) \cdot w_{data}(v)$

%     \medskip
%     \Statex \Comment{Generowanie kwestionariusza szczegółowego}
%     \State $\mathcal{S}(v) \leftarrow
%            \textsc{AktywujSekcje}\bigl(v,\,
%            \mathcal{A}(v),\,
%            \textit{reżim}(v_0)\bigr)$

%     \medskip
%     \Statex \Comment{Sprawdzenie triggera TIA}
%     \If{$\mathcal{C}_{TIA}(v)$}
%         \State $\mathcal{S}(v) \leftarrow
%                \mathcal{S}(v) \cup
%                \{\texttt{PODDOSTAWCY\_TIA}\}$
%     \EndIf

%     \medskip
%     \Statex \Comment{Zbieranie odpowiedzi i~dowodów}
%     \State $\mathcal{R}[v] \leftarrow
%            \textsc{ZbierzOdpowiedzi}\bigl(v,\,
%            \mathcal{S}(v)\bigr)$

%     \medskip
%     \Statex \Comment{Obliczenie wyników sekcji z obsługą flag blokujących}
%     \ForAll{$\mathcal{S}_k \in \mathcal{S}(v)$}
%         \If{$\exists\, q_j \in \mathcal{S}_k :
%              \kappa(q_j) = \texttt{BLOKUJĄCE}
%              \;\wedge\; a_j = 0$}
%             \State $S(\mathcal{S}_k) \leftarrow 0$
%             \State $\mathcal{F} \leftarrow \mathcal{F} \cup
%                    \{(v,\, \mathcal{S}_k,\,
%                    \textsc{BLOKUJĄCY\_SEKCJA})\}$
%         \Else
%             \State $S(\mathcal{S}_k) \leftarrow
%                    \dfrac{%
%                      \displaystyle\sum_{q_j \in \mathcal{S}_k}
%                      S(q_j) \cdot \omega(q_j)%
%                    }{%
%                      \displaystyle\sum_{q_j \in \mathcal{S}_k}
%                      \omega(q_j)%
%                    }$
%         \EndIf
%     \EndFor

%     \medskip
%     \Statex \Comment{Obliczenie scoringu obiektu}
%     \State $S_{total}(v) \leftarrow
%            \dfrac{%
%              \displaystyle\sum_{\mathcal{S}_k \in \mathcal{S}(v)}
%              S(\mathcal{S}_k) \cdot \rho_k%
%            }{%
%              \displaystyle\sum_{\mathcal{S}_k \in \mathcal{S}(v)}
%              \rho_k%
%            }$

% \EndWhile

% \end{algorithmic}
% \end{algorithm}

% \begin{algorithm}[H]
% \caption{Pełny algorytm audytu ORION -- część 2/2:
%          Flagi, rekomendacje i~agregacja}
% \label{alg:orion-czesc2}
% \begin{algorithmic}[1]
% \setcounter{ALG@line}{42}

%     \medskip
%     \Statex \Comment{Sprawdzenie flagi blokującej na poziomie obiektu}
%     \If{$\exists\, (v,\, \mathcal{S}_k,\,
%          \textsc{BLOKUJĄCY\_SEKCJA}) \in \mathcal{F}$}
%         \State $\mathcal{F} \leftarrow \mathcal{F} \cup
%                \{(v,\, \textsc{BLOKUJĄCY\_OBJ})\}$
%         \State $S_{total}(v) \leftarrow
%                \min\bigl(S_{total}(v),\; 0{,}49\bigr)$
%                \Statex \Comment{Wymuszenie klasy \textsc{NIEZGODNY}}
%     \EndIf

%     \medskip
%     \Statex \Comment{Sprawdzenie flagi TIA}
%     \If{$\text{FLAG}_{TIA}(v)$}
%         \State $\mathcal{F} \leftarrow \mathcal{F} \cup
%                \{(v,\, \textsc{BLOKUJĄCY\_TIA})\}$
%         \State $S_{total}(v) \leftarrow
%                \min\bigl(S_{total}(v),\; 0{,}49\bigr)$
%                \Statex \Comment{Wymuszenie klasy \textsc{NIEZGODNY}}

%         \medskip
%         \Statex \Comment{Dynamiczne rozszerzanie grafu ICT}
%         \ForAll{$u \in \textsc{IdentyfikujPoddostawców}(v)$}
%             \If{$u \notin V$}
%                 \State $V \leftarrow V \cup \{u\}$
%                 \State $E \leftarrow E \cup \{(v,\, u)\}$
%                 \State $\textsc{Enqueue}(Q_{BFS},\, u)$
%             \EndIf
%         \EndFor
%     \EndIf

%     \medskip
%     \Statex \Comment{Generowanie rekomendacji z priorytetami}
%     \If{$(v,\, \textsc{BLOKUJĄCY\_OBJ}) \in \mathcal{F}$
%         \textbf{ or }
%         $(v,\, \textsc{BLOKUJĄCY\_TIA}) \in \mathcal{F}$}
%         \State $\textit{priorytet} \leftarrow \textbf{P0}$
%     \ElsIf{$\exists\, q_j :
%             \omega(q_j) = \textsc{OBOWIĄZKOWE}
%             \;\wedge\; S(q_j) < 0{,}3$}
%         \State $\textit{priorytet} \leftarrow \textbf{P1}$
%     \ElsIf{$S_{total}(v) \in [0{,}25;\; 0{,}50)$}
%         \State $\textit{priorytet} \leftarrow \textbf{P2}$
%     \Else
%         \State $\textit{priorytet} \leftarrow \textbf{P3}$
%     \EndIf
%     \State $\mathcal{REC} \leftarrow \mathcal{REC} \cup
%            \textsc{GenerujRekomendacje}\bigl(v,\,
%            \mathcal{R}[v],\,
%            S_{total}(v),\,
%            \mathcal{F},\,
%            \textit{priorytet}\bigr)$

%     \medskip
%     \Statex \Comment{Sprawdzenie warunku dostawcy kluczowego}
%     \If{$\textsc{IsKeyVendor}\bigl(v,\,
%         w_{total}(v),\, n_{KV}^{*}\bigr)$}
%         \State $\mathcal{K} \leftarrow \mathcal{K} \cup \{v\}$
%     \EndIf

%     \medskip
%     \Statex \Comment{Standardowe przejście BFS -- dodanie sąsiadów}
%     \ForAll{$(v,\, u) \in E$}
%         \If{$u \notin \textit{Visited}$}
%             \State $\textsc{Enqueue}(Q_{BFS},\, u)$
%         \EndIf
%     \EndFor

% \medskip
% \Statex \Comment{\textbf{Faza 4 końcowa: Agregacja wyniku podmiotu}}
% \State $S_{podmiot} \leftarrow
%        \dfrac{%
%          \displaystyle\sum_{v \in V \setminus \{v_0\}}
%          S_{total}(v) \cdot w_{total}(v)
%          \cdot \lambda^{\,depth(v)}%
%        }{%
%          \displaystyle\sum_{v \in V \setminus \{v_0\}}
%          w_{total}(v) \cdot \lambda^{\,depth(v)}%
%        }$

% \medskip
% \Statex \Comment{Nadpisanie wyniku gdy aktywna flaga krytyczna obiektu}
% \If{$\exists\, v \in V :
%      (v,\, \textsc{BLOKUJĄCY\_OBJ}) \in \mathcal{F}$}
%     \State $S_{podmiot} \leftarrow
%            \min\bigl(S_{podmiot},\; 0{,}24\bigr)$
%            \Statex \Comment{Wymuszenie klasy \textsc{KRYTYCZNIE NIEZGODNY}}
% \EndIf

% \medskip
% \Statex \Comment{\textbf{Faza 5: Raport końcowy}}
% \State \Return $S_{podmiot}$,\ $\mathcal{R}$,\
%        $\mathcal{K}$,\ $\mathcal{F}$,\ $\mathcal{REC}$

% \end{algorithmic}
% \end{algorithm}
\section{Wyniki audytu}

W celu praktycznej weryfikacji metodologii ORION przeprowadzono audyt technologiczno-prawny w jednym z banków spółdzielczych działających na terenie Polski. Na wniosek instytucji dane identyfikacyjne banku zostały zanonimizowane; wyniki audytu zostały natomiast przedstawione w pełnym zakresie merytorycznym. Audyt objął kompletny cykl klasyfikacyjny przewidziany metodologią: identyfikację reżimu regulacyjnego, klasyfikację ekspozycji sieciowej, ocenę poziomu krytyczności oraz określenie modelu utrzymania infrastruktury i jurysdykcji podmiotu hostującego.

Przedmiotem badania był system Cyfrowego Trwałego Nośnika (CTN), pełniący kluczową funkcję w procesie masowej komunikacji banku z klientami-konsumentami. Podstawowym zadaniem systemu jest dystrybucja dokumentów oraz oświadczeń woli w sposób zgodny z restrykcyjnymi wymogami prawnymi. Architektura systemu realizuje postanowienia Unijnej Dyrektywy o Prawach Konsumenta (Dyrektywa 2011/83/UE) oraz krajowych przepisów implementujących jej założenia, w szczególności Ustawy o prawach konsumenta. Regulacje te nakładają na przedsiębiorców, w tym instytucje sektora finansowego, obowiązek przekazywania konsumentom informacji (takich jak regulaminy, taryfy opłat i prowizji czy formularze informacyjne) na trwałym nośniku, co gwarantuje integralność treści oraz wyklucza możliwość ich jednostronnej modyfikacji przez bank.

\subsection{Klasyfikacja podmiotu}

W pierwszej kolejności zespół audytorski ustalił reżim regulacyjny, jakiemu podlega badana instytucja. Bank spółdzielczy, jako podmiot sektora bankowego, został jednoznacznie zaklasyfikowany jako \textbf{podmiot finansowy} w rozumieniu Rozporządzenia DORA (\textit{Digital Operational Resilience Act}). Kwalifikacja ta wynika bezpośrednio z art.~2 ust.~1 lit.~a DORA, który na pierwszym miejscu katalogu podmiotów objętych rygorem odporności cyfrowej wymienia instytucje kredytowe.

Klasyfikacja ta determinuje zakres obowiązków banku w obszarze zarządzania ryzykiem ICT: obejmuje wymóg wdrożenia ram zarządzania ryzykiem (art.~5--16 DORA), prowadzenia rejestru umów z dostawcami ICT (art.~28 DORA) oraz przeprowadzania zaawansowanych testów odporności operacyjnej, w tym testów penetracyjnych opartych na analizie zagrożeń: TLPT (\textit{Threat-Led Penetration Testing}).

\subsection{Klasyfikacja atrybutu ekspozycji sieciowej}

Kolejnym krokiem było określenie poziomu ekspozycji sieciowej systemu CTN z zastosowaniem ankiety dostępności metodologii ORION. Analiza architektury systemu wykazała, że dokumenty bankowe oraz interfejs użytkownika są udostępniane masowemu odbiorcy bezpośrednio za pośrednictwem sieci Internet i standardowych przeglądarek internetowych, bez wymogu połączenia VPN ani dostępu do sieci wewnętrznej banku.

Taka architektura wyczerpuje przesłanki \textit{Kryterium D1: Dostęp spoza sieci organizacji} (system jest osiągalny z dowolnego węzła sieci publicznej) oraz \textit{Kryterium D4: Publiczny adres lub domena}, usługa jest opublikowana pod publicznie rozwiązywalną nazwą domenową. Publiczna ekspozycja systemu rozszerza wektor ataku i implikuje konieczność wdrożenia zaawansowanych zabezpieczeń perymetrycznych, w szczególności zapory aplikacyjnej (WAF), mechanizmów mitygacji ataków DDoS oraz inspekcji ruchu TLS.

W rezultacie system CTN został zaklasyfikowany jako \textbf{system publiczny}.

\subsection{Klasyfikacja atrybutu poziomu krytyczności}

Ocena krytyczności systemu CTN została przeprowadzona z zastosowaniem sześciokryterialnej ankiety metodologii ORION (Kryteria K1--K6). Poniższa tabela przedstawia wyniki punktacji dla każdego z kryteriów wraz z uzasadnieniem.

\begin{table}[H]
\centering
\caption{Ocena punktowa krytyczności systemu CTN}
\renewcommand{\arraystretch}{1.4}
\begin{tabular}{p{2.1cm} p{5.2cm} p{1.4cm} p{5cm}}
\toprule
\textbf{Kryterium} & \textbf{Opis} & \textbf{Wynik} & \textbf{Uzasadnienie} \\
\midrule
K1 & Wpływ finansowy i operacyjny niedostępności & NIE & Awaria nie zatrzymuje transakcji; BCP przewiduje zastępczy kanał papierowy \\
\midrule
K2 & Przetwarzanie danych osobowych lub wrażliwych & TAK & Spersonalizowane dokumenty i dane PII klientów (imię, numer umowy, dane adresowe) \\
\midrule
K3 & Podleganie regulacjom zewnętrznym & TAK & Nadzór KNF (Rekomendacja~D) i wymogi Prezesa UOKiK dot.\ trwałego nośnika \\
\midrule
K4 & Zależności innych systemów produkcyjnych & NIE & Brak systemów operacyjnie uzależnionych od CTN \\
\midrule
K5 & Ryzyko wycieku lub utraty danych & TAK & Naruszenie poufności rodzi obowiązek zgłoszenia incydentu do UODO i ryzyko reputacyjne \\
\midrule
K6 & Obsługa procesów podstawowej działalności & NIE & Awaria CTN nie blokuje realizacji transakcji płatniczych ani procesów core-bankingowych \\
\midrule
\multicolumn{2}{l}{\textbf{Suma odpowiedzi twierdzących}} & \textbf{3 TAK} & \textbf{System ważny} \\
\bottomrule
\end{tabular}
\end{table}

Wynik 3 odpowiedzi twierdzących (K2, K3, K5) plasuje system CTN w kategorii \textbf{systemu ważnego} zgodnie z matrycą klasyfikacji metodologii ORION. Kluczowe znaczenie dla ostatecznej kwalifikacji ma zasada dominacji operacyjnej: brak spełnienia kryteriów K1 i K6 (wyznaczników bezpośredniego wpływu na ciągłość podstawowej działalności) wyklucza możliwość uznania systemu za krytyczny, pomimo spełnienia trzech pozostałych kryteriów. Wdrożone plany ciągłości działania (BCP) przewidują procedurę awaryjną polegającą na tymczasowym zastąpieniu kanału cyfrowego klasyczną korespondencją papierową, co pozwala zachować pełną zgodność prawną nawet przy długotrwałej niedostępności systemu informatycznego.

Status systemu ważnego nie zwalnia organizacji z obowiązków nadzorczych. Nakłada on bezwzględny wymóg: ciągłego monitorowania systemu, regularnego przeprowadzania testów penetracyjnych, audytów zgodności regulacyjnej oraz utrzymania aktualnej dokumentacji przetwarzania danych osobowych zgodnie z art.~30 RODO. Parametry RTO i RPO nie muszą być skalibrowane w reżimie minutowym, jednak powinny zostać formalnie określone i zweryfikowane w ramach ćwiczeń BCP.

\subsection{Klasyfikacja atrybutu utrzymania w chmurze}

Ostatnim krokiem procedury audytowej było ustalenie modelu utrzymania infrastruktury systemu CTN oraz identyfikacja jurysdykcji podmiotu hostującego, z zastosowaniem ankiety klasyfikacji atrybutu utrzymania w chmurze metodologii ORION (Kryteria C1-C6). Wyniki dla każdego kryterium przedstawia poniższa tabela.

\begin{table}[H]
\centering
\caption{Ocena punktowa klasyfikacji modelu utrzymania systemu CTN}
\renewcommand{\arraystretch}{1.4}
\begin{tabular}{p{2.1cm} p{5.2cm} p{1.4cm} p{5cm}}
\toprule
\textbf{Kryterium} & \textbf{Opis} & \textbf{Wynik} & \textbf{Uzasadnienie} \\
\midrule
C1 & Model wdrożenia infrastruktury & TAK & Infrastruktura zaimplementowana u zewnętrznego polskiego dostawcy chmury \\
\midrule
C2 & Odpowiedzialność operacyjna & TAK & Dostępność zasobów i warstwa sprzętowa w gestii dostawcy na podstawie SLA \\
\midrule
C3 & Umowa outsourcingowa lub chmurowa & TAK & Umowa SaaS/PaaS i kontrakt outsourcingu IT z dostawcą \\
\midrule
C4 & Lokalizacja i kontrola danych & TAK & Dane produkcyjne i kopie zapasowe w centrum danych na terytorium RP \\
\midrule
C5 & Jurysdykcja podmiotu hostującego & TAK & Wyłącznie kapitał europejski; pełne TIA; brak ekspozycji na US CLOUD Act i FISA 702 \\
\midrule
C6 & Wsparcie techniczne & TAK & Helpdesk i administracja świadczone wyłącznie przez zespół z terytorium EOG; geoblocking VPN \\
\midrule
\multicolumn{2}{l}{\textbf{Suma odpowiedzi twierdzących}} & \textbf{6 TAK} & \textbf{Chmura EOG} \\
\bottomrule
\end{tabular}
\end{table}

Uzyskanie kompletnego zestawu sześciu odpowiedzi twierdzących, w tym pozytywnych wyników dla kryteriów C5 i C6 dotyczących jurysdykcji i wsparcia technicznego, przesądza o klasyfikacji systemu CTN jako \textbf{systemu utrzymywanego w chmurze EOG}. Jest to klasyfikacja o najwyższym poziomie bezpieczeństwa jurydycznego: podmiot hostujący oparty wyłącznie na kapitale europejskim i niemający siedziby w państwach objętych regulacjami o ekstraterytorialnym dostępie do danych (takim jak US CLOUD Act czy FISA Section~702) eliminuje ryzyko nieautoryzowanego transferu danych do państw trzecich.

Klasyfikacja ta zapewnia jednocześnie pełną zgodność z wytycznymi i komunikatami chmurowymi Urzędu Komisji Nadzoru Finansowego (UKNF) oraz z wymogami Rozporządzenia DORA w zakresie zarządzania ryzykiem ze strony zewnętrznych dostawców usług ICT (art.~28-30 DORA). Narzędzia zdalnego dostępu serwisowego skonfigurowane z restrykcyjnym geoblockingiem, uniemożliwiające zestawienie sesji administracyjnej spoza obszaru EOG, stanowią dodatkową warstwę zabezpieczenia weryfikowalną przez audytorów bezpośrednio w logach dostępowych.

\subsection{Klasyfikacja atrybutu typów przetwarzanych danych}

Klasyfikacja typów przetwarzanych danych została przeprowadzona zgodnie z tabelą~\ref{tab:typy-danych} metodologii ORION. System CTN przetwarza spersonalizowane dokumenty bankowe zawierające dane identyfikacyjne klientów (imię i nazwisko, numer umowy, dane adresowe) oraz informacje objęte tajemnicą bankową w rozumieniu art.~104 ustawy Prawo bankowe, a więc dane o rachunkach, historię transakcji i informacje o zobowiązaniach kredytowych. Tajemnica bankowa stanowi informację prawnie chronioną na mocy ustawy sektorowej, co w klasyfikacji ORION odpowiada najwyższej kategorii wrażliwości.

System CTN został zatem sklasyfikowany jako przetwarzający dane kategorii \textbf{\texttt{TAJEMNICA}} (symbol $D_3$, mnożnik wagowy $w_{data} = 2{,}0$). Klasyfikacja ta ma bezpośrednie przełożenie na wagę artefaktu w agregacji końcowej oraz aktywuje sekcję pytań dotyczących danych poufnych. Naruszenie bezpieczeństwa systemu przetwarzającego tajemnicę bankową rodzi po stronie instytucji obowiązek zgłoszenia incydentu do organu nadzoru (KNF) oraz potencjalną odpowiedzialność odszkodowawczą wobec klientów, których dane zostały ujawnione osobom nieuprawnionym.

\subsection{Klasyfikacja atrybutu kraju pochodzenia}

Atrybut kraju pochodzenia dostawcy systemu CTN został wyznaczony na podstawie weryfikacji statusu prawnego i kapitałowego podmiotu świadczącego usługę utrzymania oraz przeprowadzającego wdrożenie. Dostawca jest spółką zarejestrowaną i prowadzącą działalność operacyjną wyłącznie na terytorium Rzeczypospolitej Polskiej, bez powiązań kapitałowych z podmiotami spoza Europejskiego Obszaru Gospodarczego. Zarówno wsparcie techniczne, jak i administracja systemu są świadczone przez personel zlokalizowany w Polsce.

Atrybut przyjął wartość \textbf{\texttt{PL}}. Konsekwencją tej klasyfikacji jest brak obowiązku przeprowadzenia oceny skutków transferu danych (TIA), ponieważ żaden etap przetwarzania nie wiąże się z przekazaniem danych do państwa trzeciego w rozumieniu RODO. Dostawca nie podlega jurysdykcji przepisów o ekstraterytorialnym dostępie do danych (takich jak US CLOUD Act czy FISA Section~702), co eliminuje ryzyko nieautoryzowanego udostępnienia danych organom obcego państwa.

\subsection{Analiza aktywowanych sekcji pytań}

Zbiór atrybutów wyznaczonych w poprzednich krokach determinuje zakres sekcji pytań aktywowanych przez system ORION dla danego audytu. Poniżej przedstawiono wynik tej aktywacji wraz z uzasadnieniem dla każdej sekcji.

Podmiot audytowany (bank spółdzielczy) został sklasyfikowany jako \texttt{FINANSOWY}, system CTN jako \texttt{PUBLICZNY} o poziomie krytyczności \texttt{WAŻNY}, utrzymywany w trybie \texttt{CHMURA\_EOG}, przetwarzający dane kategorii \texttt{TAJEMNICA} ($D_3$), dostarczany przez podmiot z kraju \texttt{PL}. Zidentyfikowano jednego zewnętrznego dostawcę ICT, działającego w jurysdykcji EOG.

Na podstawie powyższych atrybutów aktywowano następujące sekcje pytań:

\begin{itemize}
  \item \textbf{\texttt{ORGANIZACJA}} (aktywacja stała dla obiektów typu \texttt{PODMIOT}): ocena struktury zarządzania ryzykiem ICT, polityk bezpieczeństwa i nadzoru zarządczego nad funkcjami cyfrowymi.

  \item \textbf{\texttt{BEZPIECZENSTWO\_FIZYCZNE}} (aktywacja stała dla obiektów typu \texttt{PODMIOT}): weryfikacja zabezpieczeń fizycznych siedziby i pomieszczeń serwerowni podmiotu audytowanego.

  \item \textbf{\texttt{BEZPIECZENSTWO\_TECHNICZNE}} (aktywacja stała dla obiektów typu \texttt{PODMIOT}): ocena technicznych środków ochrony infrastruktury, zarządzania podatnościami oraz procedur reagowania na incydenty.

  \item \textbf{\texttt{APLIKACJA}} (aktywacja warunkowa dla obiektów typu \texttt{SYSTEM\_INFORMATYCZNY}): weryfikacja bezpieczeństwa warstwy aplikacyjnej systemu CTN, obejmująca m.in.\ mechanizmy uwierzytelniania, zarządzanie sesjami i ochronę przed atakami na warstwę aplikacji (OWASP Top~10).

  \item \textbf{\texttt{CHMURA}} (aktywacja warunkowa: $\textit{utrzymanie\_chmura} \in \{\texttt{CHMURA\_EOG}, \texttt{CHMURA\_K3}\}$): ocena zgodności wynikającej ze specyfiki przetwarzania w środowisku zewnętrznym, w tym weryfikacja modelu zarządzania kluczami kryptograficznymi, izolacji danych między klientami dostawcy oraz prawa do audytu.

  \item \textbf{\texttt{DANE\_POUFNE}} (aktywacja warunkowa: $\textit{typ\_danych}(v) \in \{D_1, D_2, D_3\}$): rozszerzona ocena środków ochrony danych objętych tajemnicą bankową, obejmująca szyfrowanie w spoczynku i podczas transmisji, kontrolę dostępu oraz procedury reagowania na naruszenia ochrony danych.

  \item \textbf{\texttt{DOSTAWCY}} (aktywacja warunkowa: $\textit{utrzymanie\_chmura}(v) \in \{\texttt{CHMURA\_EOG},\\\texttt{CHMURA\_K3}\}$): ocena zarządzania ryzykiem zewnętrznego dostawcy ICT, w tym due diligence przedkontraktowe, kompletność wymogów umownych (SLA, prawo do audytu, klauzula exit) oraz monitorowanie sub-procesorów.

  \item \textbf{\texttt{EXIT\_PLAN}} (aktywacja warunkowa: $\textit{typ obiektu}(v) = \texttt{DOSTAWCA}\; \land\; \textit{klasyfikacja}$ $\textit{podmiotu}(v) = \texttt{FINANSOWY}$): weryfikacja planu wyjścia z relacji z dostawcą, obejmująca zdolność do pełnego eksportu danych w formacie otwartym, zmierzony czas migracji (RTE) oraz identyfikację alternatywnego dostawcy.

  \item \textbf{\texttt{PLAN\_CIAGLOSCI\_DZIALANIA}} (aktywacja warunkowa: $\textit{krytyczność}(v) \in \{\texttt{KRYTYCZNY}, \\\texttt{WAŻNY}\}$ oraz klasyfikacja podmiotu \texttt{FINANSOWY}): ocena planów BCP i DRP, wskaźników RTO i RPO, scenariuszy odtwarzania po awarii dostawcy chmury oraz kopii zapasowych izolowanych od infrastruktury produkcyjnej.
\end{itemize}

Następujące sekcje nie zostały aktywowane z uwagi na niespełnienie warunków progowych:

\begin{itemize}
  \item \textbf{\texttt{TIA}} (warunek: $\textit{utrzymanie\_chmura}(v) = \texttt{CHMURA\_K3}$): sekcja dotycząca oceny skutków transferu danych nie została aktywowana, ponieważ system jest utrzymywany w \texttt{CHMURA\_EOG}. Żaden etap przetwarzania nie obejmuje transferu danych do państwa trzeciego, wobec czego obowiązek przeprowadzenia TIA wymagany przez RODO nie powstaje.

  \item \textbf{\texttt{KONCENTRACJA\_RYZYKA}} (warunek: dostawca obsługuje $\geq 3$ systemy podmiotu): zidentyfikowano jednego zewnętrznego dostawcę ICT, obsługującego wyłącznie system CTN. Próg koncentracji nie został osiągnięty, wobec czego ryzyko nadmiernego uzależnienia od jednego podmiotu nie wymaga pogłębionej oceny w ramach tego audytu.
\end{itemize}

\subsection{Wyniki sekcji ORGANIZACJA}

Sekcja ORGANIZACJA obejmuje dziewięć obszarów tematycznych wymaganych przez regulacje NIS2 i DORA, dotyczących ładu zarządczego, zarządzania ryzykiem ICT oraz zdolności operacyjnej do obsługi incydentów. Ocena każdego obszaru była przeprowadzana metodą cross-validation, tzn. każde pytanie było zadawane niezależnie trzem perspektywom: prawno-compliance, technicznej (dział IT) oraz użytkownika końcowego. Zgodność odpowiedzi ze wszystkich trzech perspektyw stanowiła warunek udzielenia odpowiedzi twierdzącej bez kary za niepotwierdzenie deklaracji.

Bank spółdzielczy uzyskał kompletne, pozytywne wyniki we wszystkich dziewięciu obszarach, bez żadnych niezgodności i bez rozbieżności między perspektywami. Szczegółowe ustalenia dla każdego obszaru przedstawiono poniżej.

\textbf{Strategia cyberbezpieczeństwa} (NIS2 art.~20 ust.~1, DORA art.~5): zarząd banku formalnie zatwierdził strategię cyberbezpieczeństwa i ponosi za nią odpowiedzialność prawną. Dokument jest dostępny dla pracowników, a jego istnienie zostało potwierdzone niezależnie przez radcę prawnego, dział IT oraz losowo wybranych użytkowników końcowych, co świadczy o realnym wdrożeniu, a nie jedynie formalnym istnieniu dokumentu.

\textbf{Szkolenia zarządu} (NIS2 art.~20 ust.~2): członkowie zarządu odbyli obowiązkowe szkolenia z cyberbezpieczeństwa. Dokumentacja szkoleń (lista uczestników, daty, zakresy tematyczne) była dostępna do wglądu audytorów, co eliminuje ryzyko zakwalifikowania szkolenia jako deklaratywnego.

\textbf{Ramy zarządzania ryzykiem ICT} (NIS2 art.~21 ust.~1, DORA art.~6): organizacja wdrożyła udokumentowane ramy zarządzania ryzykiem ICT zatwierdzone przez zarząd, obejmujące rejestr ryzyk z przypisanymi właścicielami i planami mitygacji. Rejestr jest aktualizowany co najmniej raz w roku. Wszyscy respondenci (compliance, IT, użytkownicy) potrafili wskazać osobę odpowiedzialną za bezpieczeństwo systemów IT, co świadczy o wysokiej świadomości organizacyjnej.

\textbf{Ocena ryzyka łańcucha dostaw} (NIS2 art.~21 ust.~2 lit.~d, DORA art.~28): organizacja przeprowadza regularną ocenę ryzyka uwzględniającą łańcuch dostaw ICT i dostawców zewnętrznych. Ocena obejmuje wszystkich dostawców z klasyfikacją krytyczności i udokumentowanymi wynikami. Użytkownicy potwierdzili wiedzę o tym, które zewnętrzne podmioty mają dostęp do systemów lub danych organizacji.

\textbf{Zgłaszanie poważnych incydentów w ciągu 24 godzin} (NIS2 art.~23 ust.~4 lit.~a, DORA art.~19): bank posiada udokumentowaną procedurę umożliwiającą zgłoszenie poważnego incydentu do organu krajowego w wymaganym terminie 24 godzin od wykrycia. System ticketowy wymusza checkpointy czasowe z automatycznym przypomnieniem. Użytkownicy końcowi wykazali się znajomością kroków postępowania w sytuacji podejrzanego zdarzenia, co potwierdza skuteczność szkoleń.

\textbf{Klasyfikacja incydentów} (NIS2 art.~23 ust.~3, DORA art.~18): organizacja dysponuje procedurą klasyfikacji incydentów pozwalającą ocenić, czy dane zdarzenie spełnia kryteria incydentu poważnego lub znaczącego w rozumieniu regulacji. Użytkownicy potrafili odróżnić zwykłą awarię od zdarzenia o charakterze bezpieczeństwa.

\textbf{Powiadamianie klientów o incydentach} (NIS2 art.~23 ust.~4 lit.~c, DORA art.~14): bank posiada zatwierdzoną procedurę komunikacji kryzysowej z gotowymi szablonami i zdefiniowanymi kanałami komunikacji z klientami. Procedura jest możliwa do uruchomienia w krótkim czasie. Użytkownicy potwierdzili świadomość istnienia procedury powiadamiania.

\textbf{Rejestr czynności przetwarzania} (RODO art.~30): organizacja prowadzi aktualny rejestr czynności przetwarzania danych osobowych (ROPA) w formie elektronicznej, obejmujący system CTN z opisem celów, kategorii danych, odbiorców i podstaw prawnych. Rejestr jest dostępny dla Inspektora Ochrony Danych i aktualizowany przy zmianach systemu lub procesu przetwarzania.

\textbf{Rejestr kluczowych dostawców ICT} (DORA art.~28 ust.~2, art.~29 ust.~1): organizacja prowadzi aktualny rejestr zewnętrznych dostawców ICT z klasyfikacją ich krytyczności i przypisaniem do funkcji krytycznych lub istotnych. Rejestr obejmuje dane kontraktowe, zakres usług i daty przeglądów. Użytkownicy wykazali świadomość liczby zewnętrznych dostawców IT obsługujących organizację.

Kompletna zgodność we wszystkich dziewięciu obszarach, potwierdzona niezależnie przez trzy perspektywy, wskazuje na wysoki poziom dojrzałości organizacyjnej banku w zakresie ładu zarządzania ryzykiem ICT. Brak rozbieżności między odpowiedziami działu compliance, działu IT i użytkowników końcowych świadczy o tym, że wdrożone procedury i polityki funkcjonują realnie, a nie tylko jako dokumentacja formalna. Jest to wynik szczególnie istotny w kontekście przeciwdziałania zjawisku compliance theater: potwierdzone, weryfikowalne dowody zastępują deklaratywne zapewnienia, co stanowi jeden z kluczowych postulatów metodologii ORION.

\subsection{Wyniki sekcji BEZPIECZENSTWO\_TECHNICZNE}

Sekcja BEZPIECZENSTWO\_TECHNICZNE jest aktywowana dla obiektów typu \texttt{PODMIOT} i w audycie CTN obejmowała ocenę banku jako organizacji. Obejmuje osiem obszarów dotyczących technicznych środków ochrony infrastruktury ICT, zarządzania dostępem, monitorowania oraz ciągłości procesów bezpieczeństwa, wymaganych przez NIS2, DORA i RODO. We wszystkich ośmiu obszarach bank uzyskał wyniki pozytywne, bez żadnych uchybień i bez rozbieżności między perspektywami. Jest to jeden z silniejszych wyników całego audytu, wskazujący na dojrzałość operacyjną działu IT oraz na skuteczność szkoleń i kultury bezpieczeństwa w organizacji.

\textbf{Testy bezpieczeństwa i testy penetracyjne} (NIS2 art.~21 ust.~2 lit.~m, DORA art.~24--25): organizacja przeprowadza regularne testy bezpieczeństwa, testy podatności oraz testy penetracyjne zgodnie z rocznym planem obejmującym harmonogram i zakres. Wyniki testów są dokumentowane i powiązane z procesem naprawczym. Użytkownicy potwierdzili wiedzę o tym, że zewnętrzne podmioty cyklicznie weryfikują odporność systemów banku.

\textbf{Monitoring i wykrywanie anomalii} (NIS2 art.~21 ust.~2 lit.~b, DORA art.~10): wdrożone mechanizmy ciągłego monitorowania pokrywają systemy krytyczne i generują alerty w czasie rzeczywistym. Za obsługę alertów odpowiada dedykowany zespół, co zostało potwierdzone przez użytkowników, którzy byli świadomi istnienia takiego zespołu i jego roli.

\textbf{Zasada minimalnych uprawnień} (DORA art.~9, NIS2 art.~21): dostęp do systemów krytycznych jest nadawany zgodnie z zasadą least privilege i regularnie przeglądany. Stosowane jest narzędzie klasy PAM (Privileged Access Management) z automatycznym przeglądem uprawnień i alertowaniem przy próbach eskalacji. Użytkownicy potwierdzili, że nie posiadają dostępu do zasobów wykraczających poza zakres ich obowiązków.

\textbf{Szkolenia i świadomość pracowników} (NIS2 art.~21 ust.~2 lit.~g): pracownicy przechodzą regularne szkolenia z cyberbezpieczeństwa dostosowane do roli (odrębne ścieżki dla personelu IT, biznesowego i zarządu). Skuteczność szkoleń jest mierzona przez symulowane kampanie phishingowe. Użytkownicy potwierdzili udział w ćwiczeniach tego rodzaju w ciągu ostatnich 12 miesięcy.

\textbf{Offboarding i usuwanie dostępu} (DORA art.~9): organizacja posiada udokumentowany i zautomatyzowany proces dezaktywacji kont przy odejściu pracownika. Dezaktywacja następuje w dniu odejścia i obejmuje wszystkie systemy bez konieczności ręcznej interwencji administratora w poszczególnych aplikacjach. Użytkownicy wykazali świadomość wymogów w tym zakresie.

\textbf{Bezpieczeństwo poczty elektronicznej i załączników} (RODO art.~32, NIS2 art.~21): organizacja posiada politykę regulującą wysyłanie danych wrażliwych pocztą elektroniczną, z wymogiem szyfrowania załączników zawierających dane klientów. Podczas weryfikacji z perspektywy użytkownika nie odnotowano przypadków wysyłania niezabezpieczonych plików z danymi osobowymi.

\textbf{Segmentacja sieci i kontrola ruchu bocznego} (DORA art.~9, NIS2 art.~21): systemy krytyczne są logicznie odizolowane od pozostałej infrastruktury sieciowej przez segmentację VLAN uniemożliwiającą niekontrolowane przemieszczanie się potencjalnego atakującego między segmentami (lateral movement). Polityki firewalla wymuszają kontrolę ruchu wschód-zachód między segmentami. Użytkownicy rozumieją cel izolacji systemów, co świadczy o skuteczności komunikacji wewnętrznej działu IT.

\textbf{Zarządzanie certyfikatami TLS} (DORA art.~9, NIS2 art.~21 lit.~h): organizacja prowadzi inwentarz certyfikatów TLS z datami wygasania i wdrożonym procesem automatycznego odnawiania. Stosowane narzędzie monitorujące generuje alerty z wyprzedzeniem co najmniej 30 dni przed wygaśnięciem certyfikatu. Użytkownicy nie odnotowali przypadków ostrzeżeń przeglądarki o wygasłym certyfikacie na systemach banku.

Pełna zgodność we wszystkich ośmiu obszarach technicznych, potwierdzona spójnie przez wszystkie trzy perspektywy, świadczy o wysokiej kulturze bezpieczeństwa technicznego w organizacji. Szczególnie istotny jest wynik w obszarze szkoleń i świadomości: użytkownicy końcowi wykazali praktyczną wiedzę z zakresu cyberbezpieczeństwa, co bezpośrednio przekłada się na ograniczenie ryzyka ataków socjotechnicznych (phishing, vishing) jako najczęstszego wektora inicjalnego kompromitacji systemów finansowych.


\subsection{Wyniki sekcji BEZPIECZENSTWO\_FIZYCZNE}

Sekcja BEZPIECZENSTWO\_FIZYCZNE obejmuje osiem obszarów dotyczących ochrony fizycznej obiektów, infrastruktury serwerowej i nośników danych, wymaganych przez regulacje NIS2, DORA oraz RODO. Spośród ośmiu obszarów sześć zostało ocenionych pozytywnie przez wszystkie trzy perspektywy (compliance, IT, użytkownicy). Dwa obszary sklasyfikowane jako zalecane ujawniły nieprawidłowości o wspólnej przyczynie, omówione szczegółowo poniżej.

\textbf{Bezpieczeństwo fizyczne centrów danych} (DORA art.~9, NIS2 art.~21): polityki bezpieczeństwa fizycznego centrum danych spełniają wymogi w zakresie poufności, integralności i ochrony dostępu. Stosowana jest kontrola dostępu (karty dostępu) z rejestrowaniem każdego wejścia, monitoring środowiskowy oraz CCTV. Użytkownicy potwierdzili, że wejście do serwerowni wymaga autoryzacji i jest każdorazowo rejestrowane.

\textbf{Kontrola dostępu fizycznego do obiektów} (DORA art.~9, NIS2 art.~21): organizacja posiada formalną politykę kontroli dostępu fizycznego do pomieszczeń przetwarzających dane. Dostęp do serwerowni jest logowany elektronicznie, bez możliwości wejścia wyłącznie na podstawie klucza mechanicznego. Wynik potwierdzony zgodnie przez wszystkie trzy perspektywy.

\textbf{Monitoring CCTV i rejestracja wideo} (RODO art.~5 ust.~1 lit.~e, NIS2 art.~21): obszary krytyczne są objęte monitoringiem kamer z retencją nagrań zgodną z polityką bezpieczeństwa i wymogami RODO. Kamery obejmują punkty wejścia do stref krytycznych, a nagrania są chronione przed modyfikacją.

\textbf{Zasilanie i redundancja energetyczna} (DORA art.~11, NIS2 art.~21): obiekty krytyczne posiadają redundantne zasilanie z niezależnych źródeł oraz zasilacze UPS. Urządzenia zasilające są regularnie testowane pod obciążeniem.

\textbf{Wynoszenie i transport urządzeń} (RODO art.~32, NIS2 art.~21): organizacja posiada procedurę wynoszenia urządzeń poza siedzibę z rejestrem wydanych sprzętów. Urządzenia przenośne mają aktywowane szyfrowanie całego dysku oraz możliwość zdalnego wyczyszczenia w przypadku utraty lub kradzieży.

\textbf{Niszczenie i utylizacja nośników danych} (RODO art.~5, NIS2 art.~21): bank posiada procedurę bezpiecznego niszczenia nośników danych, obejmującą dyski twarde, nośniki przenośne oraz wydruki. Dyski są niszczone fizycznie lub nadpisywane skuteczną metodą przed utylizacją.

\textbf{Strefy bezpieczeństwa i segmentacja fizyczna} (NIS2 art.~21, ISO 27001 A.11): w tym obszarze stwierdzono niezgodność. Obiekty nie są w wystarczającym stopniu podzielone na odrębne strefy bezpieczeństwa z różnymi poziomami kontroli dostępu. W trakcie audytu z perspektywy użytkownika ustalono, że gość lub dostawca wchodzący do biura ma możliwość swobodnego poruszania się po budynku bez dodatkowej autoryzacji na kolejnych granicach stref. Oznacza to, że osoba posiadająca dostęp do przestrzeni biurowej może potencjalnie zbliżyć się do obszarów technicznych bez skutecznej przeszkody fizycznej lub proceduralnej.

\textbf{Ochrona przed kradzieżą i sabotażem fizycznym} (NIS2 art.~21, ISO 27001 A.11): stwierdzono niezgodność powiązaną przyczynowo z poprzednim obszarem. Podczas audytu zidentyfikowano niewystarczający nadzór nad osobami wizytującymi, które potencjalnie uzyskiwały nadmierny dostęp fizyczny do pomieszczeń zawierających sprzęt IT. Weryfikacja techniczna wykazała, że szafy rack z serwerami są zamknięte na klucz, a porty USB zablokowane, jednak brak segmentacji stref sprawia, że sam dostęp do szafy pozostaje fizycznie możliwy dla osoby nieupoważnionej, jeśli nie jest ona nadzorowana przez pracownika.

Oba zidentyfikowane ryzyka mają charakter zaleceń (nie wymagań obowiązkowych w rozumieniu NIS2 i DORA) i wynikają z tej samej luki: braku formalnej procedury nadzoru nad gośćmi oraz braku fizycznej segmentacji stref o różnym poziomie wrażliwości. Jako działania naprawcze wskazano wdrożenie rejestru gości z obowiązkowym eskortowaniem poza strefą publiczną, wprowadzenie bramek lub drzwi z osobną autoryzacją na granicy strefy biurowej i technicznej oraz rozważenie instalacji systemu zarządzania gośćmi (visitor management system) rejestrującego każde wejście i wyjście w czasie rzeczywistym.

\subsection{Wyniki sekcji APLIKACJA}

Sekcja APLIKACJA obejmuje dziewięć obszarów dotyczących bezpieczeństwa warstwy aplikacyjnej systemu CTN, w tym zarządzania dostępem, szyfrowania, kopii zapasowych, logowania i zarządzania podatnościami, wymaganych przez NIS2, DORA i RODO. Spośród dziewięciu obszarów sześć zostało ocenionych pozytywnie przez wszystkie trzy perspektywy. Trzy obszary ujawniły nieprawidłowości o zróżnicowanym charakterze: jedno dotyczyło wyłącznie luki w wiedzy użytkowników (uzupełnionej na miejscu), jedno stanowiło pytanie blokujące (usunięte przed zakończeniem audytu), a jedno wskazuje na brak technicznego wymuszenia polityki haseł.

\textbf{Uprawnienia i dostępy aplikacji dostawcy} (DORA art.~9, NIS2 art.~21): konta serwisowe aplikacji dostawcy posiadają ograniczony zakres uprawnień w usłudze katalogowej, bez dostępu do systemów niezwiązanych z obsługą CTN, i są regularnie przeglądane.

\textbf{Dostęp zdalny dostawcy} (DORA art.~28): zdalny dostęp dostawcy do systemu jest kontrolowany, logowany i ograniczony czasowo. Sesje serwisowe wymagają jawnej akceptacji po stronie organizacji, co wyklucza możliwość nieautoryzowanego dostępu przez dostawcę bez wiedzy banku.

\textbf{Rejestr zasobów ICT} (DORA art.~8): system CTN figuruje w oficjalnym rejestrze zasobów ICT z przypisaną klasyfikacją krytyczności, właścicielem biznesowym i właścicielem technicznym. Użytkownicy potwierdzili wiedzę o tym, kto odpowiada za bezpieczeństwo aplikacji.

\textbf{Retencja i ochrona logów aplikacji} (DORA art.~10, NIS2 art.~21): logi aplikacji są centralnie agregowane i przechowywane przez wymagany okres, chronione przed modyfikacją w magazynie immutable. W trakcie audytu perspektywa użytkownika nie była w stanie samodzielnie potwierdzić faktu przechowywania logów, co wskazuje na lukę w komunikacji wewnętrznej na temat procedur retencji. Podczas sesji audytowej mechanizm logowania został zademonstrowany przez dział IT i zweryfikowany przez audytorów bezpośrednio w systemie. Uchybienie zostało ocenione jako formalne (brak wiedzy użytkownika, nie brak technicznej implementacji) i nie skutkuje obniżeniem oceny technicznej obszaru.

\textbf{Szyfrowanie komunikacji} (DORA art.~9, NIS2 art.~21 lit.~h, RODO art.~32): cała komunikacja aplikacji jest szyfrowana przy użyciu protokołu TLS 1.2 lub nowszego. Skanowanie konfiguracji potwierdziło brak przestarzałych protokołów (TLS 1.0/1.1, SSL). Interfejs użytkownika jest dostępny wyłącznie przez HTTPS.

\textbf{MFA dla wszystkich użytkowników} (DORA art.~9, NIS2 art.~21): w momencie rozpoczęcia audytu stwierdzono brak obowiązkowego uwierzytelnienia wieloskładnikowego (MFA) dla kont użytkowników systemu CTN, w tym kont uprzywilejowanych. Pytanie to zostało sklasyfikowane przez metodologię ORION jako \textbf{pytanie blokujące} (wymóg obowiązkowy), co oznacza, że jego niespełnienie automatycznie obniża wynik sekcji niezależnie od wyników pozostałych pytań. Uchybienie zostało niezwłocznie zgłoszone producentowi systemu. W trakcie trwania audytu dostawca wdrożył stosowną zmianę konfiguracyjną, a audytorzy zweryfikowali jej działanie technicznie. Wynik obszaru został skorygowany do pozytywnego po potwierdzeniu wdrożenia.

\textbf{Zarządzanie hasłami i polityką haseł} (NIS2 art.~21, DORA art.~9): organizacja posiada formalną politykę haseł określającą minimalną długość, złożoność oraz zakaz ponownego użycia poprzednich haseł. Polityka jest wymuszana technicznie przez dostawcę tożsamości. Stwierdzono jednak, że system nie wymusza zmiany hasła po upływie określonego czasu (brak mechanizmu wygasania haseł). Jest to niezgodność o charakterze wymagania obowiązkowego: brak cyklicznej rotacji haseł zwiększa ryzyko długotrwałego wykorzystywania przejętych poświadczeń bez wiedzy organizacji. Jako działanie naprawcze wskazano konfigurację polityki wygasania haseł po maksymalnie 90 dniach z obowiązkową zmianą przy kolejnym logowaniu.

\textbf{Backup} (DORA art.~12, NIS2 art.~21 ust.~2 lit.~c): kopie zapasowe są przechowywane w geograficznie odrębnej lokalizacji, retencja spełnia wymogi regulacyjne, a procedura odtwarzania była regularnie testowana z udokumentowanymi wynikami. Użytkownicy wykazali świadomość procedur odtwarzania systemu po awarii.

\textbf{Zarządzanie podatnościami} (NIS2 art.~21 ust.~2 lit.~m, DORA art.~9): organizacja posiada formalny proces skanowania podatności CVE i zarządzania aktualizacjami z udokumentowanymi terminami naprawy przypisanymi do właścicieli. Skanowanie systemów krytycznych jest przeprowadzane co najmniej raz w miesiącu.

Najistotniejszym ustaleniem sekcji jest wykrycie braku MFA jako pytania blokującego. Fakt jego usunięcia w trakcie audytu dokumentuje jedną z kluczowych właściwości metodologii ORION: identyfikacja konkretnego, weryfikowalnego uchybienia i jego naprawa w czasie rzeczywistym jest bardziej wartościowym wynikiem audytu niż brak niezgodności wynikający z powierzchownej oceny. Otwartą kwestią pozostaje brak wymuszania rotacji haseł, wymagający działania naprawczego po stronie organizacji.

\subsection{Wyniki sekcji CHMURA}

Sekcja CHMURA obejmuje sześć obszarów dotyczących zgodności wynikającej ze specyfiki przetwarzania danych w środowisku zewnętrznym, w tym wymagań RODO, DORA i NIS2 wobec relacji z dostawcą usługi chmurowej. Sekcja została aktywowana z uwagi na klasyfikację systemu CTN jako \texttt{CHMURA\_EOG}. Wszystkie sześć obszarów uzyskało wyniki pozytywne. Dostawca infrastruktury jest polskim centrum danych działającym wyłącznie w jurysdykcji EOG, co w istotny sposób upraszcza część wymagań w porównaniu z dostawcami podlegającymi prawu państw trzecich.

Charakterystycznym ustaleniem tej sekcji był niski poziom wiedzy użytkowników o obowiązkach wynikających z modelu chmurowego, przy jednoczesnej pełnej zgodności po stronie technicznej i prawno-compliance. Perspektywa użytkownika ujawniła, że pracownicy banku nie byli świadomi, jakie obowiązki w zakresie bezpieczeństwa spoczywają na organizacji jako administratorze danych, a jakie przejmuje dostawca chmury jako podmiot przetwarzający. Luka ta nie skutkuje niezgodnością formalną, jednak wskazuje na potrzebę przeprowadzenia dedykowanego szkolenia z zakresu modelu współodpowiedzialności (shared responsibility model) obowiązującego w środowiskach chmurowych.

\textbf{Dokumentacja DPIA dla przetwarzania w chmurze} (RODO art.~35): przed uruchomieniem systemu CTN przeprowadzono i udokumentowano ocenę skutków dla ochrony danych (DPIA), obejmującą opis przepływów danych osobowych, zidentyfikowane ryzyka i zastosowane środki techniczne. Dokument jest aktualny i był aktualizowany po ostatniej istotnej zmianie infrastruktury.

\textbf{Szyfrowanie at-rest} (DORA art.~9, RODO art.~32): dane organizacji przechowywane w systemie są szyfrowane w spoczynku. Obszar ten, sklasyfikowany jako zalecenie (nie wymóg obowiązkowy), dotyczy modelu zarządzania kluczami szyfrującymi. Organizacja potwierdziła wdrożenie szyfrowania po stronie dostawcy, ze świadomością, że pełna kontrola nad kluczami po stronie organizacji (model HYOK) pozostaje dobrą praktyką wartą rozważenia przy kolejnym przeglądzie architektury.

\textbf{Kary umowne dla dostawcy} (NIS2 art.~34, DORA art.~30): umowa z dostawcą zawiera klauzulę odpowiedzialności i regresu, umożliwiającą organizacji dochodzenie zwrotu kar administracyjnych nałożonych przez regulatora (KNF, UODO) w przypadku, gdy ich przyczyną był błąd lub zaniedbanie po stronie dostawcy. Użytkownicy nie byli świadomi tej klauzuli, co jest typowym zjawiskiem dla postanowień umownych o charakterze prawno-regulacyjnym, których znajomość nie jest wymagana operacyjnie od personelu liniowego.

\textbf{Prawo do audytu dostawcy} (DORA art.~30 ust.~2 lit.~f, RODO art.~28 ust.~3 lit.~h): umowa z dostawcą gwarantuje organizacji prawo do przeprowadzenia audytu, w tym audytu na miejscu (onsite), bez możliwości zastąpienia go wyłącznie raportem własnym dostawcy. Polskie centrum danych jako podmiot działający w jurysdykcji krajowej nie zgłasza w tym zakresie ograniczeń proceduralnych ani prawnych, charakterystycznych dla dostawców z krajów trzecich, gdzie realizacja audytu onsite bywa utrudniona lub niemożliwa. Użytkownicy zostali podczas audytu poinformowani o przysługującym organizacji prawie do weryfikacji dostawcy.

\textbf{Procedura reagowania na incydenty po stronie dostawcy} (DORA art.~19 ust.~3, NIS2 art.~23 ust.~4): umowa z dostawcą określa maksymalny czas powiadomienia organizacji o incydencie bezpieczeństwa, umożliwiający dotrzymanie ustawowego obowiązku zgłoszenia do CSIRT lub KNF w ciągu 24 godzin. Powiadomienie trafia do wyznaczonego punktu kontaktowego i zawiera wymagany zakres informacji o zdarzeniu. Użytkownicy nie znali szczegółów procedury, jednak znajomość ta nie leży w zakresie ich obowiązków operacyjnych.

\textbf{Izolacja danych między klientami} (DORA art.~9, RODO art.~32, ISO 27017): dostawca zapewnia pełną logiczną izolację danych organizacji od danych innych klientów (tenantów). Dane systemu CTN są przechowywane w dedykowanej przestrzeni przypisanej wyłącznie do banku. Organizacja posiada dokumentację potwierdzającą mechanizmy izolacji. Użytkownicy zostali zapoznani z ryzykiem wynikającym z modelu multi-tenancy i znaczeniem technicznej izolacji danych w środowiskach współdzielonych.

Wyniki sekcji CHMURA potwierdzają, że wybór krajowego dostawcy działającego w jurysdykcji EOG istotnie ogranicza ryzyko regulacyjne związane z przetwarzaniem w chmurze. Brak ekspozycji na przepisy o ekstraterytorialnym dostępie do danych (US CLOUD Act, FISA Section~702) eliminuje potrzebę przeprowadzenia TIA, a dostępność audytu onsite u dostawcy jest gwarantowana umownie i nieobarczona ograniczeniami prawnymi. Główną rekomendacją wynikającą z tej sekcji jest uzupełnienie wiedzy użytkowników na temat modelu współodpowiedzialności w środowiskach chmurowych, tak aby personel operacyjny rozumiał granicę między obowiązkami banku a obowiązkami dostawcy w przypadku incydentu lub kontroli regulatora.

\subsection{Wyniki sekcji DANE\_POUFNE}

Sekcja DANE\_POUFNE obejmuje cztery obszary dotyczące ochrony danych osobowych i tajemnicy bankowej, wymaganych przez RODO. Sekcja została aktywowana z uwagi na klasyfikację systemu CTN jako przetwarzającego dane kategorii \texttt{TAJEMNICA} ($D_3$). We wszystkich czterech obszarach organizacja uzyskała wyniki pozytywne, a odpowiedzi były spójne we wszystkich trzech perspektywach. Bank wykazał się pełną wiedzą proceduralną i techniczną w zakresie ochrony danych szczególnie wrażliwych, co jest wynikiem wyróżniającym się na tle typowych obserwacji audytowych w sektorze finansowym, gdzie formalne polityki są często niedostatecznie internalizowane przez personel operacyjny.

\textbf{Zgłaszanie naruszeń danych osobowych do UODO} (RODO art.~33): organizacja posiada odrębną procedurę zgłaszania naruszeń ochrony danych osobowych do Prezesa UODO w ciągu 72 godzin, wyraźnie rozróżnioną od procedury zgłoszeń ICT do CSIRT i KNF. Istnieją oddzielne ścieżki eskalacji dla obu typów zdarzeń, co eliminuje ryzyko błędnego zakwalifikowania naruszenia danych osobowych wyłącznie jako incydentu technicznego i pominięcia obowiązku zgłoszenia do organu ochrony danych. Użytkownicy wykazali świadomość obowiązku powiadamiania poszkodowanych osób i znali procedurę jego realizacji.

\textbf{Inspektor Ochrony Danych} (RODO art.~37--39): organizacja powołała Inspektora Ochrony Danych, zarejestrowała go w UODO i udostępniła jego dane kontaktowe klientom oraz pracownikom. IOD posiada bezpośredni dostęp do zarządu i działa niezależnie, bez podlegania instrukcjom w zakresie wykonywania swoich zadań ustawowych. Użytkownicy wiedzieli, kto pełni funkcję IOD i jak można się z nim skontaktować, co świadczy o skutecznej komunikacji wewnętrznej w tym obszarze.

\textbf{Rejestr Czynności Przetwarzania} (RODO art.~30): organizacja prowadzi aktualny rejestr czynności przetwarzania (RCP) w formie elektronicznej, obejmujący cele przetwarzania, podstawy prawne, kategorie danych i odbiorców oraz okresy retencji i stosowane środki bezpieczeństwa dla każdej czynności. Rejestr jest aktualizowany w ramach procesu zarządzania zmianami, dzięki czemu odzwierciedla aktualny stan systemu. Użytkownicy potwierdzili znajomość jego istnienia i rozumieli cel jego prowadzenia.

\textbf{Harmonogram retencji i usuwania danych osobowych} (RODO art.~5): organizacja posiada formalny harmonogram retencji z określonymi maksymalnymi okresami przechowywania danych osobowych i mechanizmem ich automatycznego usuwania lub anonimizacji po upływie tych okresów. Proces jest logowany, co umożliwia jego weryfikację przez audytorów. Użytkownicy rozumieli, że dane klientów nie mogą być przechowywane bezterminowo i znali ogólne zasady obowiązującej polityki retencji.

Pełna zgodność we wszystkich czterech obszarach, z wysokim poziomem świadomości we wszystkich trzech perspektywach, potwierdza dojrzałość organizacji w zakresie ochrony danych osobowych. Szczególnie istotna jest spójność między procedurami zgłaszania naruszeń do UODO i do CSIRT (oddzielne ścieżki eskalacji), która eliminuje jedno z najczęstszych uchybień proceduralnych obserwowanych w audytach sektora finansowego, polegające na myleniu incydentu ICT z naruszeniem ochrony danych osobowych w rozumieniu RODO.

\subsection{Wyniki sekcji DOSTAWCY}

Sekcja DOSTAWCY obejmuje siedem obszarów dotyczących zarządzania ryzykiem zewnętrznych dostawców ICT, wymaganych przez DORA, NIS2 i wytyczne EBA. Sekcja została aktywowana z uwagi na klasyfikację systemu CTN jako utrzymywanego w trybie \texttt{CHMURA\_EOG}. W toku audytu zidentyfikowano jednego zewnętrznego dostawcę ICT, krajową spółkę świadczącą usługę utrzymania i wdrożenia systemu. Weryfikacja umów i dokumentacji dostawcy nie wykazała żadnych niezgodności. Wszystkie siedem obszarów uzyskało ocenę pozytywną.

\textbf{Retencja i ochrona logów systemów krytycznych} (DORA art.~10, NIS2 art.~21): logi systemów krytycznych są przechowywane przez wymagany okres i chronione przed modyfikacją. Centralna agregacja logów oraz retencja w magazynie immutable zostały zweryfikowane przez audytorów bezpośrednio w środowisku dostawcy.

\textbf{Due diligence i ocena ryzyka dostawcy} (DORA art.~28, EBA/GL/2019/02): przed zawarciem umowy z dostawcą systemu CTN przeprowadzono due diligence obejmujące weryfikację posiadanych certyfikatów bezpieczeństwa, lokalizacji przetwarzania danych, listy sub-procesorów oraz wyników testów penetracyjnych środowiska dostawcy. Użytkownicy rozumieli zasadność weryfikacji nowych dostawców ICT przed nawiązaniem współpracy.

\textbf{Prawa audytu dostawcy} (DORA art.~30 ust.~2 lit.~f, EBA/GL/2019/02): umowa z dostawcą zawiera zapis gwarantujący organizacji prawo do przeprowadzenia audytu, w tym audytu na miejscu (onsite), bez możliwości zastąpienia go wyłącznie raportem własnym dostawcy. Organizacja posiada dostęp do logów infrastrukturalnych i konfiguracji bezpieczeństwa środowiska dostawcy, a ostatni raport z audytu był dostępny do wglądu audytorów z potwierdzeniem zaadresowania wszystkich obserwacji.

\textbf{Transfer Impact Assessment} (DORA art.~28, Komunikat UKNF 2020): dostawca jest podmiotem krajowym działającym w jurysdykcji EOG, wobec czego ocena skutków transferu (TIA) w rozumieniu przepisów o transferze danych do państw trzecich nie jest wymagana. Dokumentacja dostawcy potwierdza, że żaden etap przetwarzania nie angażuje podmiotów spoza EOG, w tym na poziomie sub-procesorów.

\textbf{Klauzule SLA i gwarantowana dostępność} (DORA art.~30 ust.~2): umowa z dostawcą zawiera mierzalne i egzekwowalne wskaźniki poziomu usług (SLA) z gwarantowaną dostępnością dostosowaną do krytyczności systemu CTN, wraz z sankcjami umownymi za niedotrzymanie uzgodnionych parametrów. Użytkownicy wykazali znajomość dopuszczalnych okien niedostępności systemu wynikających z umowy.

\textbf{Klauzule powiadomień o zmianach u dostawcy} (DORA art.~30 ust.~2 lit.~e): umowa zobowiązuje dostawcę do uprzedniego informowania organizacji o istotnych zmianach, w tym zmianach w sub-procesorach, lokalizacji danych, architekturze bezpieczeństwa i strukturze właścicielskiej, z zachowaniem minimalnego okresu wyprzedzenia. Brak powiadomienia w wymaganym terminie uprawnia organizację do rozwiązania umowy bez kary.

\textbf{Klauzule odpowiedzialności i regres od dostawcy} (DORA art.~30 ust.~2 lit.~d, NIS2 art.~34): umowa zawiera klauzulę umożliwiającą organizacji dochodzenie regresem kosztów kar regulacyjnych nałożonych na bank przez KNF lub UODO w przypadku, gdy ich przyczyną był błąd lub zaniedbanie po stronie dostawcy. Zakres klauzuli odpowiedzialności nie jest ograniczony wyłącznie do wartości ostatniej faktury.

Wyniki sekcji DOSTAWCY potwierdzają, że bank prowadzi relację z dostawcą systemu CTN zgodnie z wymaganiami DORA w zakresie zarządzania ryzykiem stron trzecich. Krajowy charakter dostawcy i jego działalność wyłącznie w jurysdykcji EOG istotnie upraszcza zarządzanie ryzykiem w łańcuchu dostaw: brak sub-procesorów z krajów trzecich eliminuje konieczność prowadzenia rozbudowanego monitoringu na poziomie poddostawców, który jest typowym źródłem niezgodności w przypadku dostawców korzystających z infrastruktury globalnych hiperskalerów.

\subsection{Wyniki sekcji EXIT\_PLAN}

Sekcja EXIT\_PLAN obejmuje sześć obszarów dotyczących gotowości organizacji do zakończenia współpracy z dostawcą systemu CTN w sposób kontrolowany, bez utraty danych i bez zakłócenia ciągłości działania, wymaganych przez DORA art.~28. Sekcja została aktywowana z uwagi na klasyfikację audytowanego podmiotu jako \texttt{FINANSOWY} i obecność zewnętrznego dostawcy ICT. Pięć spośród sześciu obszarów uzyskało ocenę pozytywną. Zidentyfikowano jedno uchybienie w obszarze testowania planu, omówione poniżej. Bank wykazał się wysoką świadomością proceduralną i prawną, natomiast audyt ujawnił lukę w wiedzy użytkowników o tym, czym jest exit plan i jakie uprawnienia organizacji z niego wynikają, którą uzupełniono w toku sesji audytowej.

\textbf{Istnienie i zatwierdzenie exit planu} (DORA art.~28 ust.~2 lit.~f): organizacja posiada formalny exit plan zatwierdzony przez zarząd, traktowany jako dokument żywy i aktualizowany przy każdej istotnej zmianie systemu lub dostawcy. Dokument zawiera konkretne kroki techniczne, narzędzia migracyjne i sekwencje wyłączeń, a nie jedynie deklaratywny opis intencji. Użytkownicy potwierdzili wiedzę o istnieniu planu i jego ogólnym zakresie.

\textbf{Kompletność i możliwość eksportu danych} (DORA art.~30 ust.~2 lit.~h): dla systemu CTN zdefiniowano i przetestowano procedurę pełnego eksportu danych do formatu niezależnego od narzędzi własnościowych dostawcy. Czas i kompletność eksportu są udokumentowane w ramach testów ciągłości działania. Użytkownicy początkowo nie byli pewni, czy dane można pobrać w formacie standardowym, jednak kwestia ta została wyjaśniona i zademonstrowana podczas sesji audytowej jako element podnoszenia wiedzy o uprawnieniach przysługujących organizacji wobec dostawcy.

\textbf{Czas i koszty wyjścia} (DORA art.~28): exit plan zawiera udokumentowaną analizę czasu pełnej migracji (RTE) oraz kosztów wyjścia zatwierdzoną przez zarząd, obejmującą czas migracji danych, koszt wdrożenia alternatywnego systemu, koszty równoległego utrzymania obu środowisk oraz ryzyko utraty ciągłości działania. Użytkownicy nie posiadali szczegółowej wiedzy o szacowanych czasach migracji, co jest akceptowalne, ponieważ informacje te są adresowane do zarządu i działu IT, a nie do personelu operacyjnego.

\textbf{Alternatywny dostawca i przenośność usługi} (DORA art.~28): organizacja zidentyfikowała co najmniej jednego alternatywnego dostawcę zdolnego do przejęcia funkcji systemu CTN i oceniła jego zdolność do onboardingu. Architektura systemu opiera się na otwartych standardach, bez uzależnienia od własnościowych formatów danych lub interfejsów API uniemożliwiających migrację. Użytkownicy zostali poinformowani o istnieniu zidentyfikowanego dostawcy alternatywnego, co wzmacnia ich poczucie ciągłości obsługi.

\textbf{Bezpieczne usunięcie danych u dostawcy po wyjściu} (RODO art.~28 ust.~3 lit.~g, DORA art.~30): umowa z dostawcą zawiera klauzulę zobowiązującą go do wydania certyfikatu zniszczenia danych po zakończeniu współpracy, obejmującego kopie zapasowe, logi i dane archiwalne. Procedura bezpiecznego usunięcia danych jest opisana w exit planie z przypisaną odpowiedzialnością i terminami. Użytkownicy zostali zapoznani z obowiązkiem uzyskania takiego certyfikatu jako elementem ochrony danych klientów po zakończeniu umowy.

\textbf{Testowanie exit planu} (DORA art.~28): stwierdzono niezgodność. Exit plan nie był dotychczas testowany w żadnej formie, ani jako rzeczywiste przełączenie środowiska, ani jako ćwiczenie stołowe (tabletop exercise). DORA art.~28 wymaga regularnego testowania planów wyjścia z udokumentowanymi wynikami dostępnymi dla audytorów. Brak testów oznacza, że poprawność procedur migracyjnych, kompletność eksportu danych i realność zadeklarowanych parametrów RTE pozostają niezweryfikowane empirycznie. Jako działanie naprawcze wskazano przeprowadzenie w najbliższym cyklu rocznym testu eksportu danych do formatu niezależnego od dostawcy, z udokumentowaniem czasu i kompletności eksportu, oraz opracowanie harmonogramu regularnych ćwiczeń z exit planu.

Wyniki sekcji EXIT\_PLAN wskazują na wysoką dojrzałość organizacji w zakresie dokumentacji i struktury planu wyjścia, z jednym istotnym uchybieniem: brakiem jakiegokolwiek testowania planu. Pozostałe pięć obszarów (istnienie i zatwierdzenie planu, możliwość eksportu danych, analiza czasu i kosztów wyjścia, identyfikacja alternatywnego dostawcy, procedura bezpiecznego usunięcia danych) było kompletnych i zgodnych z wymogami DORA. Priorytetowym działaniem naprawczym jest przeprowadzenie pierwszego testu exit planu we współpracy z dostawcą, obejmującego rzeczywisty eksport danych, tak aby plan przestał być wyłącznie dokumentem, a stał się zweryfikowaną procedurą operacyjną. Rekomendowanym działaniem uzupełniającym jest podniesienie świadomości personelu operacyjnego o tym, czym jest exit plan i jakie uprawnienia (prawo do eksportu danych, prawo do certyfikatu zniszczenia, prawo do audytu) przysługują organizacji wobec dostawcy.

\subsection{Wyniki sekcji PLAN\_CIAGLOSCI\_DZIALANIA}

Sekcja PLAN\_CIAGLOSCI\_DZIALANIA obejmuje pięć obszarów dotyczących gotowości organizacji do utrzymania ciągłości działania i odtworzenia funkcji krytycznych po poważnym incydencie ICT, wymaganych przez DORA art.~11 i NIS2 art.~21. Sekcja została aktywowana z uwagi na klasyfikację systemu CTN jako \texttt{WAŻNY} i klasyfikację podmiotu jako \texttt{FINANSOWY}. Wszystkie pięć obszarów uzyskało ocenę pozytywną. Podobnie jak w sekcjach CHMURA i EXIT\_PLAN, wyniki wymagały uzupełnienia wiedzy użytkowników, którzy nie zawsze rozumieli czym jest plan ciągłości działania i jakie uprawnienia oraz obowiązki się z nim wiążą. Dział compliance i dział IT wykazały się wysoką świadomością i znajomością dokumentacji.

\textbf{Istnienie i zatwierdzenie planów BCP i DRP} (DORA art.~11, NIS2 art.~21 ust.~2 lit.~c): organizacja posiada formalne plany ciągłości działania (BCP) i odtwarzania po awarii (DRP) zatwierdzone przez zarząd, obejmujące funkcje krytyczne ICT. Dokumenty zawierają konkretne kroki techniczne, sekwencje przełączeń i drzewa decyzyjne dla poszczególnych scenariuszy, nie zaś ogólne deklaracje intencji. Użytkownicy potwierdzili ogólną wiedzę o procedurach awaryjnych, w tym o zastępczym trybie pracy w przypadku niedostępności systemu CTN.

\textbf{Scenariusze zagrożeń w BCP i DRP} (DORA art.~11, NIS2 art.~21): plany obejmują konkretne scenariusze zagrożeń, w tym atak ransomware, utratę centrum danych, niedostępność dostawcy chmury i awarię kaskadową w łańcuchu dostaw. Dla każdego scenariusza zdefiniowano odrębną ścieżkę odtwarzania z konkretnymi krokami technicznymi i przypisanymi odpowiedzialnościami. Użytkownicy zostali zapoznani z podstawowymi scenariuszami w trakcie sesji audytowej, co pozwoliło uzupełnić wiedzę operacyjną personelu niezbędną do skutecznego reagowania na zdarzenia.

\textbf{Testowanie BCP i DRP} (DORA art.~11 ust.~6): plany były testowane w wymaganym cyklu rocznym, z udokumentowanymi wynikami dostępnymi dla audytorów. Testy obejmowały rzeczywiste przełączenie na środowisko zapasowe z mierzonym RTO i RPO, a nie jedynie ćwiczenie stołowe (tabletop exercise). Użytkownicy w większości nie uczestniczyli bezpośrednio w ćwiczeniach awaryjnych, co wskazuje na potrzebę włączenia szerszego grona pracowników operacyjnych w kolejne testy, tak by posiadali praktyczną znajomość procedur przełączenia.

\textbf{Komunikacja kryzysowa} (NIS2 art.~23, DORA art.~14): BCP zawiera plan komunikacji kryzysowej z gotowymi szablonami i upoważnieniami do komunikowania się na zewnątrz. Organizacja dysponuje alternatywnym kanałem komunikacji działającym niezależnie od głównej infrastruktury IT, umożliwiającym koordynację działań w przypadku niedostępności poczty elektronicznej i systemów komunikacji wewnętrznej. Użytkownicy nie zawsze posiadali numery kontaktowe kluczowych osób poza środowiskami cyfrowymi, co zostało wskazane jako obszar wymagający uzupełnienia w ramach szkoleń z procedur kryzysowych.

\textbf{Scenariusz awarii kaskadowej w łańcuchu dostaw} (DORA art.~11, art.~29 ust.~1): organizacja opracowała scenariusz awarii kaskadowej uwzględniający drzewo zależności między dostawcami różnych poziomów i potencjalny wpływ awarii jednego węzła na dostępność funkcji krytycznych. Scenariusz był symulowany w ramach testów BCP. Użytkownicy zostali zapoznani z koncepcją ryzyka kaskadowego podczas sesji audytowej, uzyskując praktyczne zrozumienie tego, w jaki sposób awaria jednego systemu lub dostawcy może propagować się na inne elementy infrastruktury banku.

Wyniki sekcji PLAN\_CIAGLOSCI\_DZIALANIA potwierdzają dojrzałość organizacyjną banku w zakresie zarządzania ciągłością działania. Szczególnie istotne jest przeprowadzenie testów z rzeczywistym przełączeniem środowiska, a nie jedynie ćwiczeń stołowych, co odpowiada wzorcowi empirycznej weryfikacji obserwowanemu w sekcji EXIT\_PLAN. Wspólnym mianownikiem rekomendacji dla obu tych sekcji jest potrzeba systematycznego włączania personelu operacyjnego w ćwiczenia i szkolenia, tak aby procedury kryzysowe były znane i zrozumiałe nie tylko na poziomie zarządu i działu IT, ale również przez pracowników bezpośrednio obsługujących audytowane systemy.

\subsection{Podsumowanie wyników audytu i obliczenie $S_{total}$}

\subsubsection*{Krok 1: Zestawienie wszystkich odpowiedzi}

Audyt systemu CTN objął łącznie 62 pytania rozmieszczone w 9 aktywnych sekcjach. Spośród nich 58 uzyskało odpowiedź twierdzącą (TAK), a 4 odpowiedź negatywną (NIE). Jedno pytanie blokujące (brak MFA) zostało rozwiązane w trakcie audytu i ujęte jako pozytywne w wyniku końcowym.

\begin{table}[H]
\centering
\caption{Zestawienie odpowiedzi audytowych systemu CTN według sekcji}
\label{tab:podsumowanie-odpowiedzi}
\renewcommand{\arraystretch}{1.4}
\begin{tabular}{p{6.5cm}p{1.5cm}p{1.5cm}p{1.5cm}p{2cm}}
\toprule
\textbf{Sekcja} & \textbf{Pytań} & \textbf{TAK} & \textbf{NIE} & \textbf{Niezgodność} \\
\midrule
ORGANIZACJA & 9 & 9 & 0 & brak \\
\midrule
BEZPIECZENSTWO\_FIZYCZNE & 8 & 6 & 2 & zalecane \\
\midrule
BEZPIECZENSTWO\_TECHNICZNE & 8 & 8 & 0 & brak \\
\midrule
APLIKACJA & 9 & 8 & 1 & obowiązkowe \\
\midrule
CHMURA & 6 & 6 & 0 & brak \\
\midrule
DANE\_POUFNE & 4 & 4 & 0 & brak \\
\midrule
DOSTAWCY & 7 & 7 & 0 & brak \\
\midrule
EXIT\_PLAN & 6 & 5 & 1 & obowiązkowe \\
\midrule
PLAN\_CIAGLOSCI\_DZIALANIA & 5 & 5 & 0 & brak \\
\midrule
\textbf{Łącznie} & \textbf{62} & \textbf{58} & \textbf{4} & \\
\bottomrule
\end{tabular}
\end{table}

Zidentyfikowane niezgodności:
\begin{itemize}
  \item \textbf{Strefy bezpieczeństwa i segmentacja fizyczna} (zalecane) oraz \textbf{Ochrona przed kradzieżą i sabotażem fizycznym} (zalecane): brak segmentacji stref i niewystarczający nadzór nad gośćmi.
  \item \textbf{Zarządzanie hasłami: wymuszanie rotacji} (obowiązkowe): brak technicznego mechanizmu cyklicznej zmiany haseł.
  \item \textbf{Testowanie exit planu} (obowiązkowe): plan wyjścia nie był dotychczas testowany.
\end{itemize}

\subsubsection*{Krok 2: Wagi pytań i wzór na wynik sekcji}

Zgodnie z równaniem~\eqref{eq:omega}, każdemu pytaniu przypisywana jest waga $\omega(q_j)$ wynikająca z jego obligatoryjności:

\[
\omega(q_j) =
\begin{cases}
2{,}0 & \text{pytanie obowiązkowe} \\
1{,}0 & \text{pytanie zalecane} \\
0{,}5 & \text{pytanie informacyjne}
\end{cases}
\]

Wynik sekcji wyznaczany jest jako ważona średnia wyników pytań (równanie~\eqref{eq:wynik-sekcji-avg}), przy założeniu pełnej jakości dowodu ($d_{ij} = 1{,}0$) dla wszystkich odpowiedzi twierdzących i zerowej ($d_{ij} = 0$) dla negatywnych:

\[
S(\mathcal{S}_k) =
\frac{\sum_{q_j \in \mathcal{S}_k} S(q_j) \cdot \omega(q_j)}
     {\sum_{q_j \in \mathcal{S}_k} \omega(q_j)}
\]

\subsubsection*{Krok 3: Obliczenie wyników poszczególnych sekcji}

\textbf{ORGANIZACJA}: 9 pytań obowiązkowych, wszystkie TAK.
\[
S(\text{ORG}) = \frac{9 \times 2{,}0 \times 1{,}0}{9 \times 2{,}0} = \frac{18}{18} = 1{,}000
\]

\textbf{BEZPIECZENSTWO\_FIZYCZNE}: 5 obowiązkowych TAK, 1 zalecane TAK (CCTV), 2 zalecane NIE (strefy, ochrona).
\[
S(\text{BF}) = \frac{5 \times 2{,}0 \times 1{,}0 + 1 \times 1{,}0 \times 1{,}0 + 2 \times 1{,}0 \times 0{,}0}{5 \times 2{,}0 + 1 \times 1{,}0 + 2 \times 1{,}0} = \frac{10 + 1 + 0}{10 + 1 + 2} = \frac{11}{13} \approx 0{,}846
\]

\textbf{BEZPIECZENSTWO\_TECHNICZNE}: 5 obowiązkowych TAK, 3 zalecane TAK.
\[
S(\text{BT}) = \frac{5 \times 2{,}0 \times 1{,}0 + 3 \times 1{,}0 \times 1{,}0}{5 \times 2{,}0 + 3 \times 1{,}0} = \frac{10 + 3}{10 + 3} = \frac{13}{13} = 1{,}000
\]

\textbf{APLIKACJA}: 9 obowiązkowych, 8 TAK, 1 NIE (rotacja haseł). Pytanie MFA zostało rozwiązane w trakcie audytu i wliczone jako TAK.
\[
S(\text{APP}) = \frac{8 \times 2{,}0 \times 1{,}0 + 1 \times 2{,}0 \times 0{,}0}{9 \times 2{,}0} = \frac{16}{18} \approx 0{,}889
\]

\textbf{CHMURA}: 4 obowiązkowe TAK, 2 zalecane TAK.
\[
S(\text{CHM}) = \frac{4 \times 2{,}0 \times 1{,}0 + 2 \times 1{,}0 \times 1{,}0}{4 \times 2{,}0 + 2 \times 1{,}0} = \frac{8 + 2}{8 + 2} = \frac{10}{10} = 1{,}000
\]

\textbf{DANE\_POUFNE}: 1 obowiązkowe TAK, 3 zalecane TAK.
\[
S(\text{DP}) = \frac{1 \times 2{,}0 \times 1{,}0 + 3 \times 1{,}0 \times 1{,}0}{1 \times 2{,}0 + 3 \times 1{,}0} = \frac{2 + 3}{2 + 3} = \frac{5}{5} = 1{,}000
\]

\textbf{DOSTAWCY}: 6 obowiązkowych TAK, 1 zalecane TAK.
\[
S(\text{DOS}) = \frac{6 \times 2{,}0 \times 1{,}0 + 1 \times 1{,}0 \times 1{,}0}{6 \times 2{,}0 + 1 \times 1{,}0} = \frac{12 + 1}{12 + 1} = \frac{13}{13} = 1{,}000
\]

\textbf{EXIT\_PLAN}: 4 obowiązkowe TAK, 1 zalecane TAK, 1 informacyjne TAK, 1 obowiązkowe NIE (brak testu).
\[
S(\text{EP}) = \frac{4 \times 2{,}0 \times 1{,}0 + 1 \times 1{,}0 \times 1{,}0 + 1 \times 0{,}5 \times 1{,}0 + 1 \times 2{,}0 \times 0{,}0}{4 \times 2{,}0 + 1 \times 1{,}0 + 1 \times 0{,}5 + 1 \times 2{,}0} = 
\]
\[
S\frac{8 + 1 + 0{,}5 + 0}{8 + 1 + 0{,}5 + 2} = \frac{9{,}5}{11{,}5} \approx 0{,}826
\]
\vspace{0.5em}

\textbf{PLAN\_CIAGLOSCI\_DZIALANIA}: 4 obowiązkowe TAK, 1 zalecane TAK.
\[
S(\text{BCP}) = \frac{4 \times 2{,}0 \times 1{,}0 + 1 \times 1{,}0 \times 1{,}0}{4 \times 2{,}0 + 1 \times 1{,}0} = \frac{8 + 1}{8 + 1} = \frac{9}{9} = 1{,}000
\]

\subsubsection*{Krok 4: Obliczenie $S_{total}$ systemu CTN}

Zgodnie z równaniem~\eqref{eq:wynik-obiektu}, wynik obiektu jest średnią arytmetyczną wyników wszystkich aktywnych sekcji ($|\mathcal{S}(v)| = 9$):

\[
S_{total}(\text{CTN}) = \frac{1}{9} \sum_{\mathcal{S}_k} S(\mathcal{S}_k)
\]
\\
\[
= \frac{1{,}000 + 0{,}846 + 1{,}000 + 0{,}889 + 1{,}000 + 1{,}000 + 1{,}000 + 0{,}826 + 1{,}000}{9}
= \frac{8{,}561}{9}
\approx \mathbf{0{,}951}
\]

\begin{table}[H]
\centering
\caption{Wyniki sekcji i końcowy wynik systemu CTN}
\label{tab:wyniki-sekcji}
\renewcommand{\arraystretch}{1.4}
\begin{tabular}{p{6.5cm}p{2cm}p{2.5cm}p{2.5cm}}
\toprule
\textbf{Sekcja} & \textbf{$S(\mathcal{S}_k)$} & \textbf{Niezgodności} & \textbf{Charakter} \\
\midrule
ORGANIZACJA & $1{,}000$ & 0 & \\
BEZPIECZENSTWO\_FIZYCZNE & $0{,}846$ & 2 & zalecane \\
BEZPIECZENSTWO\_TECHNICZNE & $1{,}000$ & 0 & \\
APLIKACJA & $0{,}889$ & 1 & obowiązkowe \\
CHMURA & $1{,}000$ & 0 & \\
DANE\_POUFNE & $1{,}000$ & 0 & \\
DOSTAWCY & $1{,}000$ & 0 & \\
EXIT\_PLAN & $0{,}826$ & 1 & obowiązkowe \\
PLAN\_CIAGLOSCI\_DZIALANIA & $1{,}000$ & 0 & \\
\midrule
\textbf{$S_{total}$} & $\mathbf{0{,}951}$ & \textbf{4} & \\
\bottomrule
\end{tabular}
\end{table}

\subsubsection*{Krok 5: Klasyfikacja końcowa}

Wynik $S_{total}(\text{CTN}) = 0{,}951$ mieści się w przedziale $[0{,}85;\ 1{,}0]$, co odpowiada klasyfikacji \textbf{\textsc{ZGODNY}} według skali interpretacyjnej metodologii ORION (tabela~\ref{tab:skala-wyniku}).

Żadna z aktywnych sekcji nie uzyskała wyniku zerowego spowodowanego flagą blokującą (pytanie MFA zostało rozwiązane przed zamknięciem audytu). Nie aktywowano flagi \textsc{TIA}, ponieważ dostawca działa w jurysdykcji EOG. Wynik końcowy nie jest nadpisywany przez żaden mechanizm korekcyjny opisany w równaniach~\eqref{eq:wynik-sekcji-blok} i~\eqref{eq:flag-block-sekcja}.

Dwie spośród czterech niezgodności dotyczą pytań zalecanych (nie obowiązkowych), wobec czego ich wpływ na wynik jest ograniczony przez niższe wagi ($\omega = 1{,}0$ wobec $\omega = 2{,}0$). Dwie niezgodności w pytaniach obowiązkowych (rotacja haseł, brak testowania exit planu) obniżają wyniki odpowiednich sekcji, lecz nie zmieniają klasyfikacji końcowej. Ich usunięcie w ramach działań naprawczych podniosłoby $S_{total}$ do wartości $1{,}000$.

\subsubsection*{Graf ICT systemu CTN}

Poniższy graf ICT przedstawia topologię audytowanego środowiska oraz wyniki poszczególnych sekcji w odwzorowaniu na obiekty audytowe. Każdy węzeł zawiera identyfikator obiektu, jego typ i atrybuty klasyfikacyjne oraz wyniki aktywnych sekcji. Krawędź z etykietą wagi $w$ prowadzi od systemu informatycznego do obsługującego go dostawcy; brak linii przerywanej potwierdza, że nie aktywowano triggera TIA.

\begin{figure}[htbp]
  \centering
  \begin{tikzpicture}[
    podmiot/.style={rectangle, rounded corners=6pt,
      draw=black, thick, fill=white,
      minimum width=11cm, minimum height=0.9cm,
      align=center, font=\small\bfseries},
    sys/.style={rectangle, rounded corners=3pt,
      draw=black, thick, fill=white,
      minimum width=11cm,
      align=center, font=\small},
    vendor/.style={ellipse, draw=black, thick, fill=white,
      minimum width=11cm, minimum height=1.1cm,
      align=center, font=\small},
    verdict/.style={rectangle, rounded corners=4pt,
      draw=black, double, double distance=2pt, thick, fill=white,
      minimum width=11cm, minimum height=0.9cm,
      align=center, font=\small\bfseries},
    dep/.style={-{Stealth}, thick},
    lbl/.style={font=\scriptsize, inner sep=2pt}
  ]
  % v0: Bank (PODMIOT)
  \node[podmiot] (v0) at (0,0)
    {$v_0$\quad Bank (PODMIOT)\quad $[\checkmark]$ ZGODNY};
  % v1: System CTN (SYSTEM_INFORMATYCZNY)
  \node[sys] (v1) at (0,-4.2) {%
    $v_1$\quad\textbf{System CTN}\quad\texttt{SYSTEM\_INFORMATYCZNY}, \texttt{CHMURA\_EOG}\\[5pt]
    {\footnotesize
    \begin{tabular}{@{}ccccccccc@{}}
      ORG & BF & BT & APP & CHM & DP & DOS & EP & BCP \\[2pt]
      $1{,}000$ & $0{,}846$ & $1{,}000$ & $0{,}889$ & $1{,}000$ &
      $1{,}000$ & $1{,}000$ & $0{,}826$ & $1{,}000$ \\
    \end{tabular}}\\[5pt]
    $S_{total} = 0{,}951$\quad$[\checkmark]$\quad\textsc{ZGODNY}%
  };
  % d1: Dostawca krajowy (DOSTAWCA)
  \node[vendor] (d1) at (0,-9.5) {%
    $d_1$\quad Dostawca krajowy\quad\texttt{DOSTAWCA}, \texttt{CHMURA\_EOG}
    \quad DOS: $1{,}000$\quad$[\checkmark]$ \textsc{ZGODNY}%
  };
  % krawędzie
  \draw[dep] (v0) -- (v1);
  \draw[dep] (v1) -- node[lbl, right] {$w\!=\!2$ (WAŻNY)} (d1);
  % box z wynikiem końcowym
  \node[verdict] at (0,-12.2) {%
    $S_{total}(\text{CTN}) = 0{,}951$\quad\textsc{ZGODNY} $[\,0{,}85;\ 1{,}0\,]$
    \quad Brak \texttt{FLAG\_TIA}\quad Brak aktywnych blokad%
  };
  \end{tikzpicture}
  \caption{Graf ICT audytowanego systemu CTN. Skróty sekcji: ORG (Organizacja),
    BF (Bezpieczeństwo fizyczne), BT (Bezpieczeństwo techniczne),
    APP (Aplikacja), CHM (Chmura), DP (Dane poufne), DOS (Dostawcy),
    EP (Exit Plan), BCP (Plan ciągłości działania).
    Waga $w\!=\!2$ odpowiada klasyfikacji WAŻNY obsługiwanego systemu.}
  \label{fig:graf-ict-ctn}
\end{figure}



\chapter{Zakończenie}

Rosnąca złożoność środowisk ICT w organizacjach nadzorowanych, wzmocniona przez wejście w życie rozporządzenia DORA i implementację dyrektywy NIS2, stworzyła wyraźną lukę pomiędzy wymaganiami regulacyjnymi a dostępnymi narzędziami ich weryfikacji. Istniejące metodyki audytu (ISO 27001, COBIT, NIST CSF) dostarczają ogólnych ram kontrolnych, lecz nie oferują skwantyfikowanego, powtarzalnego mechanizmu oceny ryzyka uwzględniającego jednocześnie strukturę łańcucha dostaw ICT, klasyfikację przetwarzanych danych, specyfikę jurysdykcji dostawców chmurowych oraz wagę poszczególnych wymagań regulacyjnych. Niniejsza praca odpowiada na ten deficyt, proponując metodologię ORION (Ocena Ryzyka Informatycznego Organizacji Nadzorowanych), zaprojektowaną z myślą o organizacjach objętych regulacjami sektorowymi i unijną legislacją z zakresu cyfrowej odporności operacyjnej.

Centralnym elementem ORION jest model grafu ważonego $G = (V, E, W)$, w którym wierzchołki reprezentują obiekty audytowe trzech typów (PODMIOT, SYSTEM INFORMATYCZNY, DOSTAWCA), a krawędzie kierowane odzwierciedlają zależności usługowe z wagami odpowiadającymi krytyczności obsługiwanych systemów. Model ten umożliwia propagację ryzyka wzdłuż łańcucha dostaw i automatyczne rozszerzanie zakresu audytu o poddostawców w przypadku aktywacji triggera TIA (Third-party Impact Analysis), co bezpośrednio adresuje wymogi art. 28 DORA dotyczące mapowania zależności zewnętrznych. Mechanizm walidacji krzyżowej (każde pytanie zadawane niezależnie trzem perspektywom: compliance/prawnej, IT oraz użytkownika) ogranicza ryzyko konfabulacji i formalnego spełniania wymagań bez ich realnego wdrożenia. Wynik końcowy $S_{total}$ jest wielkością skwantyfikowaną, wyznaczaną na podstawie jednolitego wzoru ważonego z jasno zdefiniowaną skalą pięciostopniową, co czyni ocenę porównywalną między audytami i organizacjami.

Przeprowadzony audyt pilotażowy systemu CTN w instytucji bankowej potwierdził wykonalność i spójność metodologii w warunkach rzeczywistych. Spośród 62 pytań kwalifikacyjnych 58 uzyskało odpowiedź pozytywną. Cztery niezgodności (dwie obowiązkowe dotyczące rotacji haseł i braku testu exit planu, dwie zalecane z zakresu bezpieczeństwa fizycznego) przełożyły się na wynik $S_{total} = 0{,}951$, klasyfikowany jako ZGODNY w przedziale $[0{,}85;\ 1{,}0]$. Żadna sekcja nie aktywowała flagi blokującej. Trigger TIA nie został wyzwolony, ponieważ dostawca systemu działa w jurysdykcji EOG, a wszystkie umowy spełniały wymagania DORA w zakresie klauzul audytowych i prawa do wypowiedzenia. Wynik wskazuje, że organizacja prowadzi zarządzanie ryzykiem ICT na dojrzałym poziomie, jednocześnie precyzyjnie wskazując dwa obszary wymagające korekty w ramach działań naprawczych.

Jedną z celowych właściwości projektowych ORION jest jego adaptowalność. Choć poniższy przypadek użycia dotyczy sektora finansowego (banku nadzorowanego przez KNF), architektura metodologii nie zakłada żadnych elementów specyficznych wyłącznie dla bankowości. Biblioteka sekcji (bezpieczeństwo techniczne, chmura, dane poufne, ciągłość działania, zarządzanie dostawcami) stanowi zbiór bloków, które mogą być selektywnie włączane lub wyłączane na podstawie atrybutów obiektu audytowego i katalogu regulacji mającego zastosowanie do danej branży. Szpital przetwarzający dane medyczne klasy D3 aktywuje sekcje dotyczące ochrony danych poufnych z inną wagą niż operator logistyczny, lecz korzysta z tego samego aparatu formalnego i tych samych procedur zbierania dowodów. Dostawca energii objęty NIS2 jako podmiot kluczowy, podmiot administracji publicznej prowadzący rejestr ludności, platforma przetwarzania danych medycznych, czy firma fintech objęta DORA, wszystkie te organizacje mają odmienną specyfikę ryzyka, lecz wspólny mianownik: konieczność wykazania organom nadzoru, że ryzyko ICT jest zarządzane w sposób mierzalny i udokumentowany. ORION dostarcza do tego celu jednolity język.

Żaden system oceny ryzyka nie może być lepszy niż jakość danych, na których operuje, a ta zależy przede wszystkim od kompetencji i rzetelności audytora. Kluczowe znaczenie ma konsekwentne przestrzeganie zasady niezależności perspektyw w walidacji krzyżowej: zbyt swobodna interpretacja odpowiedzi jednej perspektywy jako potwierdzenia dla pozostałych wypacza wynik i podważa sens całego mechanizmu. Równie istotna jest wytrwałość w dążeniu do weryfikacji dowodów (logi, konfiguracje, umowy, wyniki testów) zamiast poprzestania na deklaracjach słownych. Przypadek wykrycia braku MFA w trakcie audytu CTN, niezgodności naprawionej jeszcze przed zamknięciem kwestionariusza, ilustruje, jak rzetelna inspekcja może mieć bezpośredni, pozytywny wpływ na bezpieczeństwo systemu. Metodologia ORION tworzy ramy do przeprowadzenia takiej inspekcji, lecz jej realna wartość zależy od gotowości audytora do konfrontowania się z trudnymi obserwacjami i formułowania wniosków bez ulegania presji organizacyjnej.

Autor identyfikuje kilka kierunków dalszego rozwoju metodologii. Po pierwsze, automatyzacja zbierania i walidacji dowodów technicznych (logi SIEM, wyniki skanerów podatności, raporty certyfikatów TLS) może istotnie skrócić czas trwania fazy kwestionariuszowej i ograniczyć ryzyko błędów ludzkich. Po drugie, mechanizm propagacji ryzyka w grafie ICT może zostać rozszerzony o modele probabilistyczne uwzględniające historyczne dane o incydentach, co pozwoliłoby na szacowanie oczekiwanej straty operacyjnej powiązanej z konkretnym poddostawcą. Po trzecie, opracowanie branżowych szablonów pytań (dla sektora zdrowotnego, energetycznego, administracji publicznej) pozwoliłoby skrócić etap konfiguracji audytu przy zachowaniu wspólnego aparatu obliczeniowego. Po czwarte, integracja z rejestrami CMDB i systemami GRC (Governance, Risk and Compliance) umożliwiłaby ciągłą, automatyczną aktualizację grafu ICT bez konieczności inicjowania pełnego audytu przy każdej zmianie w środowisku.

Metodologia ORION nie pretenduje do bycia jedynym ani wyczerpującym narzędziem oceny ryzyka ICT w organizacjach nadzorowanych. Jest natomiast propozycją wypełnienia konkretnej luki: braku ustrukturyzowanej, skwantyfikowanej i regulacyjnie zakorzenionej metody łączącej audyt organizacyjny z analizą łańcucha dostaw w jednolitym modelu grafowym. Zaprezentowany przypadek użycia potwierdza, że taki model jest wykonalny w praktyce i generuje wnioski działające dwukierunkowo: dając organom nadzoru czytelny wskaźnik zgodności, a organizacji audytowanej precyzyjny priorytet działań naprawczych. Dalszy rozwój metodologii, jej walidacja na większej próbie organizacji i ewentualna standaryzacja w postaci normy branżowej lub wytycznych regulatora, pozostają zadaniami otwartymi i stanowią naturalny kierunek kontynuacji niniejszych badań.

\chapter{Słownik pojęć}

\csvreader[
separator=semicolon,
head to column names
]{slownik.csv}{}{%
\noindent\textbf{\skrot} (\rozwiniecie)\\
\wyjasnienie

\smallskip
\noindent\rule{\textwidth}{0.4pt}
\medskip
}



















% \chapter{Standary wytwarzania}
% \section{Rola narzędzi wspomagających mapowanie wymagań}
% \section{Zgodność IT jako proces ciągły (DevSecOps)}
% \section{Praktyczne zastosowanie: Niwelowanie luki komunikacyjnej}

% \chapter{Podsumowanie i przyszłość}

% \section{AI Act i jego wpływ na Fintech/MedTech}
% \section{Post-quantum cryptography (PQC) jako nadchodzące wyzwanie}

% \appendix
% \chapter{Dodatkowe rozważania techniczne}
% \section{Formalna weryfikacja kodu i analizatory SAST/DAST}
% \section{Governance as Code: Open Policy Agent (OPA)}

\printbibliography
\end{document}